% Generated by task tex-catalogue from dossiers/code-spine.md (section D, section G and the
% appendix), dossiers/code-pubinput.md (sections C and D) and dossiers/code-layer1.md
% (sections I and J). Every probe source and output is the dossier's.
\section{The probes on the built code}\label{sec:gen-probes}

The probes of the three code dossiers: what each establishes, its source, the command that ran it and its recorded output. All ran against leanerVM \texttt{main} at \texttt{b435631} with the old pins (ArkLib \texttt{dca90385}, CompPoly \texttt{3468b38c}, VCVio \texttt{f9dc47d9}), before the checkout moved to \texttt{144c5aa} on 2026-09-30; none can run at the new pins without a rebuild. A probe that was written and not run, or argued on paper, says so in its heading. \enquote{Accepted} means that the file elaborated with exit status 0, so that every \texttt{\#guard} evaluated to \texttt{true} and every \texttt{example} and theorem was checked by the kernel.

\subsection{The spine: non-vacuity by probe (D)}\label{sec:gen-probes-spine}

\par\noindent{\small\textit{Cross-references.} A reference such as \enquote{(D.5)} or \enquote{section B} in this part is to a section of the dossier \file{dossiers/code-spine.md}, whose codes the headings keep in parentheses: A (catalogue) is \cref{sec:gen-cat-spine}; B (audit surface) is \cref{sec:gen-surface-spine}; C (the statements) is \cref{sec:gen-statements-spine}; D (non-vacuity by probe), G (what was not verified) and the appendix (the probes) is \cref{sec:gen-probes-spine}; E (conformance) is \cref{sec:gen-conf-spine}; F (findings) is \cref{sec:gen-findings-spine}.}\par

\begin{quote}\small
\textbf{Ground moved during the review.} The checkout now holds \texttt{main} at \texttt{144c5aa} (Lean 4.34.1, all pins moved). Everything here is about \texttt{b435631}; every Lean line number is \texttt{git show b435631:<path>}. All seven probes (\texttt{P1} to \texttt{P6}, with \texttt{P3} in three files) were run and recorded \textbf{before} the move, at \texttt{b435631} with the old pins; their outputs are in the appendix. Three probes (\texttt{P2Relation}, \texttt{P3aSeams}, \texttt{P5PassThrough}) use the numeral \texttt{(\allowbreak{}2 : K)}, which at the old pin is the polynomial \texttt{x} (nonzero, not Boolean) and at the new CompPoly pin is \texttt{0}; they must be re-run with \texttt{K.\allowbreak{}ofBits 2} after a rebuild, and the conclusions they support are marked accordingly. One probe was planned and not written (the read-everything phase, D.3 (d)); the sibling dossier reaches the same conclusion on paper.
\end{quote}

All probes ran against leanerVM \texttt{main} at \texttt{b435631} (ArkLib \texttt{dca90385}, CompPoly \texttt{3468b38c}, VCVio \texttt{f9dc47d9}), each under the shared lock. Their full sources and outputs are in the appendix. \enquote{Accepted} below means the file elaborated with exit status 0, so every \texttt{\#guard} evaluated to \texttt{true} and every \texttt{example} and theorem was checked by the kernel.

Library objects used in this section, introduced once:

\begin{itemize}
\item \textbf{\texttt{OracleReduction}} (ArkLib): a pair of an honest prover and a verifier for a fixed message schedule (\texttt{ProtocolSpec n}: for each of \texttt{n} rounds, who speaks and the type of what is sent). The verifier may query the input oracles and the prover's messages.
\item \textbf{\texttt{Verifier.\allowbreak{}GuardedForm}} (ArkLib \texttt{S\allowbreak{}e\allowbreak{}c\allowbreak{}u\allowbreak{}r\allowbreak{}i\allowbreak{}t\allowbreak{}y\allowbreak{}/\allowbreak{}C\allowbreak{}o\allowbreak{}o\allowbreak{}r\allowbreak{}d\allowbreak{}i\allowbreak{}n\allowbreak{}a\allowbreak{}t\allowbreak{}e\allowbreak{}W\allowbreak{}i\allowbreak{}s\allowbreak{}e\allowbreak{}S\allowbreak{}p\allowbreak{}e\allowbreak{}c\allowbreak{}i\allowbreak{}a\allowbreak{}l\allowbreak{}S\allowbreak{}o\allowbreak{}u\allowbreak{}n\allowbreak{}d\allowbreak{}n\allowbreak{}e\allowbreak{}s\allowbreak{}s\allowbreak{}/\allowbreak{}G\allowbreak{}u\allowbreak{}a\allowbreak{}r\allowbreak{}d\allowbreak{}e\allowbreak{}d\allowbreak{}.\allowbreak{}l\allowbreak{}e\allowbreak{}a\allowbreak{}n\allowbreak{}:\allowbreak{}1\allowbreak{}1\allowbreak{}2}): the data of a Boolean check and a verdict function, with the proof that the verifier outputs the verdict when the check passes and rejects otherwise.
\item \textbf{\texttt{Extractor.\allowbreak{}RoundByRound}} (ArkLib \texttt{S\allowbreak{}e\allowbreak{}c\allowbreak{}u\allowbreak{}r\allowbreak{}i\allowbreak{}t\allowbreak{}y\allowbreak{}/\allowbreak{}R\allowbreak{}o\allowbreak{}u\allowbreak{}n\allowbreak{}d\allowbreak{}B\allowbreak{}y\allowbreak{}R\allowbreak{}o\allowbreak{}u\allowbreak{}n\allowbreak{}d\allowbreak{}.\allowbreak{}l\allowbreak{}e\allowbreak{}a\allowbreak{}n\allowbreak{}:\allowbreak{}7\allowbreak{}7}): not one function from transcripts to witnesses, but a family of intermediate witness types \texttt{WitMid m}, one per round boundary \texttt{m = 0 … n}, with \texttt{WitMid 0} the input witness type, a map \texttt{extractOut} from the output witness to \texttt{WitMid n}, and maps \texttt{e\allowbreak{}x\allowbreak{}t\allowbreak{}r\allowbreak{}a\allowbreak{}c\allowbreak{}t\allowbreak{}M\allowbreak{}i\allowbreak{}d\allowbreak{} m\allowbreak{} :\allowbreak{} W\allowbreak{}i\allowbreak{}t\allowbreak{}M\allowbreak{}i\allowbreak{}d\allowbreak{} (\allowbreak{}m\allowbreak{}+\allowbreak{}1\allowbreak{})\allowbreak{} →\allowbreak{} W\allowbreak{}i\allowbreak{}t\allowbreak{}M\allowbreak{}i\allowbreak{}d\allowbreak{} m} that walk back one round at a time, each seeing the transcript so far.
\item \textbf{\texttt{V\allowbreak{}e\allowbreak{}r\allowbreak{}i\allowbreak{}f\allowbreak{}i\allowbreak{}e\allowbreak{}r\allowbreak{}.\allowbreak{}K\allowbreak{}n\allowbreak{}o\allowbreak{}w\allowbreak{}l\allowbreak{}e\allowbreak{}d\allowbreak{}g\allowbreak{}e\allowbreak{}S\allowbreak{}t\allowbreak{}a\allowbreak{}t\allowbreak{}e\allowbreak{}F\allowbreak{}u\allowbreak{}n\allowbreak{}c\allowbreak{}t\allowbreak{}i\allowbreak{}o\allowbreak{}n}} (ArkLib \texttt{RoundByRound.\allowbreak{}lean:164}): a predicate \texttt{t\allowbreak{}o\allowbreak{}F\allowbreak{}u\allowbreak{}n\allowbreak{} m\allowbreak{} s\allowbreak{}t\allowbreak{}m\allowbreak{}t\allowbreak{} t\allowbreak{}r\allowbreak{}a\allowbreak{}n\allowbreak{}s\allowbreak{}c\allowbreak{}r\allowbreak{}i\allowbreak{}p\allowbreak{}t\allowbreak{} w\allowbreak{}i\allowbreak{}t\allowbreak{}M\allowbreak{}i\allowbreak{}d} that (\texttt{toFun\_\allowbreak{}empty}) is the input relation on the empty transcript, (\texttt{toFun\_\allowbreak{}next}) cannot be made true by a prover message (if it is true after the message on \texttt{w}, it was true before on \texttt{extractMid w}), and (\texttt{toFun\_\allowbreak{}full}) is true on every full transcript on which the verifier can output a statement in the output relation.
\item \textbf{\texttt{r\allowbreak{}b\allowbreak{}r\allowbreak{}K\allowbreak{}n\allowbreak{}o\allowbreak{}w\allowbreak{}l\allowbreak{}e\allowbreak{}d\allowbreak{}g\allowbreak{}e\allowbreak{}S\allowbreak{}o\allowbreak{}u\allowbreak{}n\allowbreak{}d\allowbreak{}n\allowbreak{}e\allowbreak{}s\allowbreak{}s\allowbreak{}W\allowbreak{}o\allowbreak{}r\allowbreak{}s\allowbreak{}t\allowbreak{}C\allowbreak{}a\allowbreak{}s\allowbreak{}e\allowbreak{}W\allowbreak{}i\allowbreak{}t\allowbreak{}h}} (ArkLib \texttt{RoundByRound.\allowbreak{}lean:553}): for one named extractor and one named knowledge state function, for every input statement, every challenge round and every fixed transcript prefix, the probability over the fresh challenge that the state goes from false to true is at most the error declared for that challenge.
\end{itemize}

\begin{leancode}
-- .lake/packages/Arklib/ArkLib/OracleReduction/Security/RoundByRound.lean:553-568 (ArkLib dca90385)
def rbrKnowledgeSoundnessWorstCaseWith
    (relIn : Set (StmtIn × WitIn)) (relOut : Set (StmtOut × WitOut))
    (verifier : Verifier oSpec StmtIn StmtOut pSpec)
    (WitMid : Fin (n + 1) → Type)
    (extractor : Extractor.RoundByRound oSpec StmtIn WitIn WitOut pSpec WitMid)
    (kSF : verifier.KnowledgeStateFunction init impl relIn relOut extractor)
    (rbrKnowledgeError : pSpec.ChallengeIdx → ℝ≥0) : Prop :=
  ∀ stmtIn : StmtIn,
  ∀ i : pSpec.ChallengeIdx,
  ∀ transcript : Transcript i.1.castSucc pSpec,
    Pr[fun challenge =>
      ∃ witMid,
        ¬ kSF i.1.castSucc stmtIn transcript
          (extractor.extractMid i.1 stmtIn (transcript.concat challenge) witMid) ∧
          kSF i.1.succ stmtIn (transcript.concat challenge) witMid
      | $ᵗ (pSpec.Challenge i)] ≤ rbrKnowledgeError i
\end{leancode}

Two facts about this definition drive the results of D.3. The error is any function \texttt{ChallengeIdx → ℝ≥0}: nothing bounds it by one. And the statement bounds a probability by it: with a declared error of \texttt{1} the inequality holds of every state function.

\subsubsection{\texttt{M3Holds} on the toy instance (D.1)}\label{sec:gen-spine-D1}

\textbf{What the repository's tests do} (\texttt{t\allowbreak{}e\allowbreak{}s\allowbreak{}t\allowbreak{}s\allowbreak{}/\allowbreak{}L\allowbreak{}e\allowbreak{}a\allowbreak{}n\allowbreak{}e\allowbreak{}r\allowbreak{}V\allowbreak{}M\allowbreak{}T\allowbreak{}e\allowbreak{}s\allowbreak{}t\allowbreak{}s\allowbreak{}/\allowbreak{}P\allowbreak{}r\allowbreak{}o\allowbreak{}t\allowbreak{}o\allowbreak{}c\allowbreak{}o\allowbreak{}l\allowbreak{}/\allowbreak{}S\allowbreak{}p\allowbreak{}i\allowbreak{}n\allowbreak{}e\allowbreak{}.\allowbreak{}l\allowbreak{}e\allowbreak{}a\allowbreak{}n\allowbreak{}:\allowbreak{}2\allowbreak{}6\allowbreak{}-{}\allowbreak{}9\allowbreak{}8}, read in full). The honest stack is accepted at statement \texttt{1} (\texttt{:26}). Five mutated stacks or statements fail exactly one clause each, the three other checkable clauses being asserted to hold: \texttt{badConstraint} (\texttt{:32-{}38}), \texttt{badBalance} (\texttt{:41-{}47}), \texttt{badCount} (\texttt{:50-{}56}), the wrong statement (\texttt{:59-{}61}), \texttt{badLine} (\texttt{:65-{}71}). \texttt{badBalanceSum} (\texttt{:76-{}85}) is a stack whose pushed and pulled first coordinates have the same sum in \texttt{K} and are different multisets. The fifth clause, \texttt{aux}, is \texttt{True} on the toy and is tested by nothing. The claim of the test's docstring and of acceptance test 27 (\enquote{each of the four checkable clauses failing alone}) is met.

\textbf{What is missing, and added by probe \texttt{P2Relation.\allowbreak{}lean} (accepted).}

\begin{enumerate}
\item \emph{Multiplicity.} Every mutation of the repository changes a \textbf{value}; none changes a \textbf{count}. The toy cannot: its table pushes two tuples and its boundary pulls two, whatever the stack. The probe builds two variants of the instance.

\begin{itemize}
\item \texttt{toyOnce}: the boundary is one block of height one pulling \texttt{(\allowbreak{}1,\allowbreak{} 0,\allowbreak{} …)}. On the honest stack the pushed tuples are \texttt{(\allowbreak{}1,\allowbreak{}0,\allowbreak{}…),\allowbreak{} (\allowbreak{}1,\allowbreak{}0,\allowbreak{}…)} and the pulled tuple is \texttt{(\allowbreak{}1,\allowbreak{}0,\allowbreak{}…)}: the two sides are the same \textbf{set} (each tuple of one side occurs in the other, both guards pass) and different multisets. Only \texttt{Balanced} fails; \texttt{M3Holds} fails.
\item \texttt{toyNever}: no boundary. The tuple is pushed twice and never pulled. For each of the sixteen coordinates the sum over the pushed side equals the sum over the (empty) pulled side, \texttt{1 + 1 = 0} in characteristic two: a balance summed in the field accepts. Only \texttt{Balanced} fails; \texttt{M3Holds} fails. This is the case the module docstring of \texttt{Instance.\allowbreak{}lean} names (\texttt{:30-{}31}, \enquote{in characteristic two a tuple pushed twice and never pulled would sum to zero}) and that no test exercised.
\end{itemize}
\item \emph{The side filter.} \texttt{toySwap} exchanges the two sides (the table's flush pulls, the boundary pushes). \texttt{flushTuples} and \texttt{boundaryTuples} return the tuples on the declared side and none on the other (lengths \texttt{0,\allowbreak{} 2} and \texttt{2,\allowbreak{} 0}), and the honest stack satisfies \texttt{M3Holds toySwap}. The filters \texttt{decide (\allowbreak{}f.\allowbreak{}1 = s)} and \texttt{decide (\allowbreak{}b.\allowbreak{}side = s)} are therefore not swapped.
\item \emph{The row ranges.} One tuple per row of the table (\texttt{2 \textasciicircum{} τ}) and one per row of a block (\texttt{2 \textasciicircum{} κ}, also for \texttt{κ = 0}).
\item \emph{Which cells \texttt{CountsNonzero} ranges over.} Every row of every column listed in \texttt{counts}, and no other cell: a stack with zeros in column 0, column 2 and the padding cells satisfies the clause; a zero in either cell of column 1 fails it.
\end{enumerate}

\textbf{Conclusion.} \texttt{Balanced} is a multiset equality counted in \texttt{ℕ} and the list construction is right on every case tried. No defect found in the five clauses as definitions.

\subsubsection{The seams, one by one (D.2)}\label{sec:gen-spine-D2}

Probe \texttt{P3aSeams.\allowbreak{}lean} (accepted). The repository gives the claims no decision procedure; the probe supplies them by \texttt{inferInstanceAs} on the unfolded propositions and states, by \texttt{Iff.\allowbreak{}rfl}, that the predicates it decides are the seams.

{\footnotesize
\begin{longtable}{@{}>{\raggedright\arraybackslash}p{0.077\linewidth}>{\raggedright\arraybackslash}p{0.458\linewidth}>{\raggedright\arraybackslash}p{0.410\linewidth}@{}}
\toprule
Seam & Inhabitant (on the honest stack, statement \texttt{1}) & Near misses, each rejected \\
\midrule
\endfirsthead
\toprule
Seam & Inhabitant (on the honest stack, statement \texttt{1}) & Near misses, each rejected \\
\midrule
\endhead
\bottomrule
\endfoot
\texttt{S\allowbreak{}e\allowbreak{}a\allowbreak{}m\allowbreak{}.\allowbreak{}c\allowbreak{}o\allowbreak{}m\allowbreak{}m\allowbreak{}i\allowbreak{}t} & the honest stack (repository test \texttt{:26}, and \texttt{Iff.\allowbreak{}rfl} at \texttt{:166-{}167}) & the repository's five mutations \\
\texttt{Seam.\allowbreak{}bus} & the zerocheck claim of the toy's constraint at the points \texttt{(\allowbreak{}0)} and \texttt{(\allowbreak{}1)} with value \texttt{0}, and the column claim \enquote{column 0 at the point \texttt{(\allowbreak{}x)} is \texttt{1}} & a false linear claim; a \textbf{true} claim whose term is cubic; a false column claim; the wrong statement \\
\texttt{Seam.\allowbreak{}table} & the column claim & a false column claim; the wrong statement \\
\texttt{Seam.\allowbreak{}pub} & the column claim & a false column claim. The wrong statement is \textbf{inside} the seam: the public seam no longer mentions the statement \\
\texttt{Seam.\allowbreak{}flock} & the column claim and the weighted claim \enquote{the weight selecting cell 0 pairs to \texttt{1}} & a false weighted claim; a false column claim \\
\texttt{Seam.\allowbreak{}done} & everything & none: \texttt{S\allowbreak{}e\allowbreak{}a\allowbreak{}m\allowbreak{}.\allowbreak{}d\allowbreak{}o\allowbreak{}n\allowbreak{}e\allowbreak{} t\allowbreak{}o\allowbreak{}y\allowbreak{} =\allowbreak{} S\allowbreak{}e\allowbreak{}t\allowbreak{}.\allowbreak{}u\allowbreak{}n\allowbreak{}i\allowbreak{}v} (proved in the probe) \\
\end{longtable}}

Two further facts the probe establishes, both intended by the design and both worth stating in the report because a reader of the seam table may assume the opposite.

\begin{itemize}
\item \textbf{A seam does not imply the relation before it.} On \texttt{badConstraint} (column 2 is \texttt{[2,\allowbreak{} 0]}, so the constraint \texttt{X₂² − X₂} is violated on row 0) the zerocheck claim at the point \texttt{(\allowbreak{}1)} with value \texttt{0} is true: the extension of the table \texttt{[a,\allowbreak{} 0]} is \texttt{(\allowbreak{}1 − ζ)·a}, zero at \texttt{ζ = 1}. So \texttt{(\allowbreak{}(\allowbreak{}2\allowbreak{},\allowbreak{} ⟨\allowbreak{}[\allowbreak{}z\allowbreak{}e\allowbreak{}r\allowbreak{}o\allowbreak{}A\allowbreak{}t\allowbreak{}1\allowbreak{}]\allowbreak{},\allowbreak{} [\allowbreak{}c\allowbreak{}o\allowbreak{}l\allowbreak{}0\allowbreak{}O\allowbreak{}n\allowbreak{}e\allowbreak{}]\allowbreak{}⟩\allowbreak{})\allowbreak{},\allowbreak{} b\allowbreak{}a\allowbreak{}d\allowbreak{}C\allowbreak{}o\allowbreak{}n\allowbreak{}s\allowbreak{}t\allowbreak{}r\allowbreak{}a\allowbreak{}i\allowbreak{}n\allowbreak{}t\allowbreak{})} is in \texttt{Seam.\allowbreak{}bus toy} and \texttt{badConstraint} is outside \texttt{M3Holds toy 2}. This is the zerocheck escape; the bus phase's error must pay for it.
\item \textbf{A seam does not say which claims are emitted.} \texttt{(\allowbreak{}(\allowbreak{}2\allowbreak{},\allowbreak{} ⟨\allowbreak{}[\allowbreak{}]\allowbreak{},\allowbreak{} [\allowbreak{}]\allowbreak{}⟩\allowbreak{})\allowbreak{},\allowbreak{} b\allowbreak{}a\allowbreak{}d\allowbreak{}C\allowbreak{}o\allowbreak{}n\allowbreak{}s\allowbreak{}t\allowbreak{}r\allowbreak{}a\allowbreak{}i\allowbreak{}n\allowbreak{}t\allowbreak{})} is in \texttt{Seam.\allowbreak{}bus toy}.
\end{itemize}

\subsubsection{\texttt{Phases.\allowbreak{}Complete} and \texttt{Phases.\allowbreak{}Security} (D.3)}\label{sec:gen-spine-D3}

\paragraph{(a) The repository's pass-through phases}

\texttt{trivPhases} and \texttt{trivComplete} (\texttt{tests/\allowbreak{}…/\allowbreak{}Spine.\allowbreak{}lean:103-{}119}) inhabit \texttt{Phases toy} and \texttt{Phases.\allowbreak{}Complete} with five zero-round phases that drop every claim. The test's docstring says that \texttt{Phases.\allowbreak{}Security} \enquote{has no such inhabitant} and gives the reason (\enquote{dropping every claim is … not knowledge sound}), but proves nothing. Probe \texttt{P5PassThrough.\allowbreak{}lean} (accepted; \texttt{\#print axioms no\_\allowbreak{}security} gives the three standard axioms) proves it for the bus phase, where it fails first:

\begin{leancode}
/-- The reflection hypothesis of `Phase.passThroughSecurity` is false at the bus seam. -/
theorem not_reflects : ¬ ∀ (s : K) (o : ∀ i, TheOracle toy i),
    ((dropAll s, o), ()) ∈ Seam.bus toy → ((s, o), ()) ∈ Seam.commit toy

/-- No `Phase.Security` exists for the pass-through bus phase: not only the repository's
constructor fails, the type is empty. -/
theorem no_security (S : Phase.Security toy (Phase.passThrough toy dropAll) (Seam.commit toy)
    (Seam.bus toy)) : False
\end{leancode}

The proof takes the knowledge state function of \texttt{S} at the trivial shared-oracle state, the statement \texttt{2} with the oracle \texttt{badConstraint} (outside \texttt{M3Holds}, inside the bus seam with no claim, both by \texttt{decide +kernel}), and the empty transcript; the verifier accepts it with probability one, so \texttt{toFun\_\allowbreak{}full} makes the state true at the end, which for a phase with no round is the beginning, where \texttt{toFun\_\allowbreak{}empty} says the state is \texttt{M3Holds}. Exactly where the task expected it to fail: a pass-through bus phase cannot reflect \texttt{Seam.\allowbreak{}bus} back into \texttt{M3Holds}.

\paragraph{(b) A cheap inhabitant of \texttt{Phases.\allowbreak{}Security}: five phases that check nothing, at declared error 1}

Probe \texttt{P4Junk.\allowbreak{}lean} (accepted; \texttt{\#print axioms} gives \texttt{propext}, \texttt{Classical.\allowbreak{}choice}, \texttt{Quot.\allowbreak{}sound}).

The phase \texttt{wasted f ε} has one round: the verifier draws a challenge in \texttt{E}, ignores it, makes no query, and outputs \texttt{f} of its input statement. Its verifier accepts every transcript (\texttt{wPure}, a \texttt{Verifier.\allowbreak{}PureForm}).

\begin{itemize}
\item It is perfectly complete against any two seams that \texttt{f} carries one into the other (\texttt{wastedComplete}): the hypothesis of the repository's \texttt{Phase.\allowbreak{}passThroughComplete}.
\item With declared error \texttt{1} it has a \texttt{Phase.\allowbreak{}Security} against \textbf{any} two seams, under that completeness hypothesis alone (\texttt{wastedSecurity}). The knowledge state function is the input seam before the challenge and \texttt{True} after it (\texttt{wState}, which takes no hypothesis on \texttt{f} or on the seams); the bound is \texttt{Pr[…] ≤ 1} (\texttt{w\_\allowbreak{}rbr}).
\end{itemize}

\begin{leancode}
-- probes/code-spine/P4Junk.lean
def wState (f : StmtIn → StmtOut) :
    (wVerifier I f).toVerifier.KnowledgeStateFunction init impl relIn relOut
      (wExtractor I (StmtIn := StmtIn)) where
  toFun := fun m stmt _ _ ↦ if m.val = 0 then (stmt, ()) ∈ relIn else True
  toFun_empty := fun _ _ ↦ by simp
  toFun_next := fun m hm ↦ by
    fin_cases m
    exact absurd hm (by decide)
  toFun_full := fun _ _ _ _ ↦ by simp

noncomputable def junkSecurity : junkPhases.Security where
  bus := wastedSecurity toy _ fun _ _ h ↦ ⟨by simp, by simp, by simp, h.2.2.2.1, h.2.2.2.2⟩
  table := wastedSecurity toy _ fun _ _ h ↦ ⟨by simp, h.2.2.2.1, h.2.2.2.2⟩
  pub := wastedSecurity toy _ fun _ _ h ↦ ⟨by simp, h.2.2⟩
  flock := wastedSecurity toy _ fun _ _ _ ↦ ⟨by simp, by simp⟩
  opening := wastedSecurity toy _ fun _ _ _ ↦ trivial

example {σ : Type} (init : ProbComp σ) (impl : QueryImpl []ₒ (StateT σ ProbComp)) :
    (leanVmVerifier junkPhases).toVerifier.rbrKnowledgeSoundnessWorstCase init impl (M3Rel toy)
      (Seam.done toy) (piopError junkPhases) :=
  piop_rbrKnowledgeSoundness_exists junkPhases junkSecurity init impl

theorem piopError_junk (i : junkPhases.toDef.pSpec.ChallengeIdx) : piopError junkPhases i = 1
\end{leancode}

So \texttt{Phases.\allowbreak{}Security toy} is inhabited, both master theorems apply to \texttt{junkPhases}, and the protocol they are about commits to a stack and accepts it whatever it is. The knowledge theorem is true of it because its error is \texttt{1} at each of its five challenges (\texttt{piopError\_\allowbreak{}junk}). Nothing in the construction uses the toy: the same five phases inhabit \texttt{Phases.\allowbreak{}Security I} for every instance \texttt{I}.

\textbf{Reading.} \texttt{piop\_\allowbreak{}rbrKnowledgeSoundness} has content only together with a bound on \texttt{piopError P}. No such bound is stated in the spine, and the blueprint's \texttt{piopError\_\allowbreak{}le} (Layer 10, \texttt{:1130}) is stated over \texttt{(\allowbreak{}s : Sizes)}, a parameter the spine's \texttt{piopError P} does not have (finding \emph{the declared error is unconstrained}).

\paragraph{(c) An always-rejecting verifier}

A verifier that always rejects is round-by-round knowledge sound at error zero (the state \enquote{the input relation holds} never has to become true). It is excluded from \texttt{Phases.\allowbreak{}Security} by the structure itself:

\begin{leancode}
-- LeanerVM/Protocol/ToArkLib/Component.lean:90-93
structure Security (D : Def StmtIn OStmtIn WitIn StmtOut OStmtOut WitOut)
    (relIn : Set ((StmtIn × ∀ i, OStmtIn i) × WitIn))
    (relOut : Set ((StmtOut × ∀ i, OStmtOut i) × WitOut))
    extends Complete D relIn relOut where
\end{leancode}

\texttt{Complete.\allowbreak{}complete} is perfect completeness on the input seam, so on any statement and stack in the input seam the verifier accepts the honest run with probability one. \texttt{Seam.\allowbreak{}commit toy} is inhabited (D.1). This is a reading of the definitions, not a probe: a machine-checked refutation would need the probability of an event under a failing \texttt{OptionT} computation, which adds nothing to the argument.

The exclusion is only as strong as the input seam is inhabited: on an instance whose relation is empty (D.4 exhibits one) perfect completeness is vacuous and an always-rejecting bundle is a \texttt{Phases.\allowbreak{}Security}.

\paragraph{(d) A phase that reads the whole stack (argued on paper; the probe was not written)}

In the ideal oracle model the verifier may query the stack's extension at any point of \texttt{E\textasciicircum{}μ}, any number of times: neither ArkLib's \texttt{OracleVerifier} nor the spine counts queries. At the \texttt{2\textasciicircum{}μ} points of the cube the answers are the cells. \texttt{M3Holds} and every claim's \texttt{Holds} are decidable. So for every seam there is a zero-round phase whose verifier reads the \texttt{2\textasciicircum{}μ} cells, decides the \textbf{input} seam itself, rejects if it fails and otherwise outputs any statement in the output seam (no claim at all will do). It is perfectly complete, and knowledge sound at error zero: acceptance implies the input seam. Five of them inhabit \texttt{Phases.\allowbreak{}Security I} for every \texttt{I}, with \texttt{piopError = 0} (there is no challenge).

This is \textbf{not machine-checked} here (the probe was not written: it needs an oracle-querying verifier and the unfolding of \texttt{2\textasciicircum{}μ} simulated queries, and the machine's memory limit was reached during this session). The sibling dossier \texttt{gt-{}table-{}pub.\allowbreak{}md} (section 6, E.4 item 6) reaches the same conclusion on paper; the two readings agree. What would verify it: a probe on an instance with \texttt{μ = 0} (one cell, one query), whose verifier is \texttt{d\allowbreak{}o\allowbreak{} l\allowbreak{}e\allowbreak{}t\allowbreak{} v\allowbreak{} ←\allowbreak{} q\allowbreak{}u\allowbreak{}e\allowbreak{}r\allowbreak{}y\allowbreak{};\allowbreak{} i\allowbreak{}f\allowbreak{} d\allowbreak{}e\allowbreak{}c\allowbreak{}i\allowbreak{}d\allowbreak{}e\allowbreak{} (\allowbreak{}r\allowbreak{}e\allowbreak{}l\allowbreak{}I\allowbreak{}n\allowbreak{} s\allowbreak{} ⟨\allowbreak{}\#\allowbreak{}v\allowbreak{}[\allowbreak{}v\allowbreak{}.\allowbreak{}l\allowbreak{}i\allowbreak{}m\allowbreak{}b\allowbreak{} 0\allowbreak{}]\allowbreak{}⟩\allowbreak{})\allowbreak{} t\allowbreak{}h\allowbreak{}e\allowbreak{}n\allowbreak{} p\allowbreak{}u\allowbreak{}r\allowbreak{}e\allowbreak{} (\allowbreak{}f\allowbreak{} s\allowbreak{})\allowbreak{} e\allowbreak{}l\allowbreak{}s\allowbreak{}e\allowbreak{} f\allowbreak{}a\allowbreak{}i\allowbreak{}l\allowbreak{}u\allowbreak{}r\allowbreak{}e}.

\textbf{What (b) and (d) mean together.} The hypotheses of the master theorems are satisfiable, for every instance, by protocols that are not leanVM's: one with no check and error 1, one with no interaction, error 0 and a verifier as expensive as the prover. This is expected of a composition theorem. What pins the phases to leanVM's is outside the spine and, today, outside the repository:

\begin{enumerate}
\item a bound on \texttt{piopError P} (excludes (b));
\item the phases' \emph{definitions} being leanVM's, which no theorem of the oracle protocol states; it is Layer 12's \texttt{verify\_\allowbreak{}iff\_\allowbreak{}compiled}, together with the fixture \enquote{a proof of the pinned Rust prover is accepted by \texttt{verify}}, that ties the composed verifier to the deployed one (excludes (d), whose transcript is empty);
\item review of each phase's \texttt{Def} against the specification.
\end{enumerate}

\subsubsection{What a layout can do (D.4)}\label{sec:gen-spine-D4}

Probe \texttt{P2Relation.\allowbreak{}lean}, last section (accepted).

\texttt{toyAlias} is the toy with the layout that reads \textbf{every} column from cells \texttt{0,\allowbreak{} 1} (\texttt{extend c z = (\allowbreak{}z,\allowbreak{} 0,\allowbreak{} 0)} for every \texttt{c}). The reading law \texttt{read\_\allowbreak{}eval} holds of it (\texttt{alias\_\allowbreak{}read\_\allowbreak{}eval}, proved by the same \texttt{simp} call as the toy's), so it is a \texttt{Layout} and \texttt{toyAlias} is an \texttt{M3Instance}. On it the three columns are one column, and \texttt{M3Holds} asks for incompatible things: balance forces cells \texttt{0,\allowbreak{} 1} to be \texttt{1,\allowbreak{} 1} and the public line on column 2 forces cell 1 to be \texttt{0}. The honest stack fails; the exhaustive guard over cells \texttt{0,\allowbreak{} 1} and the statement in \texttt{\{0,\allowbreak{} 1,\allowbreak{} 2\}} finds no satisfying stack. That the relation is empty for \textbf{every} stack and statement is the two-line argument above, on paper; the guard is evidence, not a proof.

Consequences, for the report:

\begin{itemize}
\item The master theorems hold of \texttt{toyAlias} (they hold of every instance). Completeness is then vacuous and knowledge soundness says that every accepting prover has been lucky.
\item A layout \textbf{cannot ignore the stack}. This is a consequence of the law, on paper: if \texttt{read q c} did not depend on \texttt{q}, the law at the all-zero stack and at the all-one stack would give \texttt{0 = 1}, since the extension of a constant table is that constant at every point. More precisely the law \textbf{determines \texttt{read} from \texttt{extend}}: on a cube point \texttt{x} of the column, \texttt{(\allowbreak{}read q c)[x]} is the stack's extension at \texttt{extend c x}, and an extension at a cube point is the cell. So \texttt{read} is redundant data (proposal in section B), and the only freedom of an instance is \texttt{extend}.
\item What a wrong \texttt{extend} costs is \textbf{completeness and faithfulness, never soundness} of the Lean chain. If the layout aliases or mixes columns, \texttt{M3Holds I} is a relation on the aliased columns; the adaptor's soundness lemma \texttt{satisfiedBy\_\allowbreak{}witnessOf} (\texttt{M3Holds → SatisfiedBy}) is then a theorem about that instance and, if it is proved, the chain to \texttt{ValidExecution} is intact. What breaks is \texttt{m3Holds\_\allowbreak{}stackOf} and \texttt{witnessOf\_\allowbreak{}stackOf} (no stack packs distinct columns into the same cells), and the agreement with the deployed layout (the Lean verifier rejects the Rust prover's proofs).
\item \textbf{Who rules such instances out.} Not the spine. For leanISA: \texttt{witnessOf\_\allowbreak{}stackOf} and \texttt{m3Holds\_\allowbreak{}stackOf} (blueprint Layer 3, not built) for non-vacuity, and the transcription of \texttt{witness.\allowbreak{}rs:67-{}101} with Layer 12's fixture for faithfulness. Until they exist, no statement on \texttt{main} says that the relation of any instance other than the toy is inhabited.
\end{itemize}

\subsubsection{The extractor (D.5)}\label{sec:gen-spine-D5}

\textbf{Is anything marked \texttt{noncomputable}?} Probe \texttt{P1Axioms.\allowbreak{}lean}: \texttt{Lean.\allowbreak{}isNoncomputable} is \texttt{false} for \texttt{piopExtractor}, \texttt{commitExtractor}, \texttt{Component.\allowbreak{}sendExtractor}, \texttt{Phases.\allowbreak{}Security.\allowbreak{}toDef}, \texttt{Component.\allowbreak{}Security.\allowbreak{}append}, ArkLib's \texttt{E\allowbreak{}x\allowbreak{}t\allowbreak{}r\allowbreak{}a\allowbreak{}c\allowbreak{}t\allowbreak{}o\allowbreak{}r\allowbreak{}.\allowbreak{}R\allowbreak{}o\allowbreak{}u\allowbreak{}n\allowbreak{}d\allowbreak{}B\allowbreak{}y\allowbreak{}R\allowbreak{}o\allowbreak{}u\allowbreak{}n\allowbreak{}d\allowbreak{}.\allowbreak{}a\allowbreak{}p\allowbreak{}p\allowbreak{}e\allowbreak{}n\allowbreak{}d}, \texttt{C\allowbreak{}o\allowbreak{}m\allowbreak{}p\allowbreak{}o\allowbreak{}n\allowbreak{}e\allowbreak{}n\allowbreak{}t\allowbreak{}.\allowbreak{}p\allowbreak{}a\allowbreak{}s\allowbreak{}s\allowbreak{}T\allowbreak{}h\allowbreak{}r\allowbreak{}o\allowbreak{}u\allowbreak{}g\allowbreak{}h\allowbreak{}E\allowbreak{}x\allowbreak{}t\allowbreak{}r\allowbreak{}a\allowbreak{}c\allowbreak{}t\allowbreak{}o\allowbreak{}r}, \texttt{PublicInput.\allowbreak{}extractor} and \texttt{publicInputSecurity}. It is \texttt{true} for \texttt{publicInputPhase} only (its error is a real number), as the blueprint says (\texttt{:1049}).

\textbf{Does \texttt{\#print axioms} mean anything here?} No. It reports \texttt{p\allowbreak{}r\allowbreak{}o\allowbreak{}p\allowbreak{}e\allowbreak{}x\allowbreak{}t\allowbreak{},\allowbreak{} C\allowbreak{}l\allowbreak{}a\allowbreak{}s\allowbreak{}s\allowbreak{}i\allowbreak{}c\allowbreak{}a\allowbreak{}l\allowbreak{}.\allowbreak{}c\allowbreak{}h\allowbreak{}o\allowbreak{}i\allowbreak{}c\allowbreak{}e\allowbreak{},\allowbreak{} Q\allowbreak{}u\allowbreak{}o\allowbreak{}t\allowbreak{}.\allowbreak{}s\allowbreak{}o\allowbreak{}u\allowbreak{}n\allowbreak{}d} for \texttt{piopExtractor}, for \texttt{commitExtractor} and for ArkLib's \texttt{E\allowbreak{}x\allowbreak{}t\allowbreak{}r\allowbreak{}a\allowbreak{}c\allowbreak{}t\allowbreak{}o\allowbreak{}r\allowbreak{}.\allowbreak{}R\allowbreak{}o\allowbreak{}u\allowbreak{}n\allowbreak{}d\allowbreak{}B\allowbreak{}y\allowbreak{}R\allowbreak{}o\allowbreak{}u\allowbreak{}n\allowbreak{}d\allowbreak{}.\allowbreak{}a\allowbreak{}p\allowbreak{}p\allowbreak{}e\allowbreak{}n\allowbreak{}d}, all three computable definitions, because the proofs inside them (the casts between witness types) use classical lemmas. \texttt{Classical.\allowbreak{}choice} in the axiom list of a definition does not say that its \textbf{data} was chosen classically. The blueprint does not claim otherwise; the earlier review's validation paragraph lists axioms and should not be read as evidence of computability.

\textbf{Does the repository's test prove that the extractor returns the committed stack?} Only for phases with no round. The test (\texttt{tests/\allowbreak{}…/\allowbreak{}Spine.\allowbreak{}lean:157-{}159}) is

\begin{leancode}
example (S : trivPhases.Security) :
    (piopExtractor trivPhases S).extractOut ((1 : K), fun i : Fin 0 ↦ i.elim0)
      honestFullTranscript () = honest := rfl
\end{leancode}

It is true for the reason ArkLib's \texttt{append} gives (\texttt{A\allowbreak{}p\allowbreak{}p\allowbreak{}e\allowbreak{}n\allowbreak{}d\allowbreak{}/\allowbreak{}S\allowbreak{}t\allowbreak{}a\allowbreak{}t\allowbreak{}e\allowbreak{}F\allowbreak{}u\allowbreak{}n\allowbreak{}c\allowbreak{}t\allowbreak{}i\allowbreak{}o\allowbreak{}n\allowbreak{}.\allowbreak{}l\allowbreak{}e\allowbreak{}a\allowbreak{}n\allowbreak{}:\allowbreak{}1\allowbreak{}2\allowbreak{}9\allowbreak{}-{}\allowbreak{}1\allowbreak{}3\allowbreak{}9}): when the second component has no round, the composed \texttt{extractOut} is the first component's applied to the second's. With five zero-round phases the composed \texttt{extractOut} is the commit phase's, which returns the first message. As soon as one phase has a round, the composed \texttt{extractOut} is the \textbf{last} such phase's, and its value lives in that phase's last intermediate witness type. Probe \texttt{P3cExtractor.\allowbreak{}lean}, with the repository's own \texttt{publicInputPhase} in the bundle and four pass-throughs:

\begin{leancode}
-- accepted
example : (realPubSecurity Sb St Sf So).toDef.witMid (Fin.last 3) = Unit := rfl
example : (realPubSecurity Sb St Sf So).toDef.witMid 0 = Column 3 := rfl
example (s : K) (tr : …Transcript (Fin.succ ⟨0, _⟩)) (w : …witMid (Fin.succ ⟨0, _⟩)) :
    (piopExtractor realPub (realPubSecurity Sb St Sf So)).extractMid ⟨0, _⟩
      (s, fun i : Fin 0 ↦ i.elim0) tr w = (tr ⟨0, _⟩ : Column 3) := rfl
-- rejected, as expected: the test's statement with this bundle
example (tr : realPub.toDef.pSpec.FullTranscript) :
    (piopExtractor realPub (realPubSecurity Sb St Sf So)).extractOut
      ((1 : K), fun i : Fin 0 ↦ i.elim0) tr () = honest := rfl
\end{leancode}

\begin{srccode}
P3cExtractor.lean:58:49: error: Type mismatch
  honest
has type
  Column 3
but is expected to have type
  (realPubSecurity Sb St Sf So).toDef.witMid (Fin.last realPub.toDef.n)
\end{srccode}

So for leanVM's phases \enquote{\texttt{piopExtractor … extractOut} returns the stack} is not a false statement but an ill-typed one. The stack is what the composed extractor returns \textbf{at round 0}, by its \texttt{extractMid} at the commit round, and there it is the transcript's first message whatever it is handed (third example). The spine defines no function \enquote{the stack extracted from a full transcript} (the composite of \texttt{extractOut} and the \texttt{extractMid}s down to round 0); the statement \enquote{the extractor's column \texttt{q} satisfies \texttt{M3Holds}} of the blueprint's headline is therefore not a statement about any definition of the repository (finding \emph{the extracted stack is not a definition}).

\textbf{Could a phase's \texttt{Security} supply an extractor defined by \texttt{Classical.\allowbreak{}choice} without anything noticing?} Yes as to the mechanism, and it would not matter as to the result.

\begin{itemize}
\item \emph{Mechanism.} \texttt{Component.\allowbreak{}Security.\allowbreak{}extractor} is a field of type \texttt{Extractor.\allowbreak{}RoundByRound …}. Lean's types do not record computability. A phase author can write \texttt{n\allowbreak{}o\allowbreak{}n\allowbreak{}c\allowbreak{}o\allowbreak{}m\allowbreak{}p\allowbreak{}u\allowbreak{}t\allowbreak{}a\allowbreak{}b\allowbreak{}l\allowbreak{}e\allowbreak{} d\allowbreak{}e\allowbreak{}f\allowbreak{} f\allowbreak{}o\allowbreak{}o\allowbreak{}S\allowbreak{}e\allowbreak{}c\allowbreak{}u\allowbreak{}r\allowbreak{}i\allowbreak{}t\allowbreak{}y\allowbreak{} :\allowbreak{} P\allowbreak{}h\allowbreak{}a\allowbreak{}s\allowbreak{}e\allowbreak{}.\allowbreak{}S\allowbreak{}e\allowbreak{}c\allowbreak{}u\allowbreak{}r\allowbreak{}i\allowbreak{}t\allowbreak{}y\allowbreak{} …} with a classical \texttt{extractMid}; \texttt{Phases.\allowbreak{}Security.\allowbreak{}toDef}, \texttt{piopExtractor} and both theorems accept it, \texttt{piopExtractor} itself stays computable (it is a projection of its argument), \texttt{\#print axioms} does not change, and the repository's checks (\texttt{audit-{}lean.\allowbreak{}sh} forbids \texttt{axiom}, \texttt{sorry}, \texttt{admit}, \texttt{unsafe}, \texttt{native\_\allowbreak{}decide}, per \texttt{AGENTS.\allowbreak{}md}) do not forbid \texttt{noncomputable}: \texttt{publicInputPhase} is already \texttt{noncomputable}, legitimately. ArkLib itself ships such an extractor, \texttt{E\allowbreak{}x\allowbreak{}t\allowbreak{}r\allowbreak{}a\allowbreak{}c\allowbreak{}t\allowbreak{}o\allowbreak{}r\allowbreak{}.\allowbreak{}R\allowbreak{}o\allowbreak{}u\allowbreak{}n\allowbreak{}d\allowbreak{}B\allowbreak{}y\allowbreak{}R\allowbreak{}o\allowbreak{}u\allowbreak{}n\allowbreak{}d\allowbreak{}O\allowbreak{}n\allowbreak{}e\allowbreak{}S\allowbreak{}h\allowbreak{}o\allowbreak{}t\allowbreak{}.\allowbreak{}t\allowbreak{}o\allowbreak{}R\allowbreak{}o\allowbreak{}u\allowbreak{}n\allowbreak{}d\allowbreak{}B\allowbreak{}y\allowbreak{}R\allowbreak{}o\allowbreak{}u\allowbreak{}n\allowbreak{}d\allowbreak{}O\allowbreak{}f\allowbreak{}R\allowbreak{}e\allowbreak{}l} (\texttt{RoundByRound.\allowbreak{}lean:118-{}123}, \texttt{i\allowbreak{}f\allowbreak{} h\allowbreak{} :\allowbreak{} ∃\allowbreak{} v\allowbreak{},\allowbreak{} (\allowbreak{}s\allowbreak{}t\allowbreak{}m\allowbreak{}t\allowbreak{}I\allowbreak{}n\allowbreak{},\allowbreak{} v\allowbreak{})\allowbreak{} ∈\allowbreak{} r\allowbreak{}e\allowbreak{}l\allowbreak{}I\allowbreak{}n\allowbreak{} t\allowbreak{}h\allowbreak{}e\allowbreak{}n\allowbreak{} h\allowbreak{}.\allowbreak{}c\allowbreak{}h\allowbreak{}o\allowbreak{}o\allowbreak{}s\allowbreak{}e\allowbreak{} e\allowbreak{}l\allowbreak{}s\allowbreak{}e\allowbreak{} w\allowbreak{}i\allowbreak{}t\allowbreak{}I\allowbreak{}n}). Acceptance test 24 is enforced by review only.
\item \emph{Result.} From the commit phase on, a phase's input and output witness types are \texttt{Unit} (\texttt{Phase.\allowbreak{}Def}, \texttt{Phase.\allowbreak{}lean:37-{}38}), and \texttt{eqIn} forces its \texttt{WitMid 0} to be \texttt{Unit}. Whatever a phase's extractor does, at the seam it hands \texttt{(\allowbreak{})} to the phase before it. The stack comes from the commit phase's extractor alone, which is the spine's, is computable, and reads the message. A classical phase extractor cannot produce a stack. The risk acceptance test 24 guards against (a witness chosen because one exists) is excluded by the \textbf{design} \enquote{the oracle is the witness}, not by the test.
\end{itemize}

\subsubsection{The commit phase's security (D.6)}\label{sec:gen-spine-D6}

Probe \texttt{P3bCommit.\allowbreak{}lean} (accepted), every statement by \texttt{Iff.\allowbreak{}rfl} or \texttt{rfl}:

\begin{leancode}
-- Before the message: M3Holds of the candidate witness.
example (s : K) (o : ∀ i, NoOracle i) (w : Column 3) :
    ((commitSecurity toy).kSF init impl).toFun 0 (s, o) default w ↔ M3Holds toy s w := Iff.rfl
-- After the message: M3Holds of the message, whatever the candidate witness.
example (s : K) (o : ∀ i, NoOracle i) (tr : (commitSpec toy).FullTranscript) (w : Column 3) :
    ((commitSecurity toy).kSF init impl).toFun (Fin.last 1) (s, o) tr w ↔
      M3Holds toy s (tr 0) := Iff.rfl
-- The extractor: both maps return the message.
example … : (commitExtractor toy).extractOut (s, o) tr () = tr 0 := rfl
example … : (commitExtractor toy).extractMid 0 (s, o) tr w = tr 0 := rfl
\end{leancode}

The state function says what the docstring says (\texttt{SendOracle.\allowbreak{}lean:143-{}144}). Error zero is right: the phase has no challenge (\texttt{sendOracle\_\allowbreak{}rbr} is the elimination of an empty index type), so its error function is the empty function and the \texttt{0} written in \texttt{sendOracle} (\texttt{:65}) is never applied. The extractor is \enquote{read the message}.

\subsubsection{Three facts behind sections B and C.7 (D.7)}\label{sec:gen-spine-D7}

Probe \texttt{P6Surface.\allowbreak{}lean} (accepted; both theorems with the three standard axioms).

\begin{enumerate}
\item \texttt{outputPure\_\allowbreak{}of\_\allowbreak{}def}: for every \texttt{Component.\allowbreak{}Def} whatsoever, \texttt{D.\allowbreak{}red.\allowbreak{}prover.\allowbreak{}OutputIsPure} holds, because the shared oracle specification is empty (\texttt{[]ₒ}) and a computation with no possible query is its pure result. ArkLib has this at the pin as the instance \texttt{Prover.\allowbreak{}instOutputIsPureEmpty} (\texttt{C\allowbreak{}o\allowbreak{}m\allowbreak{}p\allowbreak{}o\allowbreak{}s\allowbreak{}i\allowbreak{}t\allowbreak{}i\allowbreak{}o\allowbreak{}n\allowbreak{}/\allowbreak{}S\allowbreak{}e\allowbreak{}q\allowbreak{}u\allowbreak{}e\allowbreak{}n\allowbreak{}t\allowbreak{}i\allowbreak{}a\allowbreak{}l\allowbreak{}/\allowbreak{}N\allowbreak{}o\allowbreak{}A\allowbreak{}m\allowbreak{}b\allowbreak{}i\allowbreak{}e\allowbreak{}n\allowbreak{}t\allowbreak{}.\allowbreak{}l\allowbreak{}e\allowbreak{}a\allowbreak{}n\allowbreak{}:\allowbreak{}3\allowbreak{}9\allowbreak{}-{}\allowbreak{}4\allowbreak{}2}), in a module the spine does not import. The field \texttt{C\allowbreak{}o\allowbreak{}m\allowbreak{}p\allowbreak{}o\allowbreak{}n\allowbreak{}e\allowbreak{}n\allowbreak{}t\allowbreak{}.\allowbreak{}C\allowbreak{}o\allowbreak{}m\allowbreak{}p\allowbreak{}l\allowbreak{}e\allowbreak{}t\allowbreak{}e\allowbreak{}.\allowbreak{}o\allowbreak{}u\allowbreak{}t\allowbreak{}p\allowbreak{}u\allowbreak{}t\allowbreak{}P\allowbreak{}u\allowbreak{}r\allowbreak{}e} is therefore redundant (section B).
\item \texttt{@\allowbreak{}R\allowbreak{}e\allowbreak{}f\allowbreak{}i\allowbreak{}n\allowbreak{}e\allowbreak{}m\allowbreak{}e\allowbreak{}n\allowbreak{}t\allowbreak{} :\allowbreak{} \{\allowbreak{}S\allowbreak{}t\allowbreak{}m\allowbreak{}t\allowbreak{} :\allowbreak{} T\allowbreak{}y\allowbreak{}p\allowbreak{}e\allowbreak{}\}\allowbreak{} →\allowbreak{} \{\allowbreak{}W\allowbreak{}₁\allowbreak{} :\allowbreak{} T\allowbreak{}y\allowbreak{}p\allowbreak{}e\allowbreak{} u\allowbreak{}\_\allowbreak{}1\allowbreak{}\}\allowbreak{} →\allowbreak{} \{\allowbreak{}W\allowbreak{}₂\allowbreak{} :\allowbreak{} T\allowbreak{}y\allowbreak{}p\allowbreak{}e\allowbreak{} u\allowbreak{}\_\allowbreak{}2\allowbreak{}\}\allowbreak{} →\allowbreak{} …}: the two witness universes are independent, so a refinement from the stack (\texttt{Type}) to Clean's \texttt{EnsembleWitness} (\texttt{Type 1}) can be formed, as the blueprint's ladder needs (\texttt{:364-{}369}). Checked by \texttt{\#check} and by an example from \texttt{ℕ} to a structure in \texttt{Type 1}.
\item \texttt{knowledge\_\allowbreak{}transport}: if a verifier is knowledge sound for \texttt{R} with a named straight-line extractor at error \texttt{ε} (ArkLib's \texttt{knowledgeSoundnessWith}), then, for the same provers and the same extractor, the probability that the verifier accepts and the extracted witness, mapped by a refinement \texttt{R → S}, is invalid for \texttt{S} is at most \texttt{ε}. Three lines: the event is monotone (\texttt{probEvent\_\allowbreak{}mono}) and \texttt{map\_\allowbreak{}valid} is the pointwise step. No prover conversion and no \texttt{Nonempty} hypothesis: the game is not changed, the map is applied inside its event. This is the probabilistic transport the blueprint says is not proved (\texttt{:582-{}586}), in the form Layer 13 needs (section C.7).
\end{enumerate}

\subsubsection{What was not verified, and what would verify it (G)}\label{sec:gen-spine-G}

{\footnotesize
\begin{longtable}{@{}>{\raggedright\arraybackslash}p{0.239\linewidth}>{\raggedright\arraybackslash}p{0.355\linewidth}>{\raggedright\arraybackslash}p{0.352\linewidth}@{}}
\toprule
Claim & Status & What would verify it \\
\midrule
\endfirsthead
\toprule
Claim & Status & What would verify it \\
\midrule
\endhead
\bottomrule
\endfoot
The relation of \texttt{toyAlias} is empty for every stack and statement (D.4) & evidence by an exhaustive \texttt{\#guard} over cells \texttt{0,\allowbreak{} 1} and statements in \texttt{\{0,\allowbreak{} 1,\allowbreak{} 2\}} at the old pin; the paper argument is two lines & a theorem \texttt{∀\allowbreak{} i\allowbreak{}n\allowbreak{}p\allowbreak{}u\allowbreak{}t\allowbreak{} q\allowbreak{},\allowbreak{} ¬\allowbreak{} M\allowbreak{}3\allowbreak{}H\allowbreak{}o\allowbreak{}l\allowbreak{}d\allowbreak{}s\allowbreak{} t\allowbreak{}o\allowbreak{}y\allowbreak{}A\allowbreak{}l\allowbreak{}i\allowbreak{}a\allowbreak{}s\allowbreak{} i\allowbreak{}n\allowbreak{}p\allowbreak{}u\allowbreak{}t\allowbreak{} q}, whose proof needs CompPoly's evaluation of \texttt{X 0} on a row (no \texttt{eval\_\allowbreak{}X} lemma at \texttt{3468b38c}; the sibling dossiers may have one) \\
\texttt{L\allowbreak{}a\allowbreak{}y\allowbreak{}o\allowbreak{}u\allowbreak{}t\allowbreak{}.\allowbreak{}r\allowbreak{}e\allowbreak{}a\allowbreak{}d} is determined by \texttt{extend} (C.2) & on paper & a lemma \texttt{r\allowbreak{}e\allowbreak{}a\allowbreak{}d\allowbreak{}\_\allowbreak{}e\allowbreak{}q\allowbreak{}\_\allowbreak{}o\allowbreak{}f\allowbreak{}\_\allowbreak{}e\allowbreak{}x\allowbreak{}t\allowbreak{}e\allowbreak{}n\allowbreak{}d} from CompPoly's \texttt{e\allowbreak{}v\allowbreak{}a\allowbreak{}l\allowbreak{}\_\allowbreak{}m\allowbreak{}l\allowbreak{}e\allowbreak{}\_\allowbreak{}e\allowbreak{}q\allowbreak{}\_\allowbreak{}e\allowbreak{}v\allowbreak{}a\allowbreak{}l} (cited in \texttt{F\allowbreak{}i\allowbreak{}e\allowbreak{}l\allowbreak{}d\allowbreak{}.\allowbreak{}l\allowbreak{}e\allowbreak{}a\allowbreak{}n\allowbreak{}:\allowbreak{}5\allowbreak{}0}) \\
An always-rejecting bundle has no \texttt{P\allowbreak{}h\allowbreak{}a\allowbreak{}s\allowbreak{}e\allowbreak{}s\allowbreak{}.\allowbreak{}S\allowbreak{}e\allowbreak{}c\allowbreak{}u\allowbreak{}r\allowbreak{}i\allowbreak{}t\allowbreak{}y} on an instance with an inhabited relation (D.3 (c)) & a reading of \texttt{S\allowbreak{}e\allowbreak{}c\allowbreak{}u\allowbreak{}r\allowbreak{}i\allowbreak{}t\allowbreak{}y\allowbreak{} e\allowbreak{}x\allowbreak{}t\allowbreak{}e\allowbreak{}n\allowbreak{}d\allowbreak{}s\allowbreak{} C\allowbreak{}o\allowbreak{}m\allowbreak{}p\allowbreak{}l\allowbreak{}e\allowbreak{}t\allowbreak{}e} and of \texttt{p\allowbreak{}e\allowbreak{}r\allowbreak{}f\allowbreak{}e\allowbreak{}c\allowbreak{}t\allowbreak{}C\allowbreak{}o\allowbreak{}m\allowbreak{}p\allowbreak{}l\allowbreak{}e\allowbreak{}t\allowbreak{}e\allowbreak{}n\allowbreak{}e\allowbreak{}s\allowbreak{}s} & a probe with a rejecting zero-round verifier and \texttt{p\allowbreak{}e\allowbreak{}r\allowbreak{}f\allowbreak{}e\allowbreak{}c\allowbreak{}t\allowbreak{}C\allowbreak{}o\allowbreak{}m\allowbreak{}p\allowbreak{}l\allowbreak{}e\allowbreak{}t\allowbreak{}e\allowbreak{}n\allowbreak{}e\allowbreak{}s\allowbreak{}s\allowbreak{}\_\allowbreak{}e\allowbreak{}q\allowbreak{}\_\allowbreak{}p\allowbreak{}r\allowbreak{}o\allowbreak{}b\allowbreak{}\_\allowbreak{}o\allowbreak{}n\allowbreak{}e} \\
The read-everything phase inhabits \texttt{P\allowbreak{}h\allowbreak{}a\allowbreak{}s\allowbreak{}e\allowbreak{}s\allowbreak{}.\allowbreak{}S\allowbreak{}e\allowbreak{}c\allowbreak{}u\allowbreak{}r\allowbreak{}i\allowbreak{}t\allowbreak{}y\allowbreak{} I} at error \texttt{0} for every \texttt{I} (D.3 (d)) & on paper; agrees with \texttt{g\allowbreak{}t\allowbreak{}-{}\allowbreak{}t\allowbreak{}a\allowbreak{}b\allowbreak{}l\allowbreak{}e\allowbreak{}-{}\allowbreak{}p\allowbreak{}u\allowbreak{}b\allowbreak{}.\allowbreak{}m\allowbreak{}d} section 6 & a probe on an instance with \texttt{μ = 0}: one oracle query, \texttt{i\allowbreak{}f\allowbreak{} d\allowbreak{}e\allowbreak{}c\allowbreak{}i\allowbreak{}d\allowbreak{}e\allowbreak{} (\allowbreak{}r\allowbreak{}e\allowbreak{}l\allowbreak{}I\allowbreak{}n\allowbreak{} s\allowbreak{} ⟨\allowbreak{}\#\allowbreak{}v\allowbreak{}[\allowbreak{}v\allowbreak{}.\allowbreak{}l\allowbreak{}i\allowbreak{}m\allowbreak{}b\allowbreak{} 0\allowbreak{}]\allowbreak{}⟩\allowbreak{})\allowbreak{} t\allowbreak{}h\allowbreak{}e\allowbreak{}n\allowbreak{} p\allowbreak{}u\allowbreak{}r\allowbreak{}e\allowbreak{} (\allowbreak{}f\allowbreak{} s\allowbreak{})\allowbreak{} e\allowbreak{}l\allowbreak{}s\allowbreak{}e\allowbreak{} f\allowbreak{}a\allowbreak{}i\allowbreak{}l\allowbreak{}u\allowbreak{}r\allowbreak{}e}, with \texttt{V\allowbreak{}e\allowbreak{}r\allowbreak{}i\allowbreak{}f\allowbreak{}i\allowbreak{}e\allowbreak{}r\allowbreak{}.\allowbreak{}G\allowbreak{}u\allowbreak{}a\allowbreak{}r\allowbreak{}d\allowbreak{}e\allowbreak{}d\allowbreak{}F\allowbreak{}o\allowbreak{}r\allowbreak{}m\allowbreak{}.\allowbreak{}o\allowbreak{}f\allowbreak{}\_\allowbreak{}p\allowbreak{}r\allowbreak{}o\allowbreak{}b\allowbreak{}E\allowbreak{}v\allowbreak{}e\allowbreak{}n\allowbreak{}t\allowbreak{}\_\allowbreak{}p\allowbreak{}o\allowbreak{}s} for the state function \\
The probes that use \texttt{(\allowbreak{}2 : K)} mean the same at the new CompPoly pin & they do not (\texttt{(\allowbreak{}2\allowbreak{} :\allowbreak{} K\allowbreak{})\allowbreak{} =\allowbreak{} 0} there, brief §8); their conclusions are about \texttt{b435631} with the old pins, where \texttt{2 : K} is the polynomial \texttt{x} (\texttt{t\allowbreak{}e\allowbreak{}s\allowbreak{}t\allowbreak{}s\allowbreak{}/\allowbreak{}…\allowbreak{}/\allowbreak{}S\allowbreak{}p\allowbreak{}i\allowbreak{}n\allowbreak{}e\allowbreak{}.\allowbreak{}l\allowbreak{}e\allowbreak{}a\allowbreak{}n\allowbreak{}:\allowbreak{}3\allowbreak{}0}) & re-run \texttt{P2Relation}, \texttt{P3aSeams}, \texttt{P\allowbreak{}5\allowbreak{}P\allowbreak{}a\allowbreak{}s\allowbreak{}s\allowbreak{}T\allowbreak{}h\allowbreak{}r\allowbreak{}o\allowbreak{}u\allowbreak{}g\allowbreak{}h} with \texttt{K.\allowbreak{}ofBits 2} after the rebuild \\
Line numbers of \texttt{Field.\allowbreak{}lean} and \texttt{K\allowbreak{}n\allowbreak{}o\allowbreak{}w\allowbreak{}l\allowbreak{}e\allowbreak{}d\allowbreak{}g\allowbreak{}e\allowbreak{}A\allowbreak{}p\allowbreak{}p\allowbreak{}e\allowbreak{}n\allowbreak{}d\allowbreak{}.\allowbreak{}l\allowbreak{}e\allowbreak{}a\allowbreak{}n} & at \texttt{b435631} (the working tree had moved by the time the counting script first ran; the script was re-run on \texttt{g\allowbreak{}i\allowbreak{}t\allowbreak{} s\allowbreak{}h\allowbreak{}o\allowbreak{}w\allowbreak{} b\allowbreak{}4\allowbreak{}3\allowbreak{}5\allowbreak{}6\allowbreak{}3\allowbreak{}1\allowbreak{}:} and the numbers in sections B and C are from that run) & — \\
The earlier review's validation paragraph (\texttt{d\allowbreak{}o\allowbreak{}c\allowbreak{}s\allowbreak{}/\allowbreak{}r\allowbreak{}e\allowbreak{}v\allowbreak{}i\allowbreak{}e\allowbreak{}w\allowbreak{}s\allowbreak{}/\allowbreak{}p\allowbreak{}r\allowbreak{}o\allowbreak{}t\allowbreak{}o\allowbreak{}c\allowbreak{}o\allowbreak{}l\allowbreak{}-{}\allowbreak{}s\allowbreak{}p\allowbreak{}i\allowbreak{}n\allowbreak{}e\allowbreak{}.\allowbreak{}m\allowbreak{}d\allowbreak{}:\allowbreak{}1\allowbreak{}9\allowbreak{}-{}\allowbreak{}2\allowbreak{}8}) & not re-run (no build allowed); \texttt{\#\allowbreak{}p\allowbreak{}r\allowbreak{}i\allowbreak{}n\allowbreak{}t\allowbreak{} a\allowbreak{}x\allowbreak{}i\allowbreak{}o\allowbreak{}m\allowbreak{}s} re-established by probe P1 for the declarations listed there and more & \texttt{.\allowbreak{}/\allowbreak{}s\allowbreak{}c\allowbreak{}r\allowbreak{}i\allowbreak{}p\allowbreak{}t\allowbreak{}s\allowbreak{}/\allowbreak{}v\allowbreak{}a\allowbreak{}l\allowbreak{}i\allowbreak{}d\allowbreak{}a\allowbreak{}t\allowbreak{}e\allowbreak{}.\allowbreak{}s\allowbreak{}h} on \texttt{b435631} \\
\end{longtable}}

\subsection{The spine's probes in full}\label{sec:gen-probes-spine-files}

Every probe is a plain Lean file under \texttt{.\allowbreak{}c\allowbreak{}l\allowbreak{}a\allowbreak{}u\allowbreak{}d\allowbreak{}e\allowbreak{}/\allowbreak{}r\allowbreak{}e\allowbreak{}p\allowbreak{}o\allowbreak{}r\allowbreak{}t\allowbreak{}s\allowbreak{}/\allowbreak{}b\allowbreak{}l\allowbreak{}u\allowbreak{}e\allowbreak{}p\allowbreak{}r\allowbreak{}i\allowbreak{}n\allowbreak{}t\allowbreak{}-{}\allowbreak{}r\allowbreak{}e\allowbreak{}v\allowbreak{}i\allowbreak{}e\allowbreak{}w\allowbreak{}/\allowbreak{}p\allowbreak{}r\allowbreak{}o\allowbreak{}b\allowbreak{}e\allowbreak{}s\allowbreak{}/\allowbreak{}c\allowbreak{}o\allowbreak{}d\allowbreak{}e\allowbreak{}-{}\allowbreak{}s\allowbreak{}p\allowbreak{}i\allowbreak{}n\allowbreak{}e\allowbreak{}/}, run from the repository root (leanerVM \texttt{main} at \texttt{b435631}, ArkLib \texttt{dca90385}) with

\begin{shellcode}
flock .claude/reports/blueprint-review/logs/lean.lock lake env lean .claude/reports/blueprint-review/probes/code-spine/<File>.lean
\end{shellcode}

The output of each run is the file \texttt{<File>.\allowbreak{}out} beside it, reproduced below the source. An empty output with \texttt{exit=0} means every \texttt{\#guard}, \texttt{example} and theorem of the file was accepted.

\subsubsection{\texttt{P1Axioms.lean}: axioms and computability of the load-bearing declarations}\label{sec:gen-probe-spine-P1Axioms}

\fhead{What it establishes}Every stated theorem of the spine is proved with the kernel's three standard axioms (\texttt{propext}, \texttt{Classical.\allowbreak{}choice}, \texttt{Quot.\allowbreak{}sound}), with no \texttt{sorryAx}. \texttt{Lean.\allowbreak{}isNoncomputable} is \texttt{false} for \texttt{piopExtractor}, \texttt{commitExtractor}, \texttt{Component.\allowbreak{}sendExtractor}, \texttt{Phases.\allowbreak{}Security.\allowbreak{}toDef}, \texttt{Component.\allowbreak{}Security.\allowbreak{}append}, ArkLib's \texttt{E\allowbreak{}x\allowbreak{}t\allowbreak{}r\allowbreak{}a\allowbreak{}c\allowbreak{}t\allowbreak{}o\allowbreak{}r\allowbreak{}.\allowbreak{}R\allowbreak{}o\allowbreak{}u\allowbreak{}n\allowbreak{}d\allowbreak{}B\allowbreak{}y\allowbreak{}R\allowbreak{}o\allowbreak{}u\allowbreak{}n\allowbreak{}d\allowbreak{}.\allowbreak{}a\allowbreak{}p\allowbreak{}p\allowbreak{}e\allowbreak{}n\allowbreak{}d}, \texttt{C\allowbreak{}o\allowbreak{}m\allowbreak{}p\allowbreak{}o\allowbreak{}n\allowbreak{}e\allowbreak{}n\allowbreak{}t\allowbreak{}.\allowbreak{}p\allowbreak{}a\allowbreak{}s\allowbreak{}s\allowbreak{}T\allowbreak{}h\allowbreak{}r\allowbreak{}o\allowbreak{}u\allowbreak{}g\allowbreak{}h\allowbreak{}E\allowbreak{}x\allowbreak{}t\allowbreak{}r\allowbreak{}a\allowbreak{}c\allowbreak{}t\allowbreak{}o\allowbreak{}r}, \texttt{PublicInput.\allowbreak{}extractor} and \texttt{publicInputSecurity}, and \texttt{true} for \texttt{publicInputPhase} only (its error is a real number). \texttt{\#print axioms} reports \texttt{Classical.\allowbreak{}choice} for computable definitions too, because the proofs inside them (the casts between witness types) use classical lemmas: it is no evidence of computability (D.5).

\fhead{Source}\file{.claude/reports/blueprint-review/probes/code-spine/P1Axioms.lean}

\begin{leancode}
import LeanerVM.Protocol.Spine.Toy
import LeanerVM.Protocol.Spine.Compose
import LeanerVM.Protocol.ToArkLib.Refinement
import LeanerVM.Protocol.PublicInput

/-!
Probe P1 (code-spine): axioms and computability of the spine's load-bearing declarations.
-/

open LeanerVM.Parameters LeanerVM.Protocol LeanerVM.Protocol.Toy CompPoly CPoly OracleComp

#print axioms piop_perfectCompleteness
#print axioms piop_rbrKnowledgeSoundness
#print axioms piop_rbrKnowledgeSoundness_exists
#print axioms piopExtractor
#print axioms commitExtractor
#print axioms Component.sendExtractor
#print axioms commitSecurity
#print axioms commitComplete
#print axioms Component.Security.append
#print axioms Component.Complete.append
#print axioms Phases.Security.toDef
#print axioms M3Holds
#print axioms Component.passThroughSecurity
#print axioms publicInputSecurity
#print axioms Verifier.append_rbrKnowledgeSoundnessWorstCaseWith_of_guarded_first
#print axioms Verifier.KnowledgeStateFunction.appendGuarded
#print axioms Extractor.RoundByRound.append
#print axioms Refinement.map_option_valid
#print axioms Toy.read_eval

open Lean Elab Command in
run_cmd do
  let env ← getEnv
  for n in [``piopExtractor, ``commitExtractor, ``Component.sendExtractor,
      ``Phases.Security.toDef, ``Phases.Complete.toDef, ``Component.Security.append,
      ``Component.Complete.append, ``Extractor.RoundByRound.append, ``leanVmPiop,
      ``leanVmVerifier, ``leanVmProver, ``piopError, ``Phases.toDef, ``Component.Def.append,
      ``commitDef, ``Component.sendOracle, ``Component.sendProver, ``Component.sendVerifier,
      ``publicInputPhase, ``publicInputSecurity, ``publicInputComplete, ``commitSecurity,
      ``commitComplete, ``Component.passThrough, ``Component.passThroughSecurity,
      ``Component.passThroughExtractor, ``PublicInput.extractor, ``PublicInput.prover,
      ``PublicInput.verifier, ``Verifier.KnowledgeStateFunction.appendGuarded,
      ``Component.guardedAppend, ``M3Holds, ``Toy.toy, ``Toy.honest,
      ``Extractor.Straightline.map, ``OracleReduction.append] do
    logInfo m!"{n}: noncomputable={Lean.isNoncomputable env n}"
\end{leancode}

\fhead{Command}run from the repository root, under the shared lock:

\begin{shellcode}
flock .claude/reports/blueprint-review/logs/lean.lock lake env lean .claude/reports/blueprint-review/probes/code-spine/P1Axioms.lean
\end{shellcode}

\fhead{Recorded output}

\begin{srccode}
'LeanerVM.Protocol.piop_perfectCompleteness' depends on axioms: [propext, Classical.choice, Quot.sound]
'LeanerVM.Protocol.piop_rbrKnowledgeSoundness' depends on axioms: [propext, Classical.choice, Quot.sound]
'LeanerVM.Protocol.piop_rbrKnowledgeSoundness_exists' depends on axioms: [propext, Classical.choice, Quot.sound]
'LeanerVM.Protocol.piopExtractor' depends on axioms: [propext, Classical.choice, Quot.sound]
'LeanerVM.Protocol.commitExtractor' depends on axioms: [propext, Classical.choice, Quot.sound]
'LeanerVM.Protocol.Component.sendExtractor' depends on axioms: [propext, Quot.sound]
'LeanerVM.Protocol.commitSecurity' depends on axioms: [propext, Classical.choice, Quot.sound]
'LeanerVM.Protocol.commitComplete' depends on axioms: [propext, Classical.choice, Quot.sound]
'LeanerVM.Protocol.Component.Security.append' depends on axioms: [propext, Classical.choice, Quot.sound]
'LeanerVM.Protocol.Component.Complete.append' depends on axioms: [propext, Classical.choice, Quot.sound]
'LeanerVM.Protocol.Phases.Security.toDef' depends on axioms: [propext, Classical.choice, Quot.sound]
'LeanerVM.Protocol.M3Holds' depends on axioms: [propext, Classical.choice, Quot.sound]
'LeanerVM.Protocol.Component.passThroughSecurity' depends on axioms: [propext, Classical.choice, Quot.sound]
'LeanerVM.Protocol.publicInputSecurity' depends on axioms: [propext, Classical.choice, Quot.sound]
'LeanerVM.Protocol.Verifier.append_rbrKnowledgeSoundnessWorstCaseWith_of_guarded_first' depends on axioms: [propext,
 Classical.choice,
 Quot.sound]
'LeanerVM.Protocol.Verifier.KnowledgeStateFunction.appendGuarded' depends on axioms: [propext,
 Classical.choice,
 Quot.sound]
'Extractor.RoundByRound.append' depends on axioms: [propext, Classical.choice, Quot.sound]
'LeanerVM.Protocol.Refinement.map_option_valid' depends on axioms: [propext, Quot.sound]
'LeanerVM.Protocol.Toy.read_eval' depends on axioms: [propext, Classical.choice, Quot.sound]
LeanerVM.Protocol.piopExtractor: noncomputable=false
LeanerVM.Protocol.commitExtractor: noncomputable=false
LeanerVM.Protocol.Component.sendExtractor: noncomputable=false
LeanerVM.Protocol.Phases.Security.toDef: noncomputable=false
LeanerVM.Protocol.Phases.Complete.toDef: noncomputable=false
LeanerVM.Protocol.Component.Security.append: noncomputable=false
LeanerVM.Protocol.Component.Complete.append: noncomputable=false
Extractor.RoundByRound.append: noncomputable=false
LeanerVM.Protocol.leanVmPiop: noncomputable=false
LeanerVM.Protocol.leanVmVerifier: noncomputable=false
LeanerVM.Protocol.leanVmProver: noncomputable=false
LeanerVM.Protocol.piopError: noncomputable=false
LeanerVM.Protocol.Phases.toDef: noncomputable=false
LeanerVM.Protocol.Component.Def.append: noncomputable=false
LeanerVM.Protocol.commitDef: noncomputable=false
LeanerVM.Protocol.Component.sendOracle: noncomputable=false
LeanerVM.Protocol.Component.sendProver: noncomputable=false
LeanerVM.Protocol.Component.sendVerifier: noncomputable=false
LeanerVM.Protocol.publicInputPhase: noncomputable=true
LeanerVM.Protocol.publicInputSecurity: noncomputable=false
LeanerVM.Protocol.publicInputComplete: noncomputable=false
LeanerVM.Protocol.commitSecurity: noncomputable=false
LeanerVM.Protocol.commitComplete: noncomputable=false
LeanerVM.Protocol.Component.passThrough: noncomputable=false
LeanerVM.Protocol.Component.passThroughSecurity: noncomputable=false
LeanerVM.Protocol.Component.passThroughExtractor: noncomputable=false
LeanerVM.Protocol.PublicInput.extractor: noncomputable=false
LeanerVM.Protocol.PublicInput.prover: noncomputable=false
LeanerVM.Protocol.PublicInput.verifier: noncomputable=false
LeanerVM.Protocol.Verifier.KnowledgeStateFunction.appendGuarded: noncomputable=false
LeanerVM.Protocol.Component.guardedAppend: noncomputable=false
LeanerVM.Protocol.M3Holds: noncomputable=false
LeanerVM.Protocol.Toy.toy: noncomputable=false
LeanerVM.Protocol.Toy.honest: noncomputable=false
LeanerVM.Protocol.Extractor.Straightline.map: noncomputable=false
OracleReduction.append: noncomputable=false
exit=0
\end{srccode}

\subsubsection{\texttt{P2Relation.lean}: \texttt{M3Holds} on variants of the toy instance}\label{sec:gen-probe-spine-P2Relation}

\fhead{What it establishes}What the repository's tests leave out (D.1, D.4): the two multiplicity cases (\texttt{toyOnce}, a tuple pushed twice and pulled once; \texttt{toyNever}, pushed twice and never pulled, where a balance summed in the field accepts), each failing \texttt{Balanced} alone; the side filter (\texttt{toySwap}); the row ranges; the cells \texttt{CountsNonzero} ranges over; and a layout that reads every column from the same two cells (\texttt{toyAlias}, a \texttt{Layout}, on which the exhaustive guard finds no satisfying stack). The probe uses the numeral \texttt{(\allowbreak{}2 : K)}, which at the old CompPoly pin is the polynomial \texttt{x} and at the new pin is \texttt{0}; it must be re-run with \texttt{K.\allowbreak{}ofBits 2} after a rebuild.

\fhead{Source}\file{.claude/reports/blueprint-review/probes/code-spine/P2Relation.lean}

\begin{leancode}
import LeanerVM.Protocol.Spine.Toy
import LeanerVM.Protocol.Spine.Compose

/-!
Probe P2 (code-spine): `M3Holds` on variants of the toy instance.

* multiplicity: a tuple pushed twice and pulled once, a tuple pushed twice and never pulled;
* the side filter: the same instance with the two sides exchanged;
* the layout's freedom: every column read from the same two cells.
-/

open LeanerVM.Parameters LeanerVM.Protocol LeanerVM.Protocol.Toy CompPoly CPoly OracleComp

namespace Probe

/-! ## Multiplicity -/

/-- A boundary block of height one (`κ = 0`): pulls the single tuple `(1, 0, …)`. -/
def boundaryOnce : BoundaryBlock shape where
  κ := 0
  side := .pull
  coords := Vector.ofFn fun k ↦ if k.val = 0 then .const 1 else .const 0

/-- The toy whose boundary pulls `(1, 0, …)` once, where the table pushes it twice. -/
abbrev toyOnce : M3Instance := { toy with boundary := [boundaryOnce] }

/-- The toy with no boundary: the table pushes `(1, 0, …)` twice and nothing is pulled. -/
abbrev toyNever : M3Instance := { toy with boundary := [] }

-- Pushed twice, pulled once: the two sides are the same *set* and differ as multisets.
#guard (toyOnce.tuples honest .push).length = 2
#guard (toyOnce.tuples honest .pull).length = 1
#guard (toyOnce.tuples honest .push).map (fun t ↦ t.get 0) = [1, 1]
#guard (toyOnce.tuples honest .pull).map (fun t ↦ t.get 0) = [1]
#guard (toyOnce.tuples honest .push).all fun t ↦ (toyOnce.tuples honest .pull).any fun u ↦
  t.toList == u.toList
#guard (toyOnce.tuples honest .pull).all fun t ↦ (toyOnce.tuples honest .push).any fun u ↦
  t.toList == u.toList
-- Only the balance clause fails.
#guard toyOnce.ConstraintsVanish honest
#guard ¬ toyOnce.Balanced honest
#guard toyOnce.CountsNonzero honest
#guard toyOnce.PublicLinesHold (1 : K) honest
#guard ¬ M3Holds toyOnce (1 : K) honest

-- Pushed twice, never pulled: every coordinate sums to zero in `K` on both sides (the sum of
-- the empty side is zero), so a balance summed in the field accepts; the multiset balance
-- rejects, and only it fails.
#guard (toyNever.tuples honest .push).length = 2
#guard (toyNever.tuples honest .pull).length = 0
#guard (List.finRange 16).all fun k ↦
  ((toyNever.tuples honest .push).map fun t ↦ t.get k).sum =
    ((toyNever.tuples honest .pull).map fun t ↦ t.get k).sum
#guard toyNever.ConstraintsVanish honest
#guard ¬ toyNever.Balanced honest
#guard toyNever.CountsNonzero honest
#guard toyNever.PublicLinesHold (1 : K) honest
#guard ¬ M3Holds toyNever (1 : K) honest

/-! ## The side filter -/

/-- The toy's flush, pulled instead of pushed. -/
def flushPull : Side × Vector (CMvPolynomial 3 K) 16 := (.pull, Toy.flush.2)

/-- The toy's boundary block, pushed instead of pulled. -/
def boundaryPush : BoundaryBlock shape := { Toy.boundary with side := .push }

/-- The toy with the two sides exchanged. -/
abbrev toySwap : M3Instance :=
  { toy with
    flushes := fun _ ↦ [flushPull]
    flushes_degree := fun _ f h k ↦ by
      rw [List.mem_singleton.mp h]
      exact toy.flushes_degree 0 Toy.flush (List.mem_singleton_self _) k
    boundary := [boundaryPush] }

#guard (toySwap.tuples honest .push).length = 2
#guard (toySwap.tuples honest .pull).length = 2
#guard (toySwap.flushTuples honest .push).length = 0
#guard (toySwap.flushTuples honest .pull).length = 2
#guard (toySwap.boundaryTuples honest .push).length = 2
#guard (toySwap.boundaryTuples honest .pull).length = 0
#guard M3Holds toySwap (1 : K) honest
#guard ¬ toySwap.Balanced (⟨#v[1, 0, 1, 1, 1, 0, 0, 0]⟩ : Column 3)

/-! ## The row ranges -/

-- One pushed tuple per row of the table (2 rows), one pulled tuple per row of the block
-- (2 rows); a block of height one contributes one.
#guard (toy.flushTuples honest .push).length = 2 ^ toy.τ 0
#guard (toy.boundaryTuples honest .pull).length = 2 ^ Toy.boundary.κ
#guard (toyOnce.boundaryTuples honest .pull).length = 2 ^ boundaryOnce.κ

-- The count clause ranges over every row of the count column and over nothing else: a zero
-- in the padding (cells 6, 7) or in another column's cell does not fail it.
#guard toy.CountsNonzero honest
#guard toy.CountsNonzero (⟨#v[0, 0, 1, 1, 0, 0, 0, 0]⟩ : Column 3)
#guard ¬ toy.CountsNonzero (⟨#v[1, 1, 0, 1, 1, 0, 1, 1]⟩ : Column 3)
#guard ¬ toy.CountsNonzero (⟨#v[1, 1, 1, 0, 1, 0, 1, 1]⟩ : Column 3)

/-! ## The layout's freedom: every column read from the same cells -/

/-- Every column is the slice at cells `0, 1`. -/
def aliasSlice (q : Column 3) (_ : Col) : Column 1 :=
  ⟨#v[q.values.get ⟨0, by decide⟩, q.values.get ⟨1, by decide⟩]⟩

/-- Every column's selector is `(0, 0)`. -/
def aliasExtend (_ : Col) (z : Vector E 1) : Vector E 3 := #v[z.head, 0, 0]

/-- The reading law holds of the aliased layout. -/
theorem alias_read_eval (q : Column 3) (c : Col) (z : Vector E 1) :
    CMlPolynomialEval.eval₂Mle (aliasSlice q c).values (algebraMap K E) z =
      CMlPolynomialEval.eval₂Mle q.values (algebraMap K E) (aliasExtend c z) := by
  simp [CMlPolynomialEval.eval₂Mle, CMlPolynomialEval.evalMle,
    CMlPolynomialEval.evalMleValues, CMlPolynomialEval.evalMleStep, CMlPolynomialEval.map,
    Vector.head, Vector.tail, aliasExtend, aliasSlice]

/-- The aliased layout: a `Layout`, since the law is a reading law. -/
def aliasLayout : Layout 3 Col (fun _ ↦ 1) where
  read := aliasSlice
  extend := aliasExtend
  read_eval := alias_read_eval

/-- The toy with every column read from cells `0, 1`. -/
abbrev toyAlias : M3Instance := { toy with layout := aliasLayout }

#guard (toyAlias.column honest ⟨0, 0⟩).values.toList = (toyAlias.column honest ⟨0, 2⟩).values.toList
#guard ¬ M3Holds toyAlias (1 : K) honest
-- Balance needs cells `0, 1` to be `1, 1`; the public line needs cell `1` to be `0`.
#guard toyAlias.Balanced honest
#guard ¬ toyAlias.PublicLinesHold (1 : K) honest
#guard toyAlias.PublicLinesHold (1 : K) (⟨#v[1, 0, 1, 1, 1, 0, 0, 0]⟩ : Column 3)
#guard ¬ toyAlias.Balanced (⟨#v[1, 0, 1, 1, 1, 0, 0, 0]⟩ : Column 3)

-- Evidence (not a proof) that the relation of the aliased toy is empty: no stack with cells
-- `0, 1` in `{0, 1, 2}` satisfies it at a statement in `{0, 1, 2}`. On paper: balance forces
-- cells `0, 1` to be `1, 1`, and the public line on column 2, now cells `0, 1`, forces cell 1
-- to be `0`.
#guard ([0, 1, 2] : List K).all fun a ↦ ([0, 1, 2] : List K).all fun b ↦
  ([0, 1, 2] : List K).all fun s ↦
    decide (¬ M3Holds toyAlias s (⟨#v[a, b, 1, 1, 1, 0, 0, 0]⟩ : Column 3))

end Probe
\end{leancode}

\fhead{Command}run from the repository root, under the shared lock:

\begin{shellcode}
flock .claude/reports/blueprint-review/logs/lean.lock lake env lean .claude/reports/blueprint-review/probes/code-spine/P2Relation.lean
\end{shellcode}

\fhead{Recorded output}

\begin{srccode}
exit=0
\end{srccode}

\subsubsection{\texttt{P3aSeams.lean}: each seam, an inhabitant and its near misses}\label{sec:gen-probe-spine-P3aSeams}

\fhead{What it establishes}Each seam on the honest toy stack has an inhabitant and each of its conjuncts a near miss (the table of D.2): among them a true cubic claim rejected by the degree conjunct of \texttt{Seam.\allowbreak{}bus}, the wrong statement inside \texttt{Seam.\allowbreak{}pub}, and \texttt{Seam.\allowbreak{}done toy = Set.\allowbreak{}univ}. Two facts of the design follow: a seam does not imply the relation before it (the zerocheck escape on \texttt{badConstraint}), and a seam does not say which claims are emitted. The repository gives the claims no decision procedure; the probe supplies them by \texttt{inferInstanceAs} on the unfolded propositions. It uses \texttt{(\allowbreak{}2 : K)} and must be re-run with \texttt{K.\allowbreak{}ofBits 2} after a rebuild.

\fhead{Source}\file{.claude/reports/blueprint-review/probes/code-spine/P3aSeams.lean}

\begin{leancode}
import LeanerVM.Protocol.Spine.Toy
import LeanerVM.Protocol.Spine.Compose
import LeanerVM.Protocol.PublicInput

/-!
Probe P3a (code-spine): each seam on the toy instance, an inhabitant and near misses; the
zerocheck escape; the commit phase's state function and extractor; the type of the composed
extractor's output when a phase has rounds.
-/

open LeanerVM.Parameters LeanerVM.Protocol LeanerVM.Protocol.Toy CompPoly CPoly OracleComp

namespace Probe

/-! ## Decision procedures for the claims and the seams (the repository has none)

Spelled with `inferInstanceAs` on the unfolded proposition, as the repository's public-input
test does: the instance `by unfold …; infer_instance` elaborates and then exhausts the memory
under `#guard` (observed: exit 137 on the first version of this probe). -/

instance {I : M3Instance} (q : Column I.μ) (c : ColumnClaim I) : Decidable (c.Holds q) :=
  inferInstanceAs (Decidable (CMlPolynomialEval.eval₂Mle (I.column q c.col).values
    (algebraMap K E) c.point = c.value))

instance {I : M3Instance} (q : Column I.μ) (c : LinearClaim I) : Decidable (c.Holds q) :=
  inferInstanceAs (Decidable ((c.terms.map fun t ↦ t.eval q).sum = c.value))

instance {I : M3Instance} (q : Column I.μ) (c : WeightedClaim I) : Decidable (c.Holds q) :=
  inferInstanceAs (Decidable (c.weight.pair q = c.value))

/-- The predicate of the bus seam, on a statement and a stack. -/
def busPred (I : M3Instance) (s : I.Stmt × BusOut I) (q : Column I.μ) : Prop :=
  (∀ c ∈ s.2.linear, c.Holds q) ∧
  (∀ c ∈ s.2.linear, ∀ t ∈ c.terms, t.poly.totalDegree ≤ I.d) ∧
  (∀ c ∈ s.2.columns, c.Holds q) ∧ I.PublicLinesHold s.1 q ∧ I.aux q

/-- It is the bus seam, by definition. -/
example (I : M3Instance) (s : I.Stmt × BusOut I) (o : ∀ i, TheOracle I i) :
    ((s, o), ()) ∈ Seam.bus I ↔ busPred I s (o 0) := Iff.rfl

instance (I : M3Instance) (s : I.Stmt × BusOut I) (q : Column I.μ) :
    Decidable (busPred I s q) :=
  inferInstanceAs (Decidable ((∀ c ∈ s.2.linear, c.Holds q) ∧
    (∀ c ∈ s.2.linear, ∀ t ∈ c.terms, t.poly.totalDegree ≤ I.d) ∧
    (∀ c ∈ s.2.columns, c.Holds q) ∧ I.PublicLinesHold s.1 q ∧ I.aux q))

def tablePred (I : M3Instance) (s : I.Stmt × TableOut I) (q : Column I.μ) : Prop :=
  (∀ c ∈ s.2.columns, c.Holds q) ∧ I.PublicLinesHold s.1 q ∧ I.aux q

example (I : M3Instance) (s : I.Stmt × TableOut I) (o : ∀ i, TheOracle I i) :
    ((s, o), ()) ∈ Seam.table I ↔ tablePred I s (o 0) := Iff.rfl

instance (I : M3Instance) (s : I.Stmt × TableOut I) (q : Column I.μ) :
    Decidable (tablePred I s q) :=
  inferInstanceAs (Decidable ((∀ c ∈ s.2.columns, c.Holds q) ∧ I.PublicLinesHold s.1 q ∧
    I.aux q))

def pubPred (I : M3Instance) (s : I.Stmt × PubOut I) (q : Column I.μ) : Prop :=
  (∀ c ∈ s.2.columns, c.Holds q) ∧ I.aux q

example (I : M3Instance) (s : I.Stmt × PubOut I) (o : ∀ i, TheOracle I i) :
    ((s, o), ()) ∈ Seam.pub I ↔ pubPred I s (o 0) := Iff.rfl

instance (I : M3Instance) (s : I.Stmt × PubOut I) (q : Column I.μ) :
    Decidable (pubPred I s q) :=
  inferInstanceAs (Decidable ((∀ c ∈ s.2.columns, c.Holds q) ∧ I.aux q))

def flockPred (I : M3Instance) (s : I.Stmt × FlockOut I) (q : Column I.μ) : Prop :=
  (∀ c ∈ s.2.columns, c.Holds q) ∧ ∀ c ∈ s.2.weighted, c.Holds q

example (I : M3Instance) (s : I.Stmt × FlockOut I) (o : ∀ i, TheOracle I i) :
    ((s, o), ()) ∈ Seam.flock I ↔ flockPred I s (o 0) := Iff.rfl

instance (I : M3Instance) (s : I.Stmt × FlockOut I) (q : Column I.μ) :
    Decidable (flockPred I s q) :=
  inferInstanceAs (Decidable ((∀ c ∈ s.2.columns, c.Holds q) ∧ ∀ c ∈ s.2.weighted, c.Holds q))

/-! ## The claims of the probe -/

def eZero : E := 0
def eOne : E := 1
/-- A point of `E` outside `{0, 1}`: the image of `2 : K` (the polynomial `x`). -/
def eTwo : E := ofK 2

/-- Column 0 of the honest stack is `[1, 1]`: its extension is `1` everywhere. -/
def col0One : ColumnClaim toy := ⟨⟨0, 0⟩, #v[eTwo], eOne⟩
def col0Zero : ColumnClaim toy := ⟨⟨0, 0⟩, #v[eTwo], eZero⟩

/-- The zerocheck claim of the toy's constraint at the point `(0)`: value `0`. -/
def zeroAt0 : LinearClaim toy :=
  ⟨[⟨eOne, 0, CMvPolynomial.X 2 * CMvPolynomial.X 2 - CMvPolynomial.X 2, #v[eZero]⟩], eZero⟩
/-- The same at the point `(1)`. -/
def zeroAt1 : LinearClaim toy :=
  ⟨[⟨eOne, 0, CMvPolynomial.X 2 * CMvPolynomial.X 2 - CMvPolynomial.X 2, #v[eOne]⟩], eZero⟩
/-- The same with the wrong value. -/
def zeroWrong : LinearClaim toy :=
  ⟨[⟨eOne, 0, CMvPolynomial.X 2 * CMvPolynomial.X 2 - CMvPolynomial.X 2, #v[eZero]⟩], eOne⟩
/-- A cubic term, whose claim is true of the honest stack (column 2 is `[1, 0]`, so `X₂³` is
`[1, 0]` and its extension at `(0)` is `1`). -/
def cubicTrue : LinearClaim toy :=
  ⟨[⟨eOne, 0, CMvPolynomial.X 2 * CMvPolynomial.X 2 * CMvPolynomial.X 2, #v[eZero]⟩], eOne⟩

/-- The stack of the repository's test `badConstraint`: column 2 is `[2, 0]`, not Boolean. -/
def badConstraint : Column 3 := ⟨#v[1, 1, 1, 1, 2, 0, 0, 0]⟩

/-! ## The bus seam -/

-- An inhabitant: the zerocheck claim and a column claim, on the honest stack at statement 1.
#guard busPred toy ((1 : K), ⟨[zeroAt0, zeroAt1], [col0One]⟩) honest
-- Near misses, one per conjunct.
#guard ¬ busPred toy ((1 : K), ⟨[zeroWrong], [col0One]⟩) honest          -- a false linear claim
#guard ¬ busPred toy ((1 : K), ⟨[cubicTrue], [col0One]⟩) honest          -- a true cubic claim
#guard cubicTrue.Holds honest                                            -- (it is true)
#guard ¬ busPred toy ((1 : K), ⟨[zeroAt0], [col0Zero]⟩) honest           -- a false column claim
#guard ¬ busPred toy ((0 : K), ⟨[zeroAt0], [col0One]⟩) honest            -- the wrong statement

-- The zerocheck escape: on `badConstraint` the constraint is violated on the cube (row 0), so
-- the claim at the point `(0)` is false, and the claim at the point `(1)` is true. The bus seam
-- holds of a stack outside `M3Holds`; the bus phase's error is what pays for it.
#guard ¬ toy.ConstraintsVanish badConstraint
#guard ¬ zeroAt0.Holds badConstraint
#guard zeroAt1.Holds badConstraint
#guard busPred toy ((2 : K), ⟨[zeroAt1], [col0One]⟩) badConstraint
#guard ¬ M3Holds toy (2 : K) badConstraint

-- A bus statement with no claim holds of every stack whose public line holds: the seam does
-- not say which claims are emitted.
#guard busPred toy ((2 : K), ⟨[], []⟩) badConstraint

/-! ## The table and public-input seams -/

#guard tablePred toy ((1 : K), ⟨[col0One]⟩) honest
#guard ¬ tablePred toy ((1 : K), ⟨[col0Zero]⟩) honest
#guard ¬ tablePred toy ((0 : K), ⟨[col0One]⟩) honest
#guard pubPred toy ((1 : K), ⟨[col0One]⟩) honest
#guard ¬ pubPred toy ((1 : K), ⟨[col0Zero]⟩) honest
-- The public seam no longer sees the statement: the wrong statement is inside it.
#guard pubPred toy ((0 : K), ⟨[col0One]⟩) honest

/-! ## The Flock seam -/

/-- The weight that selects cell 0 of the stack. -/
def cell0Weight : Weight 3 where
  onCube := Vector.ofFn fun i ↦ if i.val = 0 then 1 else 0
  mle := fun r ↦ CMlPolynomialEval.evalMle (Vector.ofFn fun i ↦ if i.val = 0 then 1 else 0) r
  mle_eq := fun _ ↦ rfl

def w0One : WeightedClaim toy := ⟨cell0Weight, eOne⟩
def w0Zero : WeightedClaim toy := ⟨cell0Weight, eZero⟩

#guard flockPred toy ((1 : K), ⟨[col0One], [w0One]⟩) honest
#guard ¬ flockPred toy ((1 : K), ⟨[col0One], [w0Zero]⟩) honest
#guard ¬ flockPred toy ((1 : K), ⟨[col0Zero], [w0One]⟩) honest

/-! ## The last seam has no near miss: it is the whole set -/

example : Seam.done toy = Set.univ := Set.eq_univ_of_forall fun _ ↦ trivial

end Probe
\end{leancode}

\fhead{Command}run from the repository root, under the shared lock:

\begin{shellcode}
flock .claude/reports/blueprint-review/logs/lean.lock lake env lean .claude/reports/blueprint-review/probes/code-spine/P3aSeams.lean
\end{shellcode}

\fhead{Recorded output}

\begin{srccode}
exit=0
\end{srccode}

\subsubsection{\texttt{P3bCommit.lean}: the commit phase's state function and extractor}\label{sec:gen-probe-spine-P3bCommit}

\fhead{What it establishes}Every statement by \texttt{Iff.\allowbreak{}rfl} or \texttt{rfl}: the commit phase's knowledge state function is \texttt{M3Holds} of the candidate witness before the message and \texttt{M3Holds} of the message after it, whatever the candidate witness; both maps of its extractor return the message (D.6).

\fhead{Source}\file{.claude/reports/blueprint-review/probes/code-spine/P3bCommit.lean}

\begin{leancode}
import LeanerVM.Protocol.Spine.Toy
import LeanerVM.Protocol.Spine.Compose
import LeanerVM.Protocol.PublicInput

/-!
Probe P3b (code-spine): each seam on the toy instance, an inhabitant and near misses; the
zerocheck escape; the commit phase's state function and extractor; the type of the composed
extractor's output when a phase has rounds.
-/

open LeanerVM.Parameters LeanerVM.Protocol LeanerVM.Protocol.Toy CompPoly CPoly OracleComp

namespace Probe

/-! ## The commit phase: state function and extractor -/

section Commit

variable {σ : Type} (init : ProbComp σ) (impl : QueryImpl []ₒ (StateT σ ProbComp))

-- Before the message: `M3Holds` of the candidate witness.
example (s : K) (o : ∀ i, NoOracle i) (w : Column 3) :
    ((commitSecurity toy).kSF init impl).toFun 0 (s, o) default w ↔ M3Holds toy s w := Iff.rfl

-- After the message: `M3Holds` of the message, whatever the candidate witness.
example (s : K) (o : ∀ i, NoOracle i) (tr : (commitSpec toy).FullTranscript) (w : Column 3) :
    ((commitSecurity toy).kSF init impl).toFun (Fin.last 1) (s, o) tr w ↔
      M3Holds toy s (tr 0) := Iff.rfl

-- The extractor: both maps return the message.
example (s : K) (o : ∀ i, NoOracle i) (tr : (commitSpec toy).FullTranscript) :
    (commitExtractor toy).extractOut (s, o) tr () = tr 0 := rfl
example (s : K) (o : ∀ i, NoOracle i) (tr : (commitSpec toy).FullTranscript) (w : Column 3) :
    (commitExtractor toy).extractMid 0 (s, o) tr w = tr 0 := rfl

-- The commit phase's security is the one the composition starts from.
example : (commitSecurity toy).extractor = commitExtractor toy := rfl
example : (commitDef toy).n = 1 := rfl

end Commit

end Probe
\end{leancode}

\fhead{Command}run from the repository root, under the shared lock:

\begin{shellcode}
flock .claude/reports/blueprint-review/logs/lean.lock lake env lean .claude/reports/blueprint-review/probes/code-spine/P3bCommit.lean
\end{shellcode}

\fhead{Recorded output}

\begin{srccode}
exit=0
\end{srccode}

\subsubsection{\texttt{P3cExtractor.lean}: the composed extractor when a phase has rounds}\label{sec:gen-probe-spine-P3cExtractor}

\fhead{What it establishes}With the repository's own \texttt{publicInputPhase} in the bundle and four pass-through phases, the composed extractor's \texttt{extractOut} lives in the last phase's intermediate witness type (\texttt{Unit}); the stack is what \texttt{extractMid} returns at round 0, and there it is the transcript's first message whatever it is handed; the repository test's \texttt{rfl} statement is ill-typed for this bundle, which is the expected error of the output (D.5).

\fhead{Source}\file{.claude/reports/blueprint-review/probes/code-spine/P3cExtractor.lean}

\begin{leancode}
import LeanerVM.Protocol.Spine.Toy
import LeanerVM.Protocol.Spine.Compose
import LeanerVM.Protocol.PublicInput

/-!
Probe P3c (code-spine): the type and the value of the composed extractor when a phase has
rounds. The last declaration is expected to fail.
-/

open LeanerVM.Parameters LeanerVM.Protocol LeanerVM.Protocol.Toy CompPoly CPoly OracleComp

namespace Probe

/-! ## The composed extractor when a phase has rounds -/

/-- The toy's phases with the real public-input phase (two rounds) and four pass-throughs. -/
noncomputable def realPub : Phases toy where
  bus := Phase.passThrough toy fun s ↦ (s, ⟨[], []⟩)
  table := Phase.passThrough toy fun p ↦ (p.1, ⟨[]⟩)
  pub := publicInputPhase toy
  flock := Phase.passThrough toy fun p ↦ (p.1, ⟨[], []⟩)
  opening := Phase.passThrough toy fun _ ↦ ()

example : realPub.toDef.n = 3 := rfl

section

variable (Sb : Phase.Security toy realPub.bus (Seam.commit toy) (Seam.bus toy))
  (St : Phase.Security toy realPub.table (Seam.bus toy) (Seam.table toy))
  (Sf : Phase.Security toy realPub.flock (Seam.pub toy) (Seam.flock toy))
  (So : Phase.Security toy realPub.opening (Seam.flock toy) (Seam.done toy))

/-- A security bundle whose public-input field is the repository's. -/
def realPubSecurity : realPub.Security := ⟨Sb, St, publicInputSecurity toy, Sf, So⟩

-- The output slot of the composed extractor is the public-input phase's last intermediate
-- witness, `Unit`: it is not the stack.
example : (realPubSecurity Sb St Sf So).toDef.witMid (Fin.last 3) = Unit := rfl

-- The stack is what the composed extractor returns at round 0, from the transcript's first
-- message.
example : (realPubSecurity Sb St Sf So).toDef.witMid 0 = Column 3 := rfl

-- At round 0 the composed extractor returns the transcript's first message, whatever the
-- phases' extractors and whatever intermediate witness it is handed.
example (s : K)
    (tr : realPub.toDef.pSpec.Transcript (Fin.succ (⟨0, Nat.zero_lt_succ 2⟩ : Fin 3)))
    (w : (realPubSecurity Sb St Sf So).toDef.witMid (Fin.succ (⟨0, Nat.zero_lt_succ 2⟩ : Fin 3))) :
    (piopExtractor realPub (realPubSecurity Sb St Sf So)).extractMid
      (⟨0, Nat.zero_lt_succ 2⟩ : Fin 3) (s, fun i : Fin 0 ↦ i.elim0) tr w =
        (tr ⟨0, Nat.zero_lt_succ 0⟩ : Column 3) := rfl

-- EXPECTED TO FAIL: the repository's test statement, with the real public-input phase in the
-- bundle. The output slot is `Unit`, so the statement "`extractOut` returns the stack" does not
-- typecheck.
example (tr : realPub.toDef.pSpec.FullTranscript) :
    (piopExtractor realPub (realPubSecurity Sb St Sf So)).extractOut
      ((1 : K), fun i : Fin 0 ↦ i.elim0) tr () = honest := rfl

end

end Probe
\end{leancode}

\fhead{Command}run from the repository root, under the shared lock:

\begin{shellcode}
flock .claude/reports/blueprint-review/logs/lean.lock lake env lean .claude/reports/blueprint-review/probes/code-spine/P3cExtractor.lean
\end{shellcode}

\fhead{Recorded output}

\begin{srccode}
.claude/reports/blueprint-review/probes/code-spine/P3cExtractor.lean:58:49: error: Type mismatch
  honest
has type
  Column 3
but is expected to have type
  (realPubSecurity Sb St Sf So).toDef.witMid (Fin.last realPub.toDef.n)
exit=1
\end{srccode}

\subsubsection{\texttt{P4Junk.lean}: five phases that check nothing, at declared error 1}\label{sec:gen-probe-spine-P4Junk}

\fhead{What it establishes}\texttt{Phases.\allowbreak{}Security toy} is inhabited by five phases that draw a challenge, check nothing and declare error \texttt{1} (\texttt{junkSecurity}); both master theorems apply to them; \texttt{piopError junkPhases i = 1} at each of the five challenges (\texttt{piopError\_\allowbreak{}junk}). Nothing in the construction uses the toy (D.3 (b)).

\fhead{Source}\file{.claude/reports/blueprint-review/probes/code-spine/P4Junk.lean}

\begin{leancode}
import LeanerVM.Protocol.Spine.Toy
import LeanerVM.Protocol.Spine.Compose
import LeanerVM.Protocol.ToArkLib.GuardedVerdict
import LeanerVM.Protocol.ToArkLib.KeepOracles

/-!
Probe P4 (code-spine): `Phases.Security` of the toy instance is inhabited by five phases that
check nothing, each with one wasted challenge whose declared error is `1`.

A phase `wasted f ε`: the verifier draws one challenge in `E`, ignores it, and maps the
statement by `f`. It is complete whenever `f` carries the input seam into the output seam
(the pass-through's completeness). With declared error `ε = 1` it is round-by-round knowledge
sound for *any* two seams: the state is the input seam before the challenge and `True` after
it, and a probability is at most `1`.
-/

open LeanerVM.Parameters LeanerVM.Protocol LeanerVM.Protocol.Toy CompPoly CPoly OracleComp
  OracleSpec ProtocolSpec
open scoped NNReal ENNReal

namespace Probe

variable (I : M3Instance) {StmtIn StmtOut : Type}

/-- One round: a challenge in `E`. -/
@[reducible]
def wSpec : ProtocolSpec 1 := ⟨!v[.V_to_P], !v[E]⟩

instance : ∀ i, OracleInterface (wSpec.Message i)
  | ⟨0, h⟩ => nomatch h

instance : ∀ i, SampleableType (wSpec.Challenge i)
  | ⟨0, _⟩ => (inferInstance : SampleableType E)

/-- The prover receives the challenge and outputs the mapped statement. -/
def wProver (f : StmtIn → StmtOut) :
    OracleProver []ₒ StmtIn (TheOracle I) Unit StmtOut (TheOracle I) Unit wSpec where
  PrvState := fun _ ↦ (StmtIn × ∀ i, TheOracle I i) × Unit
  input := _root_.id
  receiveChallenge
    | ⟨0, _⟩ => fun st ↦ pure fun _ ↦ st
  sendMessage
    | ⟨0, h⟩ => nomatch h
  output := fun st ↦ pure ((f st.1.1, st.1.2), ())

/-- The verifier ignores the challenge, queries nothing and maps the statement. -/
def wVerifier (f : StmtIn → StmtOut) :
    OracleVerifier []ₒ StmtIn (TheOracle I) StmtOut (TheOracle I) wSpec where
  verify := fun s _ ↦ pure (f s)
  outputOracle := .inl (keepOracles (TheOracle I) wSpec)

theorem wVerifier_run (f : StmtIn → StmtOut) (s : StmtIn) (o : ∀ i, TheOracle I i)
    (tr : wSpec.FullTranscript) :
    (wVerifier I f).toVerifier.run (s, o) tr = pure (f s, o) := by
  simp only [Verifier.run, OracleVerifier.toVerifier]
  rw [OracleVerifier.materializeOutput_of_keepOracles _ rfl]
  rfl

/-- The verifier accepts every transcript. -/
def wPure (f : StmtIn → StmtOut) : (wVerifier I f).toVerifier.PureForm where
  verify := fun p _ ↦ (f p.1, p.2)
  verify_eq := fun ⟨s, o⟩ tr ↦ wVerifier_run I f s o tr

/-- The phase, with a declared error `ε` on its one challenge. -/
noncomputable def wasted (f : StmtIn → StmtOut) (ε : ℝ≥0) : Phase.Def I StmtIn StmtOut where
  n := 1
  pSpec := wSpec
  red := ⟨wProver I f, wVerifier I f⟩
  err := fun _ ↦ ε

theorem wProver_run_support (f : StmtIn → StmtOut) (s : StmtIn) (o : ∀ i, TheOracle I i)
    (pr : wSpec.FullTranscript × (StmtOut × ∀ i, TheOracle I i) × Unit)
    (hpr : pr ∈ support ((wProver I f).run (s, o) ())) : pr.2 = ((f s, o), ()) := by
  have h0 : wSpec.dir 0 = .V_to_P := rfl
  simp only [Prover.run, Prover.runToRound, Fin.induction_one,
    Prover.processRound_of_dir_eq_V_to_P 0 h0] at hpr
  simp only [ChallengeIdx, Challenge, wProver, id_eq, liftM_pure, bind_pure_comp, map_pure,
    pure_bind, Functor.map_map, support_map, Set.mem_image] at hpr
  obtain ⟨r, -, rfl⟩ := hpr
  rfl

variable {relIn : Set ((StmtIn × ∀ i, TheOracle I i) × Unit)}
  {relOut : Set ((StmtOut × ∀ i, TheOracle I i) × Unit)}

/-- Completeness, under the pass-through's hypothesis. -/
noncomputable def wastedComplete (f : StmtIn → StmtOut) (ε : ℝ≥0)
    (hc : ∀ s o, ((s, o), ()) ∈ relIn → ((f s, o), ()) ∈ relOut) :
    Phase.Complete I (wasted I f ε) relIn relOut where
  outputPure := ⟨_, fun _ ↦ rfl⟩
  guarded := (wPure I f).toGuardedForm
  complete := fun init impl ↦ by
    apply Reduction.perfectCompleteness_of_run_support
    intro stmtIn witIn hIn x hx
    obtain ⟨s, o⟩ := stmtIn
    obtain ⟨pr, hpr, rfl⟩ :=
      Reduction.mem_support_run_of_guarded _ (wPure I f).toGuardedForm (s, o) witIn hx
    have hout := wProver_run_support I f s o pr hpr
    have hc' : (wPure I f).toGuardedForm.check (s, o) pr.1 = true := rfl
    rw [if_pos hc']
    refine ⟨_, rfl, ?_, congrArg Prod.fst hout⟩
    have hw : pr.2.2 = () := rfl
    exact hw ▸ hc s o hIn

/-- The extractor keeps the trivial witness. -/
def wExtractor : Extractor.RoundByRound (OracleSpec.emptySpec.{0, 0})
    (StmtIn × ∀ i, TheOracle I i) Unit Unit wSpec (fun _ ↦ Unit) where
  eqIn := rfl
  extractMid := fun _ _ _ _ ↦ ()
  extractOut := fun _ _ _ ↦ ()

variable {σ : Type} (init : ProbComp σ) (impl : QueryImpl []ₒ (StateT σ ProbComp))

/-- A knowledge state function for ANY two seams: the input seam before the challenge, `True`
after it. No hypothesis on `f`, `relIn`, `relOut`. -/
def wState (f : StmtIn → StmtOut) :
    (wVerifier I f).toVerifier.KnowledgeStateFunction init impl relIn relOut
      (wExtractor I (StmtIn := StmtIn)) where
  toFun := fun m stmt _ _ ↦ if m.val = 0 then (stmt, ()) ∈ relIn else True
  toFun_empty := fun _ _ ↦ by simp
  toFun_next := fun m hm ↦ by
    fin_cases m
    exact absurd hm (by decide)
  toFun_full := fun _ _ _ _ ↦ by simp

/-- Round-by-round knowledge soundness at declared error `1`: a probability is at most one. -/
theorem w_rbr (f : StmtIn → StmtOut) :
    (wVerifier I f).toVerifier.rbrKnowledgeSoundnessWorstCaseWith init impl relIn relOut
      (fun _ ↦ Unit) (wExtractor I) (wState I init impl f) (fun _ ↦ 1) := by
  intro stmtIn i tr
  simp

/-- The security half at declared error `1`, under the completeness hypothesis alone. -/
noncomputable def wastedSecurity (f : StmtIn → StmtOut)
    (hc : ∀ s o, ((s, o), ()) ∈ relIn → ((f s, o), ()) ∈ relOut) :
    Phase.Security I (wasted I f 1) relIn relOut where
  toComplete := wastedComplete I f 1 hc
  witMid := fun _ ↦ Unit
  extractor := wExtractor I
  kSF := fun init impl ↦ wState I init impl f
  rbr := fun init impl ↦ w_rbr I init impl f

/-! ## The toy's five phases, checking nothing -/

/-- Five phases that draw a challenge and drop every claim. -/
noncomputable def junkPhases : Phases toy where
  bus := wasted toy (fun s ↦ (s, ⟨[], []⟩)) 1
  table := wasted toy (fun p ↦ (p.1, ⟨[]⟩)) 1
  pub := wasted toy (fun p ↦ (p.1, ⟨[]⟩)) 1
  flock := wasted toy (fun p ↦ (p.1, ⟨[], []⟩)) 1
  opening := wasted toy (fun _ ↦ ()) 1

/-- `Phases.Security` of the toy is inhabited by them. -/
noncomputable def junkSecurity : junkPhases.Security where
  bus := wastedSecurity toy _ fun _ _ h ↦ ⟨by simp, by simp, by simp, h.2.2.2.1, h.2.2.2.2⟩
  table := wastedSecurity toy _ fun _ _ h ↦ ⟨by simp, h.2.2.2.1, h.2.2.2.2⟩
  pub := wastedSecurity toy _ fun _ _ h ↦ ⟨by simp, h.2.2⟩
  flock := wastedSecurity toy _ fun _ _ _ ↦ ⟨by simp, by simp⟩
  opening := wastedSecurity toy _ fun _ _ _ ↦ trivial

example : junkPhases.toDef.n = 6 := rfl

/-- Their completeness halves. -/
noncomputable def junkComplete : junkPhases.Complete :=
  ⟨junkSecurity.bus.toComplete, junkSecurity.table.toComplete, junkSecurity.pub.toComplete,
    junkSecurity.flock.toComplete, junkSecurity.opening.toComplete⟩

/-- Both master theorems hold of the junk protocol. -/
example {σ : Type} (init : ProbComp σ) (impl : QueryImpl []ₒ (StateT σ ProbComp)) :
    (leanVmPiop junkPhases).perfectCompleteness init impl (M3Rel toy) (Seam.done toy) :=
  piop_perfectCompleteness junkPhases junkComplete init impl

example {σ : Type} (init : ProbComp σ) (impl : QueryImpl []ₒ (StateT σ ProbComp)) :
    (leanVmVerifier junkPhases).toVerifier.rbrKnowledgeSoundnessWorstCase init impl (M3Rel toy)
      (Seam.done toy) (piopError junkPhases) :=
  piop_rbrKnowledgeSoundness_exists junkPhases junkSecurity init impl

/-- Appending keeps a constant declared error. -/
theorem append_err_const {S₁ S₂ S₃ W₁ W₂ W₃ : Type} {ι₁ ι₂ ι₃ : Type} {O₁ : ι₁ → Type}
    {O₂ : ι₂ → Type} {O₃ : ι₃ → Type} [∀ i, OracleInterface (O₁ i)]
    [∀ i, OracleInterface (O₂ i)] [∀ i, OracleInterface (O₃ i)]
    (D₁ : Component.Def S₁ O₁ W₁ S₂ O₂ W₂) (D₂ : Component.Def S₂ O₂ W₂ S₃ O₃ W₃) (c : ℝ≥0)
    (h₁ : ∀ i, D₁.err i = c) (h₂ : ∀ i, D₂.err i = c) : ∀ i, (D₁.append D₂).err i = c := by
  intro i
  change Sum.elim D₁.err D₂.err (ChallengeIdx.sumEquiv.symm i) = c
  cases ChallengeIdx.sumEquiv.symm i with
  | inl a => exact h₁ a
  | inr b => exact h₂ b

/-- Every challenge of the junk protocol carries the declared error `1`. -/
theorem piopError_junk (i : junkPhases.toDef.pSpec.ChallengeIdx) : piopError junkPhases i = 1 := by
  unfold piopError Phases.toDef
  refine append_err_const _ _ 1 (append_err_const _ _ 1 (append_err_const _ _ 1
    (append_err_const _ _ 1 (append_err_const _ _ 1 ?_ ?_) ?_) ?_) ?_) ?_ i
  · rintro ⟨j, hj⟩
    fin_cases j
    exact absurd hj (by decide)
  all_goals intro _; rfl

#print axioms junkSecurity
#print axioms piopError_junk

end Probe
\end{leancode}

\fhead{Command}run from the repository root, under the shared lock:

\begin{shellcode}
flock .claude/reports/blueprint-review/logs/lean.lock lake env lean .claude/reports/blueprint-review/probes/code-spine/P4Junk.lean
\end{shellcode}

\fhead{Recorded output}

\begin{srccode}
'Probe.junkSecurity' depends on axioms: [propext, Classical.choice, Quot.sound]
'Probe.piopError_junk' depends on axioms: [propext, Classical.choice, Quot.sound]
exit=0
\end{srccode}

\subsubsection{\texttt{P5PassThrough.lean}: the pass-through bus phase has no security}\label{sec:gen-probe-spine-P5PassThrough}

\fhead{What it establishes}The reflection hypothesis of \texttt{Phase.\allowbreak{}passThroughSecurity} is false at the bus seam (\texttt{not\_\allowbreak{}reflects}), and the type \texttt{Phase.\allowbreak{}Security} of the pass-through bus phase is empty (\texttt{no\_\allowbreak{}security}), where the repository's test docstring only asserts it (D.3 (a)). The probe uses \texttt{(\allowbreak{}2 : K)} and must be re-run with \texttt{K.\allowbreak{}ofBits 2} after a rebuild.

\fhead{Source}\file{.claude/reports/blueprint-review/probes/code-spine/P5PassThrough.lean}

\begin{leancode}
import LeanerVM.Protocol.Spine.Toy
import LeanerVM.Protocol.Spine.Compose

/-!
Probe P5 (code-spine): the pass-through bus phase of the repository's test (`trivPhases.bus`)
has no `Phase.Security` against the spine's seams, on the toy instance.

`badConstraint` at statement `2` is outside `M3Holds` and the pass-through's output on it (no
claim) is inside the bus seam. A knowledge state function must be true at the end of an
accepting transcript and, the phase having no round, the end is the beginning, where the state
is the input seam.
-/

open LeanerVM.Parameters LeanerVM.Protocol LeanerVM.Protocol.Toy CompPoly CPoly OracleComp
  OracleSpec ProtocolSpec

namespace Probe

/-- The test's stack `badConstraint`: column 2 is `[2, 0]`. -/
def badConstraint : Column 3 := ⟨#v[1, 1, 1, 1, 2, 0, 0, 0]⟩

/-- The statement map of `trivPhases.bus`: no claim. -/
def dropAll (s : K) : K × BusOut toy := (s, ⟨[], []⟩)

/-- Outside the relation. -/
theorem bad_not_m3Holds : ¬ M3Holds toy (2 : K) badConstraint := by decide +kernel

/-- Inside the bus seam, with no claim. -/
theorem bad_mem_bus :
    ((dropAll 2, fun _ : Fin 1 ↦ badConstraint), ()) ∈ Seam.bus toy :=
  ⟨by simp [dropAll], by simp [dropAll], by simp [dropAll], by decide +kernel, trivial⟩

/-- The reflection hypothesis of `Phase.passThroughSecurity` is false at the bus seam. -/
theorem not_reflects : ¬ ∀ (s : K) (o : ∀ i, TheOracle toy i),
    ((dropAll s, o), ()) ∈ Seam.bus toy → ((s, o), ()) ∈ Seam.commit toy := fun h ↦
  bad_not_m3Holds (h 2 (fun _ ↦ badConstraint) bad_mem_bus)

/-- The empty implementation of the empty shared oracle. -/
def noImpl : QueryImpl []ₒ (StateT Unit ProbComp) := fun t ↦ PEmpty.elim t

/-- No `Phase.Security` exists for the pass-through bus phase: not only the repository's
constructor fails, the type is empty. -/
theorem no_security (S : Phase.Security toy (Phase.passThrough toy dropAll) (Seam.commit toy)
    (Seam.bus toy)) : False := by
  have ksf := S.kSF (pure ()) noImpl
  let stmt : K × ∀ i, TheOracle toy i := (2, fun _ ↦ badConstraint)
  let tr : (Phase.passThrough toy dropAll).pSpec.FullTranscript := fun i ↦ Fin.elim0 i
  have htr : tr = (default : (Phase.passThrough toy dropAll).pSpec.Transcript 0) :=
    funext fun i ↦ Fin.elim0 i
  have hpos : Pr[fun stmtOut ↦ (stmtOut, ()) ∈ Seam.bus toy | OptionT.mk do
      (simulateQ noImpl ((Phase.passThrough toy dropAll).red.verifier.toVerifier.run stmt
        tr)).run' (← (pure () : ProbComp Unit))] > 0 := by
    have hrun : (Phase.passThrough toy dropAll).red.verifier.toVerifier.run stmt tr =
        pure (dropAll 2, fun _ ↦ badConstraint) :=
      Component.passThroughVerifier_toVerifier_run (TheOracle toy) dropAll 2
        (fun _ ↦ badConstraint) tr
    rw [hrun]
    change Pr[_ | OptionT.mk (do let st ← (pure () : ProbComp Unit); (simulateQ noImpl
      (OptionT.run (pure (dropAll 2, fun _ : Fin 1 ↦ badConstraint)))).run' st)] > 0
    rw [OptionT.run_pure, simulateQ_pure]
    rw [gt_iff_lt, probEvent_pos_iff]
    refine ⟨(dropAll 2, fun _ ↦ badConstraint), ?_, bad_mem_bus⟩
    simp
  have hfull := ksf.toFun_full stmt tr () hpos
  rw [htr] at hfull
  exact bad_not_m3Holds ((ksf.toFun_empty stmt _).mpr hfull)

#print axioms no_security

end Probe
\end{leancode}

\fhead{Command}run from the repository root, under the shared lock:

\begin{shellcode}
flock .claude/reports/blueprint-review/logs/lean.lock lake env lean .claude/reports/blueprint-review/probes/code-spine/P5PassThrough.lean
\end{shellcode}

\fhead{Recorded output}

\begin{srccode}
'Probe.no_security' depends on axioms: [propext, Classical.choice, Quot.sound]
exit=0
\end{srccode}

\subsubsection{\texttt{P6Surface.lean}: three facts behind the audit surface and the transport}\label{sec:gen-probe-spine-P6Surface}

\fhead{What it establishes}The three facts of D.7: every \texttt{Component.\allowbreak{}Def} over the empty shared oracle has a pure prover output (\texttt{outputPure\_\allowbreak{}of\_\allowbreak{}def}), so the field \texttt{outputPure} is redundant; the two witness universes of \texttt{Refinement} are independent (\texttt{\#check @Refinement}); and the probabilistic transport of knowledge along a refinement holds in the event form, in three lines (\texttt{knowledge\_\allowbreak{}transport}).

\fhead{Source}\file{.claude/reports/blueprint-review/probes/code-spine/P6Surface.lean}

\begin{leancode}
import LeanerVM.Protocol.ToArkLib.Component
import LeanerVM.Protocol.ToArkLib.Refinement

/-!
Probe P6 (code-spine): three facts behind section B's proposals and section C.7.

1. The field `Component.Complete.outputPure` holds of every component: the shared oracle is
   empty, so the prover's output makes no query. (ArkLib has this at the pin as the instance
   `Prover.instOutputIsPureEmpty`, `Composition/Sequential/NoAmbient.lean:39-42`, in a module
   leanerVM does not import and whose `.olean` is not built; it is restated here.)
2. The universes of `Refinement`: whether its two witness types may live in different
   universes (`Column μ : Type`, Clean's `EnsembleWitness : Type 1`). The last `example` of that
   section is the one that decides it.
3. Knowledge transports along a refinement at the same error, with the map applied in the
   event and the game, the provers and the extractor left alone.
-/

open OracleComp OracleSpec ProtocolSpec LeanerVM.Protocol
open scoped NNReal

namespace Probe

/-! ## 1. `outputPure` is automatic -/

/-- The result of a computation that has no possible oracle query (ArkLib
`NoAmbient.lean:25-27` at `dca90385`, restated). -/
def runEmpty {α : Type} : OracleComp []ₒ α → α
  | .pure a => a
  | .liftBind t _ => isEmptyElim t

/-- Every computation over the empty specification is its pure result (`NoAmbient.lean:30-35`,
restated). -/
theorem eq_pure_runEmpty {α : Type} (oa : OracleComp []ₒ α) : oa = pure (runEmpty oa) := by
  cases oa with
  | pure a => rfl
  | queryBind t _ => exact isEmptyElim t

/-- The field `outputPure`, for every component, from its definition alone. -/
theorem outputPure_of_def {StmtIn : Type} {ιi : Type} {OStmtIn : ιi → Type} {WitIn : Type}
    {StmtOut : Type} {ιo : Type} {OStmtOut : ιo → Type} {WitOut : Type}
    [∀ i, OracleInterface (OStmtIn i)] [∀ i, OracleInterface (OStmtOut i)]
    (D : Component.Def StmtIn OStmtIn WitIn StmtOut OStmtOut WitOut) :
    D.red.prover.OutputIsPure :=
  ⟨fun st ↦ runEmpty (D.red.prover.output st),
    fun st ↦ eq_pure_runEmpty (D.red.prover.output st)⟩

#print axioms outputPure_of_def

/-! ## 3. Knowledge transports along a refinement, in the event -/

variable {ι : Type} {oSpec : OracleSpec ι} {StmtIn WitIn StmtOut WitOut : Type} {n : ℕ}
  {pSpec : ProtocolSpec n} [∀ i, SampleableType (pSpec.Challenge i)]
  {σ : Type} (init : ProbComp σ) (impl : QueryImpl oSpec (StateT σ ProbComp))

/-- If `E` extracts a witness of `R` except with probability `ε`, then the same extractor
followed by the map of a refinement from `R` to `S` yields a witness of `S` except with
probability `ε`. -/
theorem knowledge_transport {R : Set (StmtIn × WitIn)} {relOut : Set (StmtOut × WitOut)}
    {W₂ : Type} {S : Set (StmtIn × W₂)} (f : Refinement R S)
    {verifier : Verifier oSpec StmtIn StmtOut pSpec}
    {E : Extractor.Straightline oSpec StmtIn WitIn WitOut pSpec} {ε : ℝ≥0}
    (h : verifier.knowledgeSoundnessWith init impl R relOut E ε)
    (stmtIn : StmtIn) (witIn : WitIn)
    (prover : Prover oSpec StmtIn WitIn StmtOut WitOut pSpec) :
    let pImpl : QueryImpl (oSpec + [pSpec.Challenge]ₒ) (StateT σ ProbComp) :=
      impl.addLift challengeQueryImpl
    let exec := do
      let ⟨⟨⟨transcript, ⟨_, witOut⟩⟩, stmtOut⟩, proveQueryLog, verifyQueryLog⟩
        ← (Reduction.mk prover verifier).runWithLog stmtIn witIn
      let extractedWitIn? ←
        liftM (E stmtIn witOut transcript proveQueryLog.fst verifyQueryLog).run
      return (stmtIn, extractedWitIn?, stmtOut, witOut)
    Pr[fun ⟨stmtIn, extractedWitIn?, stmtOut, witOut⟩ =>
        (∀ w ∈ extractedWitIn?, (stmtIn, f.map stmtIn w) ∉ S) ∧ (stmtOut, witOut) ∈ relOut
      | OptionT.mk do (simulateQ pImpl exec.run).run' (← init)] ≤ ε :=
  le_trans (probEvent_mono fun _ _ hx ↦
    ⟨fun w hw hR ↦ hx.1 w hw (f.map_valid _ _ hR), hx.2⟩) (h stmtIn witIn prover)

#print axioms knowledge_transport

/-! ## 2. The universes of `Refinement` -/

#check @Refinement

/-- A witness type in `Type 1`, as Clean's `EnsembleWitness` is. -/
structure Big : Type 1 where
  carrier : Type

/-- Decides the question: a refinement from a relation whose witness is in `Type` to one whose
witness is in `Type 1`. Accepted if the two witness universes of `Refinement` are independent,
rejected if `{W₁ W₂ : Type _}` gave them one universe. -/
example : Refinement (Stmt := Unit) {p : Unit × ℕ | True} {p : Unit × Big | True} where
  map := fun _ n ↦ ⟨Fin n⟩
  map_valid := fun _ _ _ ↦ trivial

end Probe
\end{leancode}

\fhead{Command}run from the repository root, under the shared lock:

\begin{shellcode}
flock .claude/reports/blueprint-review/logs/lean.lock lake env lean .claude/reports/blueprint-review/probes/code-spine/P6Surface.lean
\end{shellcode}

\fhead{Recorded output}

\begin{srccode}
'Probe.outputPure_of_def' depends on axioms: [propext, Classical.choice, Quot.sound]
'Probe.knowledge_transport' depends on axioms: [propext, Classical.choice, Quot.sound]
@Refinement : {Stmt : Type} → {W₁ : Type u_1} → {W₂ : Type u_2} → Set (Stmt × W₁) → Set (Stmt × W₂) → Type (max u_1 u_2)
(two linter warnings on unused binder names, omitted)
exit=0
\end{srccode}

\subsection{The public-input phase: mutation probes (C)}\label{sec:gen-probes-pubinput}

\par\noindent{\small\textit{Cross-references.} A reference such as \enquote{(D.5)} or \enquote{section B} in this part is to a section of the dossier \file{dossiers/code-pubinput.md}, whose codes the headings keep in parentheses: A (catalogue) is \cref{sec:gen-cat-pubinput}; B (faithfulness), E (the flag \lean{PublicLine.sent}) and F (statement checks) is \cref{sec:gen-conf-pubinput}; C (mutation probes) and D (the top limb) is \cref{sec:gen-probes-pubinput}; G (findings) is \cref{sec:gen-findings-pubinput}.}\par

Every probe of this part is a copy of the phase's module with one mutation (C.0); the dossier reproduces the mutation and the declarations appended to the copy, and the rest of each file is the module at \texttt{b435631}. Each was run with the command of C.0 on its own file, from the repository root, under the shared lock; the probes marked \enquote{written, not run} were not compiled, because the environment was lost (C.9).

\subsubsection{Method (C.0)}\label{sec:gen-pubinput-C0}

\textbf{How a probe is built.} The phase's module has private imports and private lemmas (\texttt{lineClaim\_\allowbreak{}holds\_\allowbreak{}iff}, \texttt{line\_\allowbreak{}challenge\_\allowbreak{}unique}, \texttt{pooled\_\allowbreak{}mem\_\allowbreak{}pub}, \texttt{bad\_\allowbreak{}challenge\_\allowbreak{}unique}, \texttt{simulateQ\_\allowbreak{}queryValues}, \texttt{prover\_\allowbreak{}run\_\allowbreak{}support}), which a plain file cannot reach. So each probe is a \textbf{copy of the whole module, kept a \texttt{module} file}, under the namespace \texttt{LeanerVM.\allowbreak{}Protocol.\allowbreak{}Probe<k>}, with one mutation applied by exact string replacement, and with new declarations appended after the copied text. The script is \texttt{p\allowbreak{}r\allowbreak{}o\allowbreak{}b\allowbreak{}e\allowbreak{}s\allowbreak{}/\allowbreak{}c\allowbreak{}o\allowbreak{}d\allowbreak{}e\allowbreak{}-{}\allowbreak{}p\allowbreak{}u\allowbreak{}b\allowbreak{}i\allowbreak{}n\allowbreak{}p\allowbreak{}u\allowbreak{}t\allowbreak{}/\allowbreak{}t\allowbreak{}o\allowbreak{}o\allowbreak{}l\allowbreak{}s\allowbreak{}/\allowbreak{}m\allowbreak{}u\allowbreak{}t\allowbreak{}a\allowbreak{}t\allowbreak{}e\allowbreak{}.\allowbreak{}p\allowbreak{}y}; the replacements of mutation \texttt{k} are \texttt{p\allowbreak{}r\allowbreak{}o\allowbreak{}b\allowbreak{}e\allowbreak{}s\allowbreak{}/\allowbreak{}c\allowbreak{}o\allowbreak{}d\allowbreak{}e\allowbreak{}-{}\allowbreak{}p\allowbreak{}u\allowbreak{}b\allowbreak{}i\allowbreak{}n\allowbreak{}p\allowbreak{}u\allowbreak{}t\allowbreak{}/\allowbreak{}t\allowbreak{}o\allowbreak{}o\allowbreak{}l\allowbreak{}s\allowbreak{}/\allowbreak{}m\allowbreak{}<\allowbreak{}k\allowbreak{}>\allowbreak{}.\allowbreak{}p\allowbreak{}y}; the shared text (the two refutation lemmas, the pool that reads the message) is \texttt{p\allowbreak{}r\allowbreak{}o\allowbreak{}b\allowbreak{}e\allowbreak{}s\allowbreak{}/\allowbreak{}c\allowbreak{}o\allowbreak{}d\allowbreak{}e\allowbreak{}-{}\allowbreak{}p\allowbreak{}u\allowbreak{}b\allowbreak{}i\allowbreak{}n\allowbreak{}p\allowbreak{}u\allowbreak{}t\allowbreak{}/\allowbreak{}t\allowbreak{}o\allowbreak{}o\allowbreak{}l\allowbreak{}s\allowbreak{}/\allowbreak{}l\allowbreak{}i\allowbreak{}b\allowbreak{}.\allowbreak{}p\allowbreak{}y}. The probe does not import the original \texttt{L\allowbreak{}e\allowbreak{}a\allowbreak{}n\allowbreak{}e\allowbreak{}r\allowbreak{}V\allowbreak{}M\allowbreak{}.\allowbreak{}P\allowbreak{}r\allowbreak{}o\allowbreak{}t\allowbreak{}o\allowbreak{}c\allowbreak{}o\allowbreak{}l\allowbreak{}.\allowbreak{}P\allowbreak{}u\allowbreak{}b\allowbreak{}l\allowbreak{}i\allowbreak{}c\allowbreak{}I\allowbreak{}n\allowbreak{}p\allowbreak{}u\allowbreak{}t}.

\textbf{The baseline.} \texttt{Probe0Baseline.\allowbreak{}lean} is the module with only the namespace changed (the \texttt{diff} against the source is the two namespace lines). It compiles:

\begin{srccode}
$ flock .claude/reports/blueprint-review/logs/lean.lock lake env lean \
    -D autoImplicit=false -D relaxedAutoImplicit=false \
    .claude/reports/blueprint-review/probes/code-pubinput/Probe0Baseline.lean
'LeanerVM.Protocol.Probe0.publicInputSecurity' depends on axioms: [propext, Classical.choice, Quot.sound]
'LeanerVM.Protocol.Probe0.publicInputComplete' depends on axioms: [propext, Classical.choice, Quot.sound]
exit=0
\end{srccode}

Every probe below was run with the same command on its own file, from the repository root, under the shared lock.

\textbf{What \enquote{caught} means.} For each mutation three things are reported.

\begin{itemize}
\item \emph{Unchanged proofs}: the verifier is changed, the verdict lemma \texttt{verifier\_\allowbreak{}verify} is restated for the verifier as changed (it describes what the verifier does), and \texttt{complete}, \texttt{stateFunction}, \texttt{rbr} are left as they are, except where a one-line adaptation is forced by the restated verdict. The first declaration that stops compiling is reported with the compiler's message.
\item \emph{Repair}: whether another state function (or another honest prover) proves the theorem for the mutated verifier at the same error. If it does, the mutation \textbf{is not caught}: a developer who makes the mutation can still produce \texttt{Phase.\allowbreak{}Security}, and with it the master theorems.
\item \emph{Refutation}: where no repair exists, a Lean theorem that the mutated verifier has \textbf{no} round-by-round knowledge soundness at any error below one (whatever the extractor, the intermediate witness types and the state function), or is not perfectly complete. Then the mutation \textbf{is caught}.
\end{itemize}

\textbf{The objects, for a reader who does not know ArkLib} (pin \texttt{dca90385}).

\begin{itemize}
\item A \emph{verifier} (\texttt{Verifier}) maps a statement and a full transcript to an optional output statement. An \emph{oracle verifier} (\texttt{OracleVerifier}) does the same but reads the prover's messages and the input oracles through queries; \texttt{toVerifier} turns it into a verifier by answering the queries.
\item \emph{Perfect completeness} of a reduction (a prover and a verifier) against an input relation and an output relation (\texttt{R\allowbreak{}e\allowbreak{}d\allowbreak{}u\allowbreak{}c\allowbreak{}t\allowbreak{}i\allowbreak{}o\allowbreak{}n\allowbreak{}.\allowbreak{}p\allowbreak{}e\allowbreak{}r\allowbreak{}f\allowbreak{}e\allowbreak{}c\allowbreak{}t\allowbreak{}C\allowbreak{}o\allowbreak{}m\allowbreak{}p\allowbreak{}l\allowbreak{}e\allowbreak{}t\allowbreak{}e\allowbreak{}n\allowbreak{}e\allowbreak{}s\allowbreak{}s}, \texttt{A\allowbreak{}r\allowbreak{}k\allowbreak{}L\allowbreak{}i\allowbreak{}b\allowbreak{}/\allowbreak{}O\allowbreak{}r\allowbreak{}a\allowbreak{}c\allowbreak{}l\allowbreak{}e\allowbreak{}R\allowbreak{}e\allowbreak{}d\allowbreak{}u\allowbreak{}c\allowbreak{}t\allowbreak{}i\allowbreak{}o\allowbreak{}n\allowbreak{}/\allowbreak{}S\allowbreak{}e\allowbreak{}c\allowbreak{}u\allowbreak{}r\allowbreak{}i\allowbreak{}t\allowbreak{}y\allowbreak{}/\allowbreak{}B\allowbreak{}a\allowbreak{}s\allowbreak{}i\allowbreak{}c\allowbreak{}.\allowbreak{}l\allowbreak{}e\allowbreak{}a\allowbreak{}n}): on every input in the input relation, with probability one the verifier accepts, its output is in the output relation with the prover's output witness, \textbf{and the prover's output statement equals the verifier's} (\texttt{p\allowbreak{}e\allowbreak{}r\allowbreak{}f\allowbreak{}e\allowbreak{}c\allowbreak{}t\allowbreak{}C\allowbreak{}o\allowbreak{}m\allowbreak{}p\allowbreak{}l\allowbreak{}e\allowbreak{}t\allowbreak{}e\allowbreak{}n\allowbreak{}e\allowbreak{}s\allowbreak{}s\allowbreak{}\_\allowbreak{}e\allowbreak{}q\allowbreak{}\_\allowbreak{}p\allowbreak{}r\allowbreak{}o\allowbreak{}b\allowbreak{}\_\allowbreak{}o\allowbreak{}n\allowbreak{}e}, \texttt{:164-{}175}).
\item A \emph{round-by-round extractor} (\texttt{Extractor.\allowbreak{}RoundByRound}, \texttt{A\allowbreak{}r\allowbreak{}k\allowbreak{}L\allowbreak{}i\allowbreak{}b\allowbreak{}/\allowbreak{}O\allowbreak{}r\allowbreak{}a\allowbreak{}c\allowbreak{}l\allowbreak{}e\allowbreak{}R\allowbreak{}e\allowbreak{}d\allowbreak{}u\allowbreak{}c\allowbreak{}t\allowbreak{}i\allowbreak{}o\allowbreak{}n\allowbreak{}/\allowbreak{}S\allowbreak{}e\allowbreak{}c\allowbreak{}u\allowbreak{}r\allowbreak{}i\allowbreak{}t\allowbreak{}y\allowbreak{}/\allowbreak{}R\allowbreak{}o\allowbreak{}u\allowbreak{}n\allowbreak{}d\allowbreak{}B\allowbreak{}y\allowbreak{}R\allowbreak{}o\allowbreak{}u\allowbreak{}n\allowbreak{}d\allowbreak{}.\allowbreak{}l\allowbreak{}e\allowbreak{}a\allowbreak{}n\allowbreak{}:\allowbreak{}8\allowbreak{}0\allowbreak{}-{}\allowbreak{}9\allowbreak{}1}) maps witnesses backwards, round by round, from the output witness to the input witness. Here all witnesses are \texttt{Unit}, since the stack is the oracle.
\item A \emph{knowledge state function} (\texttt{V\allowbreak{}e\allowbreak{}r\allowbreak{}i\allowbreak{}f\allowbreak{}i\allowbreak{}e\allowbreak{}r\allowbreak{}.\allowbreak{}K\allowbreak{}n\allowbreak{}o\allowbreak{}w\allowbreak{}l\allowbreak{}e\allowbreak{}d\allowbreak{}g\allowbreak{}e\allowbreak{}S\allowbreak{}t\allowbreak{}a\allowbreak{}t\allowbreak{}e\allowbreak{}F\allowbreak{}u\allowbreak{}n\allowbreak{}c\allowbreak{}t\allowbreak{}i\allowbreak{}o\allowbreak{}n}, same file \texttt{:164-{}186}) is a predicate on (round, statement, partial transcript, intermediate witness) with three laws: at round 0 it is the input relation (\texttt{toFun\_\allowbreak{}empty}); across a \textbf{prover} message, if it holds after the message it held before (\texttt{toFun\_\allowbreak{}next}); if the verifier can output a statement in the output relation, it holds of the full transcript (\texttt{toFun\_\allowbreak{}full}).
\item \emph{Round-by-round knowledge soundness for a given extractor and state function} (\texttt{V\allowbreak{}e\allowbreak{}r\allowbreak{}i\allowbreak{}f\allowbreak{}i\allowbreak{}e\allowbreak{}r\allowbreak{}.\allowbreak{}r\allowbreak{}b\allowbreak{}r\allowbreak{}K\allowbreak{}n\allowbreak{}o\allowbreak{}w\allowbreak{}l\allowbreak{}e\allowbreak{}d\allowbreak{}g\allowbreak{}e\allowbreak{}S\allowbreak{}o\allowbreak{}u\allowbreak{}n\allowbreak{}d\allowbreak{}n\allowbreak{}e\allowbreak{}s\allowbreak{}s\allowbreak{}W\allowbreak{}o\allowbreak{}r\allowbreak{}s\allowbreak{}t\allowbreak{}C\allowbreak{}a\allowbreak{}s\allowbreak{}e\allowbreak{}W\allowbreak{}i\allowbreak{}t\allowbreak{}h}, same file \texttt{:553-{}568}): for every statement, every challenge round and every transcript prefix, the probability over the fresh challenge that the state goes from false to true is at most the error of that challenge. The existential form \texttt{r\allowbreak{}b\allowbreak{}r\allowbreak{}K\allowbreak{}n\allowbreak{}o\allowbreak{}w\allowbreak{}l\allowbreak{}e\allowbreak{}d\allowbreak{}g\allowbreak{}e\allowbreak{}S\allowbreak{}o\allowbreak{}u\allowbreak{}n\allowbreak{}d\allowbreak{}n\allowbreak{}e\allowbreak{}s\allowbreak{}s\allowbreak{}W\allowbreak{}o\allowbreak{}r\allowbreak{}s\allowbreak{}t\allowbreak{}C\allowbreak{}a\allowbreak{}s\allowbreak{}e} (\texttt{:537-{}551}) says that some witness types, some extractor and some state function exist.
\end{itemize}

\begin{leancode}
-- ArkLib dca90385, ArkLib/OracleReduction/Security/RoundByRound.lean:553-568
def rbrKnowledgeSoundnessWorstCaseWith
    (relIn : Set (StmtIn × WitIn)) (relOut : Set (StmtOut × WitOut))
    (verifier : Verifier oSpec StmtIn StmtOut pSpec)
    (WitMid : Fin (n + 1) → Type)
    (extractor : Extractor.RoundByRound oSpec StmtIn WitIn WitOut pSpec WitMid)
    (kSF : verifier.KnowledgeStateFunction init impl relIn relOut extractor)
    (rbrKnowledgeError : pSpec.ChallengeIdx → ℝ≥0) : Prop :=
  ∀ stmtIn : StmtIn,
  ∀ i : pSpec.ChallengeIdx,
  ∀ transcript : Transcript i.1.castSucc pSpec,
    Pr[fun challenge =>
      ∃ witMid,
        ¬ kSF i.1.castSucc stmtIn transcript
          (extractor.extractMid i.1 stmtIn (transcript.concat challenge) witMid) ∧
          kSF i.1.succ stmtIn (transcript.concat challenge) witMid
      | $ᵗ (pSpec.Challenge i)] ≤ rbrKnowledgeError i
\end{leancode}

\textbf{The verifier under test}, as it stands (\texttt{L\allowbreak{}e\allowbreak{}a\allowbreak{}n\allowbreak{}e\allowbreak{}r\allowbreak{}V\allowbreak{}M\allowbreak{}/\allowbreak{}P\allowbreak{}r\allowbreak{}o\allowbreak{}t\allowbreak{}o\allowbreak{}c\allowbreak{}o\allowbreak{}l\allowbreak{}/\allowbreak{}P\allowbreak{}u\allowbreak{}b\allowbreak{}l\allowbreak{}i\allowbreak{}c\allowbreak{}I\allowbreak{}n\allowbreak{}p\allowbreak{}u\allowbreak{}t\allowbreak{}.\allowbreak{}l\allowbreak{}e\allowbreak{}a\allowbreak{}n\allowbreak{}:\allowbreak{}1\allowbreak{}2\allowbreak{}6\allowbreak{}-{}\allowbreak{}1\allowbreak{}3\allowbreak{}7} and \texttt{:229-{}234}):

\begin{leancode}
/-- What the verifier pools: the claims it received, then one claim per public line. -/
def pooled (s : I.Stmt × TableOut I) (r : E) : I.Stmt × PubOut I :=
  (s.1, ⟨s.2.columns ++ (I.publicLines s.1).map (lineClaim I r)⟩)

/-- The values the verifier expects in the prover's message: the values at the challenge of the
lines whose value is sent, in order. -/
def expectedValues (input : I.Stmt) (r : E) : List E :=
  ((I.publicLines input).filter (·.sent)).map (lineValue I r)

/-- The verifier's check: the prover's message is the expected values. -/
def check (s : I.Stmt × TableOut I) (r : E) (cs : List E) : Bool :=
  decide (cs = expectedValues I s.1 r)
\end{leancode}

\begin{leancode}
def verifier : OracleVerifier []ₒ (I.Stmt × TableOut I) (TheOracle I)
    (I.Stmt × PubOut I) (TheOracle I) pSpec where
  verify := fun s chals ↦ do
    let cs ← liftM queryValues
    if check I s (chals ⟨0, rfl⟩) cs then pure (pooled I s (chals ⟨0, rfl⟩)) else failure
  outputOracle := .inl (keepOracles (TheOracle I) pSpec)
\end{leancode}

The observation everything below follows from: \textbf{\texttt{pooled} does not take \texttt{cs}}. The message enters the verifier in \texttt{check} and nowhere else, and the proof of \texttt{rbr} (\texttt{PublicInput.\allowbreak{}lean:350-{}366}) never mentions \texttt{check}.

\subsubsection{The two refutation lemmas (shared by the probes) (C.1)}\label{sec:gen-pubinput-C1}

They are stated for the phase's schedule \texttt{pSpec} (a challenge in \texttt{E}, then a message \texttt{List E}) and proved in each probe that uses them; source in \texttt{tools/\allowbreak{}lib.\allowbreak{}py}, first developed in the plain file \texttt{ScratchRefute.\allowbreak{}lean} against the real \texttt{L\allowbreak{}e\allowbreak{}a\allowbreak{}n\allowbreak{}e\allowbreak{}r\allowbreak{}V\allowbreak{}M\allowbreak{}.\allowbreak{}P\allowbreak{}r\allowbreak{}o\allowbreak{}t\allowbreak{}o\allowbreak{}c\allowbreak{}o\allowbreak{}l\allowbreak{}.\allowbreak{}P\allowbreak{}u\allowbreak{}b\allowbreak{}l\allowbreak{}i\allowbreak{}c\allowbreak{}I\allowbreak{}n\allowbreak{}p\allowbreak{}u\allowbreak{}t} (exit 0).

\begin{leancode}
/-- The transcript `(r, cs)`, built as a knowledge state function reads it. -/
def tr2 (r : E) (cs : List E) : pSpec.FullTranscript :=
  Transcript.concat (m := (1 : Fin 2)) cs
    (Transcript.concat (m := (0 : Fin 2)) r (default : Transcript 0 pSpec))

/-- A guarded verifier on the schedule `pSpec` that, on one statement outside the input
relation, has for every challenge a message it accepts into the output relation is not
round-by-round knowledge sound at an error below one: whatever the extractor and the state
function. -/
theorem not_rbr {StmtIn StmtOut : Type} {relIn : Set (StmtIn × Unit)}
    {relOut : Set (StmtOut × Unit)}
    (V : Verifier []ₒ StmtIn StmtOut pSpec) (G : V.GuardedForm)
    (impl : QueryImpl []ₒ (StateT Unit ProbComp))
    (stmt : StmtIn) (hin : (stmt, ()) ∉ relIn) (msg : E → List E)
    (hacc : ∀ r, G.check stmt (tr2 r (msg r)) = true ∧
      (G.out stmt (tr2 r (msg r)), ()) ∈ relOut)
    (ε : pSpec.ChallengeIdx → ℝ≥0) (hε : ε ⟨0, rfl⟩ < 1) :
    ¬ V.rbrKnowledgeSoundnessWorstCase (pure ()) impl relIn relOut ε := by
  rintro ⟨WitMid, ext, kSF, h⟩
  have hbound := h stmt ⟨0, rfl⟩ (default : Transcript 0 pSpec)
  -- Every challenge is a bad challenge.
  have hall : ∀ r : E, ∃ witMid,
      ¬ kSF (Fin.castSucc 0) stmt (default : Transcript 0 pSpec)
          (ext.extractMid 0 stmt
            (Transcript.concat (m := (0 : Fin 2)) r
              (default : Transcript 0 pSpec)) witMid) ∧
        kSF (Fin.succ 0) stmt (Transcript.concat (m := (0 : Fin 2)) r
          (default : Transcript 0 pSpec)) witMid := by
    intro r
    obtain ⟨hc, hout⟩ := hacc r
    have hpos : Pr[fun stmtOut => (stmtOut, ()) ∈ relOut
        | OptionT.mk do
            (simulateQ impl (V.run stmt (tr2 r (msg r)))).run'
              (← (pure () : ProbComp Unit))] > 0 := by
      have hv : V.run stmt (tr2 r (msg r)) = pure (G.out stmt (tr2 r (msg r))) := by
        have := G.verify_eq stmt (tr2 r (msg r))
        rw [if_pos hc] at this
        exact this
      rw [hv]
      change Pr[_ | OptionT.mk (do let st ← (pure () : ProbComp Unit); (simulateQ impl
        (OptionT.run (pure (G.out stmt (tr2 r (msg r))) :
          OptionT (OracleComp []ₒ) StmtOut))).run' st)] > 0
      rw [OptionT.run_pure, simulateQ_pure]
      rw [gt_iff_lt, probEvent_pos_iff]
      refine ⟨G.out stmt (tr2 r (msg r)), ?_, hout⟩
      simp
    have hfull := kSF.toFun_full stmt (tr2 r (msg r)) () hpos
    have hnext := kSF.toFun_next 1 rfl stmt
      (Transcript.concat (m := (0 : Fin 2)) r (default : Transcript 0 pSpec))
      (msg r) _ hfull
    refine ⟨_, ?_, hnext⟩
    intro h0
    exact hin ((kSF.toFun_empty stmt _).mpr h0)
  -- So the bad event has probability one.
  have hone : Pr[fun challenge : E => ∃ witMid,
      ¬ kSF (Fin.castSucc 0) stmt (default : Transcript 0 pSpec)
          (ext.extractMid 0 stmt
            (Transcript.concat (m := (0 : Fin 2)) challenge
              (default : Transcript 0 pSpec)) witMid) ∧
        kSF (Fin.succ 0) stmt (Transcript.concat (m := (0 : Fin 2)) challenge
          (default : Transcript 0 pSpec)) witMid | $ᵗ E] = 1 := by
    rw [probEvent_eq_one_iff]
    exact ⟨by simp, fun r _ ↦ hall r⟩
  have hle : (1 : ℝ≥0∞) ≤ ((ε ⟨0, rfl⟩ : ℝ≥0) : ℝ≥0∞) := hone ▸ hbound
  exact absurd (ENNReal.coe_lt_one_iff.mpr hε) (not_lt.mpr hle)

/-- A reduction none of whose outcomes, on one input in the input relation, satisfies the
completeness event is not perfectly complete. -/
theorem not_perfectCompleteness {ι : Type} {oSpec : OracleSpec ι}
    {StmtIn WitIn StmtOut WitOut : Type} {n : ℕ} {pSpec : ProtocolSpec n}
    [∀ i, SampleableType (pSpec.Challenge i)]
    {σ : Type} (init : ProbComp σ) (impl : QueryImpl oSpec (StateT σ ProbComp))
    {relIn : Set (StmtIn × WitIn)} {relOut : Set (StmtOut × WitOut)}
    (reduction : Reduction oSpec StmtIn WitIn StmtOut WitOut pSpec)
    (stmtIn : StmtIn) (witIn : WitIn) (hin : (stmtIn, witIn) ∈ relIn)
    (h : ∀ x ∈ support (reduction.run stmtIn witIn).run, ∀ result, x = some result →
      ¬ ((result.2, result.1.2.2) ∈ relOut ∧ result.1.2.1 = result.2)) :
    ¬ reduction.perfectCompleteness init impl relIn relOut := by
  intro hc
  rw [Reduction.perfectCompleteness_eq_prob_one] at hc
  have h1 := hc stmtIn witIn hin
  dsimp only at h1
  have hpos := lt_of_lt_of_eq (zero_lt_one' ℝ≥0∞) h1.symm
  obtain ⟨x, hx, hev⟩ := probEvent_pos_iff.mp hpos
  rw [OptionT.mem_support_iff, OptionT.run_mk, mem_support_bind_iff] at hx
  obtain ⟨s, _, hx⟩ := hx
  exact h (some x) (support_simulateQ_run'_subset _ _ s hx) x rfl hev
\end{leancode}

The first lemma is refutation in the strong sense: it denies ArkLib's \textbf{existential} \texttt{r\allowbreak{}b\allowbreak{}r\allowbreak{}K\allowbreak{}n\allowbreak{}o\allowbreak{}w\allowbreak{}l\allowbreak{}e\allowbreak{}d\allowbreak{}g\allowbreak{}e\allowbreak{}S\allowbreak{}o\allowbreak{}u\allowbreak{}n\allowbreak{}d\allowbreak{}n\allowbreak{}e\allowbreak{}s\allowbreak{}s\allowbreak{}W\allowbreak{}o\allowbreak{}r\allowbreak{}s\allowbreak{}t\allowbreak{}C\allowbreak{}a\allowbreak{}s\allowbreak{}e}, so no choice of extractor, witness types or state function helps, and it holds for every error below one, not only \texttt{1/\allowbreak{}|E|}. It is stated for one choice of the shared-oracle initial state (\texttt{pure (\allowbreak{})}) and implementation; the phase's \texttt{Phase.\allowbreak{}Security} owes its bound for every choice, so one choice refutes it.

\subsubsection{Results so far (C.2)}\label{sec:gen-pubinput-C2}

\begin{landscape}
{\footnotesize
\begin{longtable}{@{}>{\raggedright\arraybackslash}p{0.058\linewidth}>{\raggedright\arraybackslash}p{0.242\linewidth}>{\raggedright\arraybackslash}p{0.098\linewidth}>{\raggedright\arraybackslash}p{0.101\linewidth}>{\raggedright\arraybackslash}p{0.194\linewidth}>{\raggedright\arraybackslash}p{0.172\linewidth}@{}}
\toprule
\# & Mutation & \texttt{complete} & \texttt{s\allowbreak{}t\allowbreak{}a\allowbreak{}t\allowbreak{}e\allowbreak{}F\allowbreak{}u\allowbreak{}n\allowbreak{}c\allowbreak{}t\allowbreak{}i\allowbreak{}o\allowbreak{}n}/\texttt{rbr} unchanged & Knowledge sound by another state function? & Verdict \\
\midrule
\endfirsthead
\toprule
\# & Mutation & \texttt{complete} & \texttt{s\allowbreak{}t\allowbreak{}a\allowbreak{}t\allowbreak{}e\allowbreak{}F\allowbreak{}u\allowbreak{}n\allowbreak{}c\allowbreak{}t\allowbreak{}i\allowbreak{}o\allowbreak{}n}/\texttt{rbr} unchanged & Knowledge sound by another state function? & Verdict \\
\midrule
\endhead
\bottomrule
\endfoot
1 & no check; lines' values pooled & compiles & \texttt{s\allowbreak{}t\allowbreak{}a\allowbreak{}t\allowbreak{}e\allowbreak{}F\allowbreak{}u\allowbreak{}n\allowbreak{}c\allowbreak{}t\allowbreak{}i\allowbreak{}o\allowbreak{}n\allowbreak{}.\allowbreak{}t\allowbreak{}o\allowbreak{}F\allowbreak{}u\allowbreak{}n\allowbreak{}\_\allowbreak{}f\allowbreak{}u\allowbreak{}l\allowbreak{}l} fails & \textbf{yes, same error} (\texttt{rbr\textquotesingle{}}, \texttt{p\allowbreak{}u\allowbreak{}b\allowbreak{}l\allowbreak{}i\allowbreak{}c\allowbreak{}I\allowbreak{}n\allowbreak{}p\allowbreak{}u\allowbreak{}t\allowbreak{}S\allowbreak{}e\allowbreak{}c\allowbreak{}u\allowbreak{}r\allowbreak{}i\allowbreak{}t\allowbreak{}y\allowbreak{}\textquotesingle{}}) & \textbf{not caught} \\
2 & no check; prover's values pooled & compiles & \texttt{toFun\_\allowbreak{}full} fails & no: refuted (\texttt{n\allowbreak{}o\allowbreak{}t\allowbreak{}\_\allowbreak{}k\allowbreak{}n\allowbreak{}o\allowbreak{}w\allowbreak{}l\allowbreak{}e\allowbreak{}d\allowbreak{}g\allowbreak{}e\allowbreak{}S\allowbreak{}o\allowbreak{}u\allowbreak{}n\allowbreak{}d}, toy) & caught by knowledge soundness \\
3a & first value checked; lines' values pooled & compiles & \texttt{toFun\_\allowbreak{}full} fails & \textbf{yes, same error} (\texttt{rbr\textquotesingle{}}) & \textbf{not caught} \\
3b & first value checked; prover's values pooled & compiles & \texttt{toFun\_\allowbreak{}full} fails & no: refuted (two lines sent) & caught by knowledge soundness \\
4a & the pinned verifiers' equation; lines' values pooled & compiles & \texttt{toFun\_\allowbreak{}full} fails & \textbf{yes, same error} (\texttt{rbr\textquotesingle{}}) & \textbf{not caught} \\
4b & the pinned verifiers' equation; prover's values pooled (the deployed verifier) & by C.5 and C.6's arguments &  & sound at \texttt{1/\allowbreak{}|E|} by a new argument; key lemma proved (C.13) & a sound verifier, not a mutation to catch \\
5 & lines not sent are not pooled & compiles & \texttt{toFun\_\allowbreak{}full} fails & no: refuted (a line not sent) & caught by knowledge soundness \\
any & any message, any check the message passes; lines' values pooled & proved (\texttt{completeG}) &  & \textbf{yes, same error} (\texttt{rbrG}, \texttt{securityG}) & \textbf{not caught}, in general \\
6 & claims pooled at a wrong point (paper, C.12) & fails at \texttt{l\allowbreak{}i\allowbreak{}n\allowbreak{}e\allowbreak{}C\allowbreak{}l\allowbreak{}a\allowbreak{}i\allowbreak{}m\allowbreak{}\_\allowbreak{}h\allowbreak{}o\allowbreak{}l\allowbreak{}d\allowbreak{}s\allowbreak{}\_\allowbreak{}i\allowbreak{}f\allowbreak{}f} & fails & no: both refuted on a height-4 instance & caught by completeness (unverified) \\
7a & wrong check, honest prover unchanged (written, not run, C.10) & fails at \texttt{d\allowbreak{}e\allowbreak{}c\allowbreak{}i\allowbreak{}d\allowbreak{}e\allowbreak{}\_\allowbreak{}e\allowbreak{}q\allowbreak{}\_\allowbreak{}t\allowbreak{}r\allowbreak{}u\allowbreak{}e\allowbreak{} h\allowbreak{}m\allowbreak{}s\allowbreak{}g} & compiles & (completeness refuted) & caught by completeness against a fixed prover (unverified) \\
7b & wrong check, prover adapted, lines' values pooled (run, C.4 \texttt{s\allowbreak{}w\allowbreak{}a\allowbreak{}p\allowbreak{}p\allowbreak{}e\allowbreak{}d\allowbreak{}C\allowbreak{}h\allowbreak{}e\allowbreak{}c\allowbreak{}k}) & proved &  & \textbf{yes, same error} & \textbf{not caught} \\
8 & extra check (written, not run, C.11) & fails & compiles & (completeness refuted, for every prover) & caught by completeness (unverified) \\
\end{longtable}}
\end{landscape}

In every probe the unchanged \texttt{stateFunction} fails with the same message, at its field \texttt{toFun\_\allowbreak{}full}, because its last round names the original \texttt{check} and the original \texttt{pooled}:

\begin{srccode}
error: Type mismatch
  Verifier.GuardedForm.of_probEvent_pos (guarded I) init impl stmt tr (fun stmtOut => (stmtOut, x✝) ∈ Seam.pub I) h
has type
  (guarded I).check stmt tr = true ∧ ((guarded I).out stmt tr, x✝) ∈ Seam.pub I
but is expected to have type
  if h0 : ↑(Fin.last 2) = 0 then (stmt, ()) ∈ Seam.table I
  else
    if h1 : ↑(Fin.last 2) = 1 then ((pooled I stmt.1 (tr ⟨0, ⋯⟩), stmt.2), ()) ∈ Seam.pub I
    else check I stmt.1 (tr ⟨0, ⋯⟩) (tr ⟨1, ⋯⟩) = true ∧ ((pooled I stmt.1 (tr ⟨0, ⋯⟩), stmt.2), ()) ∈ Seam.pub I
\end{srccode}

That failure alone says little: the state function is part of the proof, not of the statement (\texttt{Phase.\allowbreak{}Security} carries it as a field, \texttt{kSF}), so whoever changes the verifier changes it too. The decisive columns are the last two.

(Details of each probe: sections C.3 to C.13.)

\subsubsection{Mutation 1: the check removed, the lines' values pooled. NOT CAUGHT (C.3)}\label{sec:gen-pubinput-C3}

File \texttt{p\allowbreak{}r\allowbreak{}o\allowbreak{}b\allowbreak{}e\allowbreak{}s\allowbreak{}/\allowbreak{}c\allowbreak{}o\allowbreak{}d\allowbreak{}e\allowbreak{}-{}\allowbreak{}p\allowbreak{}u\allowbreak{}b\allowbreak{}i\allowbreak{}n\allowbreak{}p\allowbreak{}u\allowbreak{}t\allowbreak{}/\allowbreak{}P\allowbreak{}r\allowbreak{}o\allowbreak{}b\allowbreak{}e\allowbreak{}1\allowbreak{}N\allowbreak{}o\allowbreak{}C\allowbreak{}h\allowbreak{}e\allowbreak{}c\allowbreak{}k\allowbreak{}.\allowbreak{}l\allowbreak{}e\allowbreak{}a\allowbreak{}n} (replacements in \texttt{tools/\allowbreak{}m1.\allowbreak{}py}).

The mutation (the verifier reads the message and ignores it), with the verdict lemma and the guarded form restated for it:

\begin{leancode}
  verify := fun s chals ↦ do
    let _cs ← liftM queryValues
    pure (pooled I s (chals ⟨0, rfl⟩))
\end{leancode}

\begin{leancode}
theorem verifier_verify (s : I.Stmt × TableOut I) (o : ∀ i, TheOracle I i)
    (tr : pSpec.FullTranscript) :
    (verifier I).toVerifier.verify (s, o) tr = pure (pooled I s (tr 0), o) := by
  -- the original proof, ending `rw [simulateQ_optionT_bind_run, simulateQ_queryValues, pure_bind]; rfl`
def guarded : (verifier I).toVerifier.GuardedForm where
  check := fun _ _ ↦ true
  out := fun p tr ↦ (pooled I p.1 (tr 0), p.2)
  verify_eq := fun ⟨s, o⟩ tr ↦ (verifier_verify I s o tr).trans (if_pos rfl).symm
\end{leancode}

\texttt{complete} compiles with one line adapted (\texttt{h\allowbreak{}a\allowbreak{}v\allowbreak{}e\allowbreak{} h\allowbreak{}c\allowbreak{} :\allowbreak{} (\allowbreak{}g\allowbreak{}u\allowbreak{}a\allowbreak{}r\allowbreak{}d\allowbreak{}e\allowbreak{}d\allowbreak{} I\allowbreak{})\allowbreak{}.\allowbreak{}c\allowbreak{}h\allowbreak{}e\allowbreak{}c\allowbreak{}k\allowbreak{} (\allowbreak{}s\allowbreak{},\allowbreak{} o\allowbreak{})\allowbreak{} p\allowbreak{}r\allowbreak{}.\allowbreak{}1\allowbreak{} =\allowbreak{} t\allowbreak{}r\allowbreak{}u\allowbreak{}e\allowbreak{} :\allowbreak{}=\allowbreak{} r\allowbreak{}f\allowbreak{}l}). The unchanged \texttt{stateFunction} fails at \texttt{toFun\_\allowbreak{}full} (message in C.2; line 341 of the probe).

The repair, appended to the probe, is the state function with the check dropped from its last round; \texttt{rbr\textquotesingle{}} has the original proof of \texttt{rbr}, word for word:

\begin{leancode}
def stateFunction' :
    (verifier I).toVerifier.KnowledgeStateFunction init impl (Seam.table I) (Seam.pub I)
      (extractor I) where
  toFun := fun m stmt tr _ ↦
    if h0 : m.val = 0 then ((stmt, ()) ∈ Seam.table I)
    else ((pooled I stmt.1 (tr ⟨0, by omega⟩), stmt.2), ()) ∈ Seam.pub I
  toFun_empty := fun _ _ ↦ Iff.rfl
  toFun_next := fun m hm stmt tr msg w h ↦ by
    have hm1 : m = 1 := by
      fin_cases m
      · exact absurd hm (by decide)
      · rfl
    subst hm1
    exact h
  toFun_full := fun stmt tr _ h ↦
    (Verifier.GuardedForm.of_probEvent_pos (guarded I) init impl stmt tr _ h).2

theorem rbr' :
    (verifier I).toVerifier.rbrKnowledgeSoundnessWorstCaseWith init impl (Seam.table I)
      (Seam.pub I) (fun _ ↦ Unit) (extractor I) (stateFunction' I init impl)
      (fun _ ↦ error) := by
  -- the original proof of `rbr`

def publicInputSecurity' (I : M3Instance) :
    Phase.Security I (publicInputPhase I) (Seam.table I) (Seam.pub I) where
  toComplete := publicInputComplete I
  witMid := fun _ ↦ Unit
  extractor := PublicInput.extractor I
  kSF := PublicInput.stateFunction' I
  rbr := PublicInput.rbr' I
\end{leancode}

Output (exit 1 because of the one expected error):

\begin{srccode}
Probe1NoCheck.lean:341:4: error: Type mismatch   [the message of C.2]
'LeanerVM.Protocol.Probe1.PublicInput.verifier_verify' depends on axioms: [propext, Classical.choice, Quot.sound]
'LeanerVM.Protocol.Probe1.publicInputComplete' depends on axioms: [propext, Classical.choice, Quot.sound]
'LeanerVM.Protocol.Probe1.publicInputSecurity' depends on axioms: [propext, sorryAx, Classical.choice, Quot.sound]
'LeanerVM.Protocol.Probe1.publicInputSecurity'' depends on axioms: [propext, Classical.choice, Quot.sound]
\end{srccode}

\textbf{Reading.} \texttt{publicInputSecurity\textquotesingle{}} is a complete \texttt{Phase.\allowbreak{}Security} of the verifier that checks nothing, against the spine's two seams, at the error \texttt{1/\allowbreak{}|E|}, on the kernel's three axioms. It fills the field \texttt{pub} of \texttt{Phases.\allowbreak{}Security}, so \texttt{piop\_\allowbreak{}rbrKnowledgeSoundness} holds of a protocol whose public-input verifier accepts every message. This is not a defect of the theorem: the mutated verifier \textbf{is} knowledge sound. The prover's message in this phase is a function of public data; a verifier that recomputes it has no need to read it.

\subsubsection{The general form: any message, any check the message passes. NOT CAUGHT (C.4)}\label{sec:gen-pubinput-C4}

File \texttt{p\allowbreak{}r\allowbreak{}o\allowbreak{}b\allowbreak{}e\allowbreak{}s\allowbreak{}/\allowbreak{}c\allowbreak{}o\allowbreak{}d\allowbreak{}e\allowbreak{}-{}\allowbreak{}p\allowbreak{}u\allowbreak{}b\allowbreak{}i\allowbreak{}n\allowbreak{}p\allowbreak{}u\allowbreak{}t\allowbreak{}/\allowbreak{}P\allowbreak{}r\allowbreak{}o\allowbreak{}b\allowbreak{}e\allowbreak{}A\allowbreak{}n\allowbreak{}y\allowbreak{}C\allowbreak{}h\allowbreak{}e\allowbreak{}c\allowbreak{}k\allowbreak{}.\allowbreak{}l\allowbreak{}e\allowbreak{}a\allowbreak{}n} (\texttt{tools/\allowbreak{}anycheck.\allowbreak{}py}): the module unchanged, and appended to it a phase with two parameters, the prover's message \texttt{msg} and the verifier's check \texttt{chk}.

\begin{leancode}
variable (I : M3Instance)
  (msg : I.Stmt → E → List E)
  (chk : (I.Stmt × TableOut I) → E → List E → Bool)

/-- The prover sends `msg`. -/
def proverG : OracleProver []ₒ (I.Stmt × TableOut I) (TheOracle I) Unit
    (I.Stmt × PubOut I) (TheOracle I) Unit pSpec where
  PrvState
    | ⟨0, _⟩ => ((I.Stmt × TableOut I) × ∀ i, TheOracle I i) × Unit
    | _ => E × (((I.Stmt × TableOut I) × ∀ i, TheOracle I i) × Unit)
  input := _root_.id
  receiveChallenge
    | ⟨0, _⟩ => fun st ↦ pure fun r ↦ (r, st)
    | ⟨1, h⟩ => nomatch h
  sendMessage
    | ⟨0, h⟩ => nomatch h
    | ⟨1, _⟩ => fun st ↦ pure (msg st.2.1.1.1 st.1, st)
  output := fun st ↦ pure ((pooled I st.2.1.1 st.1, st.2.1.2), ())

/-- The verifier checks `chk` and pools the lines' values. -/
def verifierG : OracleVerifier []ₒ (I.Stmt × TableOut I) (TheOracle I)
    (I.Stmt × PubOut I) (TheOracle I) pSpec where
  verify := fun s chals ↦ do
    let cs ← liftM queryValues
    if chk s (chals ⟨0, rfl⟩) cs then pure (pooled I s (chals ⟨0, rfl⟩)) else failure
  outputOracle := .inl (keepOracles (TheOracle I) pSpec)
\end{leancode}

\begin{leancode}
/-- Perfect completeness, whatever the message and the check, as long as the one passes the
other. -/
theorem completeG (hchk : ∀ s r, chk s r (msg s.1 r) = true) {σ : Type} (init : ProbComp σ)
    (impl : QueryImpl []ₒ (StateT σ ProbComp)) :
    (OracleReduction.mk (proverG I msg) (verifierG I chk)).perfectCompleteness init impl
      (Seam.table I) (Seam.pub I)

/-- Round-by-round knowledge soundness at `1/|E|`, whatever the check: no hypothesis on
`chk`. -/
theorem rbrG :
    (verifierG I chk).toVerifier.rbrKnowledgeSoundnessWorstCaseWith init impl (Seam.table I)
      (Seam.pub I) (fun _ ↦ Unit) (extractor I) (stateFunctionG I chk init impl)
      (fun _ ↦ error)

/-- Both halves, for every message and every check the message passes. -/
def securityG (I : M3Instance) (msg : I.Stmt → E → List E)
    (chk : (I.Stmt × TableOut I) → E → List E → Bool)
    (hchk : ∀ s r, chk s r (msg s.1 r) = true) :
    Phase.Security I (phaseG I msg chk) (Seam.table I) (Seam.pub I)

/-- No check at all, and a prover that sends nothing. -/
def noCheckNoMessage (I : M3Instance) := securityG I (fun _ _ ↦ []) (fun _ _ _ ↦ true)
  fun _ _ ↦ rfl

/-- No check at all, and a prover that sends forty-two zeros. -/
def noCheckJunk (I : M3Instance) :=
  securityG I (fun _ _ ↦ List.replicate 42 0) (fun _ _ _ ↦ true) fun _ _ ↦ rfl

/-- The values of the lines with their cells swapped. -/
def swappedValues (I : M3Instance) (input : I.Stmt) (r : E) : List E :=
  ((I.publicLines input).filter (·.sent)).map fun l ↦ (1 + r) * ofK l.cell1 + r * ofK l.cell0

/-- A wrong check, the cells swapped, with a prover that sends what the wrong check expects. -/
def swappedCheck (I : M3Instance) :=
  securityG I (swappedValues I) (fun s r cs ↦ decide (cs = swappedValues I s.1 r))
    fun _ _ ↦ decide_eq_true rfl

/-- The specification's phase is the instance with the expected values and their check. -/
def specification (I : M3Instance) :=
  securityG I (PublicInput.expectedValues I) (PublicInput.check I) fun _ _ ↦ decide_eq_true rfl
\end{leancode}

(The proofs are the module's own with \texttt{check} replaced by \texttt{chk} and \texttt{expectedValues} by \texttt{msg}; full text in the probe file.) Output, exit 0:

\begin{srccode}
'LeanerVM.Protocol.ProbeAny.securityG' depends on axioms: [propext, Classical.choice, Quot.sound]
'LeanerVM.Protocol.ProbeAny.noCheckNoMessage' depends on axioms: [propext, Classical.choice, Quot.sound]
'LeanerVM.Protocol.ProbeAny.noCheckJunk' depends on axioms: [propext, Classical.choice, Quot.sound]
'LeanerVM.Protocol.ProbeAny.swappedCheck' depends on axioms: [propext, Classical.choice, Quot.sound]
'LeanerVM.Protocol.ProbeAny.specification' depends on axioms: [propext, Classical.choice, Quot.sound]
\end{srccode}

\textbf{Reading.} The two theorems of the phase constrain the \emph{pool} (which claims, at which point, with which value) and nothing else. They do not constrain the prover's message (its length, its content), nor the check on it. The phase as merged is the instance \texttt{specification} of a family every member of which has the same theorems. In particular:

\begin{itemize}
\item mutation 1 (no check), mutation 3 in the form \enquote{first value only} and mutation 4 in the form \enquote{the pinned verifiers' equation} are members (each also run on its own, C.5);
\item mutation 7 (a wrong check) is a member \textbf{when the honest prover is changed with it} (\texttt{swappedCheck}). Completeness is a statement about the pair (prover, verifier), and nothing fixes the honest prover's message but the verifier's check. So a wrong check on this message is caught by completeness only if the prover is held fixed (C.10).
\end{itemize}

\subsubsection{Mutations 3a and 4a: a weaker check, the lines' values pooled. NOT CAUGHT (C.5)}\label{sec:gen-pubinput-C5}

Files \texttt{Probe3a.\allowbreak{}lean}, \texttt{Probe4a.\allowbreak{}lean} (\texttt{tools/\allowbreak{}m3a.\allowbreak{}py}, \texttt{tools/\allowbreak{}m4a.\allowbreak{}py}, shared \texttt{tools/\allowbreak{}weak.\allowbreak{}py}). The verifier's check is replaced by \texttt{checkWeak}; \texttt{check} stays defined, since the unchanged state function names it.

\begin{leancode}
/-- The weakened check: the first value only. -/                       -- mutation 3a
def checkWeak (s : I.Stmt × TableOut I) (r : E) (cs : List E) : Bool :=
  decide (cs.head? = (expectedValues I s.1 r).head?)
\end{leancode}

\begin{leancode}
/-- The pinned verifiers' check: with two lines whose value is sent, one equation on the two
values, `c₀ + y·c₁ = e₀ + y·e₁`; with any other number, the check per value. -/   -- mutation 4a
def checkWeak (s : I.Stmt × TableOut I) (r : E) (cs : List E) : Bool :=
  match (I.publicLines s.1).filter (·.sent), cs with
  | [l₀, l₁], [c₀, c₁] => decide (c₀ + y * c₁ = lineValue I r l₀ + y * lineValue I r l₁)
  | ls, cs => decide (cs = ls.map (lineValue I r))
\end{leancode}

In both, \texttt{complete} compiles after the honest message is shown to pass the weak check (one line in 3a, a case split in 4a), the unchanged \texttt{stateFunction} fails at \texttt{toFun\_\allowbreak{}full} (line 354, line 365), and the state function with \texttt{checkWeak} in its last round gives \texttt{rbr\textquotesingle{}} by the original proof. Output of both (exit 1, the one expected error each):

\begin{srccode}
Probe3a.lean:354:4: error: Type mismatch   [the message of C.2]
'LeanerVM.Protocol.Probe3a.PublicInput.verifier_verify' depends on axioms: [propext, Classical.choice, Quot.sound]
'LeanerVM.Protocol.Probe3a.PublicInput.complete' depends on axioms: [propext, Classical.choice, Quot.sound]
'LeanerVM.Protocol.Probe3a.PublicInput.stateFunction' depends on axioms: [propext, sorryAx, Classical.choice, Quot.sound]
'LeanerVM.Protocol.Probe3a.PublicInput.rbr' depends on axioms: [propext, sorryAx, Classical.choice, Quot.sound]
'LeanerVM.Protocol.Probe3a.PublicInput.rbr'' depends on axioms: [propext, Classical.choice, Quot.sound]

Probe4a.lean:365:4: error: Type mismatch   [the message of C.2]
'LeanerVM.Protocol.Probe4a.PublicInput.verifier_verify' depends on axioms: [propext, Classical.choice, Quot.sound]
'LeanerVM.Protocol.Probe4a.PublicInput.complete' depends on axioms: [propext, Classical.choice, Quot.sound]
'LeanerVM.Protocol.Probe4a.PublicInput.stateFunction' depends on axioms: [propext, sorryAx, Classical.choice, Quot.sound]
'LeanerVM.Protocol.Probe4a.PublicInput.rbr' depends on axioms: [propext, sorryAx, Classical.choice, Quot.sound]
'LeanerVM.Protocol.Probe4a.PublicInput.rbr'' depends on axioms: [propext, Classical.choice, Quot.sound]
\end{srccode}

\subsubsection{Mutation 2: no check, the prover's values pooled. CAUGHT by knowledge soundness (C.6)}\label{sec:gen-pubinput-C6}

File \texttt{Probe2TrustProver.\allowbreak{}lean} (\texttt{tools/\allowbreak{}m2.\allowbreak{}py}). This is the verifier that trusts the prover. The pool that reads the message (shared by probes 2, 3b and 4b, in \texttt{tools/\allowbreak{}lib.\allowbreak{}py}):

\begin{leancode}
/-- The claims on the lines, a line whose value is sent taking the next value of the message
(zero when the message is too short), the others the value the verifier computes. -/
def claimsFrom (r : E) : List (PublicLine I.toShape) → List E → List (ColumnClaim I)
  | [], _ => []
  | l :: ls, cs =>
    if l.sent then ⟨l.col, linePoint l.pos r, cs.headD 0⟩ :: claimsFrom r ls cs.tail
    else lineClaim I r l :: claimsFrom r ls cs

/-- The pool, with the values the prover sent. -/
def pooledFrom (s : I.Stmt × TableOut I) (r : E) (cs : List E) : I.Stmt × PubOut I :=
  (s.1, ⟨s.2.columns ++ claimsFrom I r (I.publicLines s.1) cs⟩)

/-- On the expected values, the claims are the lines' claims. -/
theorem claimsFrom_expected (r : E) (ls : List (PublicLine I.toShape)) :
    claimsFrom I r ls ((ls.filter (·.sent)).map (lineValue I r)) = ls.map (lineClaim I r)

/-- On the expected values, the pool is the pool of the lines' values. -/
theorem pooledFrom_expected (s : I.Stmt × TableOut I) (r : E) :
    pooledFrom I s r (expectedValues I s.1 r) = pooled I s r
\end{leancode}

The mutation:

\begin{leancode}
  verify := fun s chals ↦ do
    let cs ← liftM queryValues
    pure (pooledFrom I s (chals ⟨0, rfl⟩) cs)
\end{leancode}

\texttt{complete} compiles (the honest message makes the two pools equal, by \texttt{pooledFrom\_\allowbreak{}expected}). The unchanged \texttt{stateFunction} fails at \texttt{toFun\_\allowbreak{}full} (line 376). No repair exists: appended to the probe, on the toy instance, with the stack \texttt{badLine} whose column 2 is \texttt{[1,\allowbreak{} 1]} where the statement \texttt{1} says \texttt{(\allowbreak{}1,\allowbreak{} 0)}:

\begin{leancode}
abbrev badLine : Column 3 := ⟨#v[1, 1, 1, 1, 1, 1, 0, 0]⟩

abbrev badStmt : (K × TableOut toy) × (∀ i, TheOracle toy i) := (((1 : K), ⟨[]⟩), fun _ ↦ badLine)

/-- The wrong stack is outside the table seam. -/
theorem badStmt_not_mem : (badStmt, ()) ∉ Seam.table toy := by
  rintro ⟨-, hl, -⟩
  have h := (hl _ (List.mem_singleton_self _)).2
  exact one_ne_zero h

/-- What a truthful prover sends: the extension of its column 2 at `(r)`. -/
def trueValue (r : E) : E :=
  CMlPolynomialEval.eval₂Mle (toy.column badLine ⟨0, 2⟩).values (algebraMap K E)
    (linePoint (n := 1) (by decide) r)

/-- The mutated verifier accepts it at every challenge, with a pool that holds of the wrong
stack. -/
theorem accepts (r : E) :
    (guarded toy).check badStmt (tr2 r [trueValue r]) = true ∧
      ((guarded toy).out badStmt (tr2 r [trueValue r]), ()) ∈ Seam.pub toy := by
  refine ⟨rfl, ?_, trivial⟩
  intro c hc
  change c ∈ [(⟨⟨0, 2⟩, linePoint (n := 1) (by decide) r, trueValue r⟩ : ColumnClaim toy)] at hc
  rw [List.mem_singleton] at hc
  subst hc
  rfl

/-- No extractor, no state function and no error below one make the mutated verifier
round-by-round knowledge sound against the two seams. -/
theorem not_knowledgeSound (ε : pSpec.ChallengeIdx → ℝ≥0) (hε : ε ⟨0, rfl⟩ < 1) :
    ¬ (verifier toy).toVerifier.rbrKnowledgeSoundnessWorstCase (pure ()) noOracle
      (Seam.table toy) (Seam.pub toy) ε :=
  not_rbr _ (guarded toy) noOracle badStmt badStmt_not_mem (fun r ↦ [trueValue r]) accepts ε hε
\end{leancode}

Output (exit 1, the one expected error):

\begin{srccode}
Probe2TrustProver.lean:376:4: error: Type mismatch   [the message of C.2]
'LeanerVM.Protocol.Probe2.PublicInput.verifier_verify' depends on axioms: [propext, Classical.choice, Quot.sound]
'LeanerVM.Protocol.Probe2.PublicInput.complete' depends on axioms: [propext, Classical.choice, Quot.sound]
'LeanerVM.Protocol.Probe2.PublicInput.stateFunction' depends on axioms: [propext, sorryAx, Classical.choice, Quot.sound]
'LeanerVM.Protocol.Probe2.PublicInput.rbr' depends on axioms: [propext, sorryAx, Classical.choice, Quot.sound]
'LeanerVM.Protocol.Probe2.PublicInput.not_knowledgeSound' depends on axioms: [propext, Classical.choice, Quot.sound]
\end{srccode}

\textbf{Reading.} With the prover's values pooled, the check is what ties the pool to the statement, and its removal is caught: there is no \texttt{Phase.\allowbreak{}Security}. This is the design of the three pinned verifiers and of the build the earlier review read.

\subsubsection{Mutation 3b: first value checked, the prover's values pooled. CAUGHT by knowledge soundness (C.7)}\label{sec:gen-pubinput-C7}

File \texttt{Probe3b.\allowbreak{}lean} (\texttt{tools/\allowbreak{}m3b.\allowbreak{}py}). Verifier: \texttt{i\allowbreak{}f\allowbreak{} c\allowbreak{}h\allowbreak{}e\allowbreak{}c\allowbreak{}k\allowbreak{}W\allowbreak{}e\allowbreak{}a\allowbreak{}k\allowbreak{} …\allowbreak{} t\allowbreak{}h\allowbreak{}e\allowbreak{}n\allowbreak{} p\allowbreak{}u\allowbreak{}r\allowbreak{}e\allowbreak{} (\allowbreak{}p\allowbreak{}o\allowbreak{}o\allowbreak{}l\allowbreak{}e\allowbreak{}d\allowbreak{}F\allowbreak{}r\allowbreak{}o\allowbreak{}m\allowbreak{} …\allowbreak{})\allowbreak{} e\allowbreak{}l\allowbreak{}s\allowbreak{}e\allowbreak{} f\allowbreak{}a\allowbreak{}i\allowbreak{}l\allowbreak{}u\allowbreak{}r\allowbreak{}e}, with \texttt{checkWeak} of mutation 3a. \texttt{complete} compiles; unchanged \texttt{stateFunction} fails at \texttt{toFun\_\allowbreak{}full} (line 390). Refuted on an instance with two lines whose value is sent, by a prover that sends the right first value and the true extension of its wrong second column:

\begin{leancode}
abbrev twoSent : M3Instance :=
  { toy with publicLines := fun v ↦
      [⟨⟨0, 0⟩, v, 1, true, by decide⟩, ⟨⟨0, 1⟩, 1, 1, true, by decide⟩] }

/-- Column 0 is `[1, 1]`, as the statement `1` says; column 1 is `[1, 0]`, not `[1, 1]`. -/
abbrev badSecond : Column 3 := ⟨#v[1, 1, 1, 0, 1, 0, 0, 0]⟩

def cheat (r : E) : List E := [(1 + r) * ofK 1 + r * ofK 1, trueValue 1 r]

theorem accepts (r : E) :
    (guarded twoSent).check badStmt (tr2 r (cheat r)) = true ∧
      ((guarded twoSent).out badStmt (tr2 r (cheat r)), ()) ∈ Seam.pub twoSent

theorem not_knowledgeSound (ε : pSpec.ChallengeIdx → ℝ≥0) (hε : ε ⟨0, rfl⟩ < 1) :
    ¬ (verifier twoSent).toVerifier.rbrKnowledgeSoundnessWorstCase (pure ()) noOracle
      (Seam.table twoSent) (Seam.pub twoSent) ε
\end{leancode}

\begin{srccode}
Probe3b.lean:390:4: error: Type mismatch   [the message of C.2]
'LeanerVM.Protocol.Probe3b.PublicInput.complete' depends on axioms: [propext, Classical.choice, Quot.sound]
'LeanerVM.Protocol.Probe3b.PublicInput.not_knowledgeSound' depends on axioms: [propext, Classical.choice, Quot.sound]
\end{srccode}

\subsubsection{Mutation 5: the claim of a line whose value is not sent is not pooled. CAUGHT by knowledge soundness (C.8)}\label{sec:gen-pubinput-C8}

File \texttt{Probe5.\allowbreak{}lean} (\texttt{tools/\allowbreak{}m5.\allowbreak{}py}). The pool keeps only the lines whose value is sent (for leanISA: the claim on \texttt{mem\_\allowbreak{}2} is dropped); prover and verifier output the same pool; the check is the original.

\begin{leancode}
/-- The mutated pool: the claims received, then one claim per line whose value is sent. -/
def pooledSent (s : I.Stmt × TableOut I) (r : E) : I.Stmt × PubOut I :=
  (s.1, ⟨s.2.columns ++ ((I.publicLines s.1).filter (·.sent)).map (lineClaim I r)⟩)
\end{leancode}

\texttt{complete} compiles (a sub-list of true claims). Unchanged \texttt{stateFunction} fails at \texttt{toFun\_\allowbreak{}full} (line 359). Refuted on the toy with its one line not sent: the honest empty message passes the check, the pool is empty, and the wrong stack is accepted at every challenge.

\begin{leancode}
abbrev noneSent : M3Instance :=
  { toy with publicLines := fun v ↦ [⟨⟨0, 2⟩, v, 0, false, by decide⟩] }

theorem accepts (r : E) :
    (guarded noneSent).check badStmt (tr2 r []) = true ∧
      ((guarded noneSent).out badStmt (tr2 r []), ()) ∈ Seam.pub noneSent

theorem not_knowledgeSound (ε : pSpec.ChallengeIdx → ℝ≥0) (hε : ε ⟨0, rfl⟩ < 1) :
    ¬ (verifier noneSent).toVerifier.rbrKnowledgeSoundnessWorstCase (pure ()) noOracle
      (Seam.table noneSent) (Seam.pub noneSent) ε
\end{leancode}

\begin{srccode}
Probe5.lean:359:4: error: Type mismatch   [the message of C.2]
'LeanerVM.Protocol.Probe5.PublicInput.complete' depends on axioms: [propext, Classical.choice, Quot.sound]
'LeanerVM.Protocol.Probe5.PublicInput.not_knowledgeSound' depends on axioms: [propext, Classical.choice, Quot.sound]
\end{srccode}

\textbf{Reading.} Given an instance that lists the line, dropping its claim is caught. Whether the instance lists it is another matter (section D).

\subsubsection{The probe environment was lost at 09:18 on 2026-09-30 (read this before trusting \enquote{not run}) (C.9)}\label{sec:gen-pubinput-C9}

Between the pause and the restart the review branch received the merge \texttt{8d3ea7d} of \texttt{main} with \texttt{1\allowbreak{}4\allowbreak{}4\allowbreak{}c\allowbreak{}5\allowbreak{}a\allowbreak{}a\allowbreak{} c\allowbreak{}h\allowbreak{}o\allowbreak{}r\allowbreak{}e\allowbreak{}(\allowbreak{}d\allowbreak{}e\allowbreak{}p\allowbreak{}s\allowbreak{})\allowbreak{}:\allowbreak{} u\allowbreak{}p\allowbreak{}g\allowbreak{}r\allowbreak{}a\allowbreak{}d\allowbreak{}e\allowbreak{} t\allowbreak{}o\allowbreak{} L\allowbreak{}e\allowbreak{}a\allowbreak{}n\allowbreak{} 4\allowbreak{}.\allowbreak{}3\allowbreak{}4\allowbreak{}.\allowbreak{}1\allowbreak{} (\allowbreak{}\#\allowbreak{}6\allowbreak{}1\allowbreak{})}, which changes \texttt{lean-{}toolchain} (\texttt{v4.\allowbreak{}33.\allowbreak{}1} to \texttt{v4.\allowbreak{}34.\allowbreak{}1}) and \texttt{lake-{}manifest.\allowbreak{}json} (ArkLib \texttt{dca90385} to \texttt{7653a901}, VCVio \texttt{f9dc47d9} to \texttt{a4232d08}, CompPoly \texttt{3468b38c} to \texttt{572f9973}, Clean \texttt{93c9d1ef} to \texttt{42fe4b26}, Mathlib \texttt{0df444a3} to \texttt{d13f23b7}, PolyFun \texttt{c0c92369} to \texttt{41d3b21d} at a new URL). The first \texttt{lake env lean} after that (mine, probe 7a) made Lake re-resolve the manifest: it checked every package out at its new revision and \textbf{deleted and re-cloned PolyFun}, whose oleans at the old revision are gone (\texttt{find /\allowbreak{} -{}name PolyFun.\allowbreak{}olean} finds nothing). The LeanerVM oleans under \texttt{.\allowbreak{}lake/\allowbreak{}build} were built by 4.33.1 and are refused by 4.34.1 (\enquote{incompatible header}); run with the 4.33.1 binary and a hand-made \texttt{LEAN\_\allowbreak{}PATH} over the surviving package builds, they fail on \texttt{u\allowbreak{}n\allowbreak{}k\allowbreak{}n\allowbreak{}o\allowbreak{}w\allowbreak{}n\allowbreak{} m\allowbreak{}o\allowbreak{}d\allowbreak{}u\allowbreak{}l\allowbreak{}e\allowbreak{} p\allowbreak{}r\allowbreak{}e\allowbreak{}f\allowbreak{}i\allowbreak{}x\allowbreak{} \textquotesingle{}\allowbreak{}P\allowbreak{}o\allowbreak{}l\allowbreak{}y\allowbreak{}F\allowbreak{}u\allowbreak{}n\allowbreak{}\textquotesingle{}}. So from that moment \textbf{no probe can be compiled against the pins without a rebuild}, which the rules forbid. The coordinator's restart message said \enquote{nothing else changed}; the merge contradicts it, and it invalidated \texttt{lake env} as the way to run probes for every agent (the brief's section 8, added afterwards, records the change and asks for no Lean run). I edited no tracked file; the change to \texttt{.\allowbreak{}lake/\allowbreak{}packages} is Lake's own reaction to the merged manifest. (The library sources under \texttt{.\allowbreak{}lake/\allowbreak{}packages/} are now at the new revisions too: every library citation in this dossier was read before 09:18, or is read with \texttt{git show <pin>:<path>} inside the package's repository.)

Probes 1, 2, 3a, 3b, 4a, 5, \texttt{ProbeAnyCheck}, \texttt{ProbeWordsLemma} and \texttt{ScratchRefute} ran before the change, and their outputs above are as printed. Probes 7a and 8 are written (files \texttt{Probe7a.\allowbreak{}lean}, \texttt{Probe8.\allowbreak{}lean}, replacements \texttt{tools/\allowbreak{}m7a.\allowbreak{}py}, \texttt{tools/\allowbreak{}m8.\allowbreak{}py}) and \textbf{not run}; probe 6 and the deployed-verifier phase (4b) are argued on paper below and not written. Each is marked. Where a conclusion below rests on an unrun probe, the conclusion is stated as what the probe is written to show, with the reason it should compile, and is flagged \textbf{unverified}.

\subsubsection{Mutation 7: a wrong check (the cells swapped). Caught by completeness ONLY IF the honest prover is held fixed; probe 7a written, not run (C.10)}\label{sec:gen-pubinput-C10}

\emph{With the prover changed with the check}: the probe \texttt{ProbeAnyCheck} (C.4, run) proves \texttt{s\allowbreak{}w\allowbreak{}a\allowbreak{}p\allowbreak{}p\allowbreak{}e\allowbreak{}d\allowbreak{}C\allowbreak{}h\allowbreak{}e\allowbreak{}c\allowbreak{}k\allowbreak{} I\allowbreak{} :\allowbreak{} P\allowbreak{}h\allowbreak{}a\allowbreak{}s\allowbreak{}e\allowbreak{}.\allowbreak{}S\allowbreak{}e\allowbreak{}c\allowbreak{}u\allowbreak{}r\allowbreak{}i\allowbreak{}t\allowbreak{}y\allowbreak{} I\allowbreak{} (\allowbreak{}p\allowbreak{}h\allowbreak{}a\allowbreak{}s\allowbreak{}e\allowbreak{}G\allowbreak{} I\allowbreak{} (\allowbreak{}s\allowbreak{}w\allowbreak{}a\allowbreak{}p\allowbreak{}p\allowbreak{}e\allowbreak{}d\allowbreak{}V\allowbreak{}a\allowbreak{}l\allowbreak{}u\allowbreak{}e\allowbreak{}s\allowbreak{} I\allowbreak{})\allowbreak{} (\allowbreak{}f\allowbreak{}u\allowbreak{}n\allowbreak{} s\allowbreak{} r\allowbreak{} c\allowbreak{}s\allowbreak{} ↦\allowbreak{} d\allowbreak{}e\allowbreak{}c\allowbreak{}i\allowbreak{}d\allowbreak{}e\allowbreak{} (\allowbreak{}c\allowbreak{}s\allowbreak{} =\allowbreak{} s\allowbreak{}w\allowbreak{}a\allowbreak{}p\allowbreak{}p\allowbreak{}e\allowbreak{}d\allowbreak{}V\allowbreak{}a\allowbreak{}l\allowbreak{}u\allowbreak{}e\allowbreak{}s\allowbreak{} I\allowbreak{} s\allowbreak{}.\allowbreak{}1\allowbreak{} r\allowbreak{})\allowbreak{})\allowbreak{})\allowbreak{} (\allowbreak{}S\allowbreak{}e\allowbreak{}a\allowbreak{}m\allowbreak{}.\allowbreak{}t\allowbreak{}a\allowbreak{}b\allowbreak{}l\allowbreak{}e\allowbreak{} I\allowbreak{})\allowbreak{} (\allowbreak{}S\allowbreak{}e\allowbreak{}a\allowbreak{}m\allowbreak{}.\allowbreak{}p\allowbreak{}u\allowbreak{}b\allowbreak{} I\allowbreak{})}. Both halves hold. \textbf{Not caught.}

\emph{With the honest prover unchanged} (file \texttt{Probe7a.\allowbreak{}lean}, \textbf{not run}): the verifier checks \texttt{checkWrong}, the prover sends \texttt{expectedValues}, the lines' values are pooled.

\begin{leancode}
/-- The values of the lines with their cells swapped: `(1 + r)·cell1 + r·cell0`. -/
def swappedValues (input : I.Stmt) (r : E) : List E :=
  ((I.publicLines input).filter (·.sent)).map fun l ↦ (1 + r) * ofK l.cell1 + r * ofK l.cell0

/-- The wrong check: the message is the swapped values. -/
def checkWrong (s : I.Stmt × TableOut I) (r : E) (cs : List E) : Bool :=
  decide (cs = swappedValues I s.1 r)
\end{leancode}

Expected: \texttt{complete} fails at \texttt{h\allowbreak{}a\allowbreak{}v\allowbreak{}e\allowbreak{} h\allowbreak{}c\allowbreak{} :\allowbreak{} (\allowbreak{}g\allowbreak{}u\allowbreak{}a\allowbreak{}r\allowbreak{}d\allowbreak{}e\allowbreak{}d\allowbreak{} I\allowbreak{})\allowbreak{}.\allowbreak{}c\allowbreak{}h\allowbreak{}e\allowbreak{}c\allowbreak{}k\allowbreak{} (\allowbreak{}s\allowbreak{},\allowbreak{} o\allowbreak{})\allowbreak{} p\allowbreak{}r\allowbreak{}.\allowbreak{}1\allowbreak{} =\allowbreak{} t\allowbreak{}r\allowbreak{}u\allowbreak{}e\allowbreak{} :\allowbreak{}=\allowbreak{} d\allowbreak{}e\allowbreak{}c\allowbreak{}i\allowbreak{}d\allowbreak{}e\allowbreak{}\_\allowbreak{}e\allowbreak{}q\allowbreak{}\_\allowbreak{}t\allowbreak{}r\allowbreak{}u\allowbreak{}e\allowbreak{} h\allowbreak{}m\allowbreak{}s\allowbreak{}g} (\texttt{h\allowbreak{}m\allowbreak{}s\allowbreak{}g\allowbreak{} :\allowbreak{} p\allowbreak{}r\allowbreak{}.\allowbreak{}1\allowbreak{} 1\allowbreak{} =\allowbreak{} e\allowbreak{}x\allowbreak{}p\allowbreak{}e\allowbreak{}c\allowbreak{}t\allowbreak{}e\allowbreak{}d\allowbreak{}V\allowbreak{}a\allowbreak{}l\allowbreak{}u\allowbreak{}e\allowbreak{}s\allowbreak{} I\allowbreak{} s\allowbreak{}.\allowbreak{}1\allowbreak{} (\allowbreak{}p\allowbreak{}r\allowbreak{}.\allowbreak{}1\allowbreak{} 0\allowbreak{})}, but the check asks for \texttt{swappedValues}); \texttt{rbr} compiles (it never reads the check). The refutation written in the probe, on the toy at the statement \texttt{1} with the honest stack (\texttt{honest}, column 2 \texttt{= [1,\allowbreak{} 0]}): the honest message is \texttt{[(\allowbreak{}1 + r)·1 + r·0]}, the swapped value is \texttt{[(\allowbreak{}1 + r)·0 + r·1]}, and \texttt{1 + r = r} gives \texttt{1 = 0}; so the check rejects at every challenge, every outcome of the run is \texttt{none}, and \texttt{not\_\allowbreak{}perfectCompleteness} (C.1, proved) concludes:

\begin{leancode}
theorem honest_rejected (r : E) : checkWrong toy honestStmt.1 r (expectedValues toy 1 r) = false := by
  apply decide_eq_false
  intro h
  change [(1 + r) * ofK 1 + r * ofK 0] = [(1 + r) * ofK 0 + r * ofK 1] at h
  have h1 : ofK (1 : K) = 1 := map_one (algebraMap K E)
  have h0 : ofK (0 : K) = 0 := map_zero (algebraMap K E)
  rw [h0, h1, List.cons.injEq] at h
  have : (1 : E) = 0 := by linear_combination h.1
  exact one_ne_zero this

/-- Not perfectly complete, for every initial state and implementation of the shared oracle. -/
theorem not_complete {σ : Type} (init : ProbComp σ) (impl : QueryImpl []ₒ (StateT σ ProbComp)) :
    ¬ (OracleReduction.mk (prover toy) (verifier toy)).perfectCompleteness init impl
      (Seam.table toy) (Seam.pub toy) := by
  refine not_perfectCompleteness init impl _ honestStmt () honestStmt_mem ?_
  intro x hx result hres
  obtain ⟨pr, hpr, rfl⟩ := Reduction.mem_support_run_of_guarded _ (guarded toy) honestStmt () hx
  obtain ⟨hmsg, -⟩ := prover_run_support toy honestStmt.1 honestStmt.2 pr hpr
  have hc : (guarded toy).check honestStmt pr.1 = false := by
    show checkWrong toy honestStmt.1 (pr.1 0) (pr.1 1) = false
    rw [hmsg]
    exact honest_rejected (pr.1 0)
  rw [hc] at hres
  exact absurd hres (by simp)
\end{leancode}

\textbf{Reading.} \enquote{A wrong check makes completeness unprovable} holds against a fixed honest prover. Nothing in the phase, the spine or the blueprint fixes the honest prover's message except the verifier's own check; the blueprint's acceptance tests contain no transcript-level fixture for this phase (the dumped-proof fixture is Layer 12's). So a wrong check together with a matching prover passes both theorems; only a differential test against the Rust prover's transcript (Layer 12's fixture) would reject it.

\emph{With the prover's values pooled} (the deployed design, not probed): a wrong check is caught whatever the prover, since the check is then what makes the pooled claims true of an honest stack: with the honest prover the check rejects; with the prover adapted to the wrong check the pooled claims carry the swapped values and are false of the honest stack (whenever \texttt{cell0 ≠ cell1}), so completeness fails either way.

\subsubsection{Mutation 8: an extra check leanVM does not make. Caught by completeness; written, not run (C.11)}\label{sec:gen-pubinput-C11}

File \texttt{Probe8.\allowbreak{}lean} (\textbf{not run}). The verifier requires, on top of the check, that the message's first value be nonzero:

\begin{leancode}
/-- The check with an extra condition: the first value of the message is nonzero. -/
def checkExtra (s : I.Stmt × TableOut I) (r : E) (cs : List E) : Bool :=
  check I s r cs && decide (cs.headD 0 ≠ 0)
\end{leancode}

Expected: \texttt{complete} fails at the same line as in 7a (the honest message passes \texttt{check} but \texttt{decide\_\allowbreak{}eq\_\allowbreak{}true hmsg} does not prove \texttt{checkExtra}); \texttt{rbr} compiles. The refutation written in the probe: the toy at the statement \texttt{0}, whose line has cells \texttt{(\allowbreak{}0,\allowbreak{} 0)}, with the honest stack \texttt{h\allowbreak{}o\allowbreak{}n\allowbreak{}e\allowbreak{}s\allowbreak{}t\allowbreak{}0\allowbreak{} :\allowbreak{}=\allowbreak{} ⟨\allowbreak{}\#\allowbreak{}v\allowbreak{}[\allowbreak{}1\allowbreak{},\allowbreak{} 1\allowbreak{},\allowbreak{} 1\allowbreak{},\allowbreak{} 1\allowbreak{},\allowbreak{} 0\allowbreak{},\allowbreak{} 0\allowbreak{},\allowbreak{} 0\allowbreak{},\allowbreak{} 0\allowbreak{}]\allowbreak{}⟩}; the line's value is \texttt{(\allowbreak{}1 + r)·0 + r·0 = 0} at every challenge, so the extra condition fails on the honest message, and (since \texttt{check} pins the message to that value) on every message: no prover passes, honest or adapted. So this mutation is caught \textbf{regardless of the prover}, unlike mutation 7.

\begin{leancode}
theorem honest_rejected (r : E) :
    checkExtra toy honestStmt.1 r (expectedValues toy 0 r) = false := by
  unfold checkExtra
  rw [Bool.and_eq_false_iff]
  right
  apply decide_eq_false
  intro h
  apply h
  change (1 + r) * ofK 0 + r * ofK 0 = 0
  have h0 : ofK (0 : K) = 0 := map_zero (algebraMap K E)
  rw [h0, mul_zero, mul_zero, add_zero]
\end{leancode}

\subsubsection{Mutation 6: the claims pooled at a wrong point. Caught by completeness (and by knowledge soundness); argued on paper, not written (C.12)}\label{sec:gen-pubinput-C12}

Argued on paper; not written. Change \texttt{lineClaim} to pool at \texttt{(\allowbreak{}r,\allowbreak{} 1,\allowbreak{} 0,\allowbreak{} …,\allowbreak{} 0)} (or any point other than \texttt{(\allowbreak{}r,\allowbreak{} 0,\allowbreak{} …,\allowbreak{} 0)}), keeping \texttt{lineValue}.

\emph{Where the proofs stop.} \texttt{lineClaim\_\allowbreak{}holds\_\allowbreak{}iff} (\texttt{PublicInput.\allowbreak{}lean:141-{}148}) rewrites with \texttt{eval₂Mle\_\allowbreak{}linePoint l.\allowbreak{}pos}, whose left side is the extension \textbf{at \texttt{linePoint}}; with another point the rewrite finds no instance and the lemma is unprovable (its statement is false: the extension at \texttt{(\allowbreak{}r,\allowbreak{} 1,\allowbreak{} 0,\allowbreak{} …)} is the line through cells 2 and 3, by the selection identity \texttt{evalMle\_\allowbreak{}append\_\allowbreak{}boolVec} at index 1, \texttt{T\allowbreak{}o\allowbreak{}C\allowbreak{}o\allowbreak{}m\allowbreak{}p\allowbreak{}P\allowbreak{}o\allowbreak{}l\allowbreak{}y\allowbreak{}/\allowbreak{}M\allowbreak{}u\allowbreak{}l\allowbreak{}t\allowbreak{}i\allowbreak{}l\allowbreak{}i\allowbreak{}n\allowbreak{}e\allowbreak{}a\allowbreak{}r\allowbreak{}.\allowbreak{}l\allowbreak{}e\allowbreak{}a\allowbreak{}n\allowbreak{}:\allowbreak{}3\allowbreak{}4\allowbreak{}2\allowbreak{}-{}\allowbreak{}3\allowbreak{}4\allowbreak{}8}). Everything downstream fails: \texttt{pooled\_\allowbreak{}mem\_\allowbreak{}pub} and so \texttt{complete}; \texttt{bad\_\allowbreak{}challenge\_\allowbreak{}unique} and so \texttt{rbr}.

\emph{Completeness is refuted}, on an instance with a column of height 4 (log-height 2, so that a second coordinate exists; the toy's columns have height 2, on which every point \texttt{(\allowbreak{}r)} is the line point, so the toy cannot witness this mutation): one line with cells \texttt{(\allowbreak{}0,\allowbreak{} 0)}, the honest stack \texttt{[0,\allowbreak{} 0,\allowbreak{} 1,\allowbreak{} 1]}. It is in the table seam. The pooled claim has value \texttt{lineValue = 0} and the extension at \texttt{(\allowbreak{}r,\allowbreak{} 1)} is \texttt{(\allowbreak{}1 -{} r)·1 + r·1 = 1 ≠ 0} at every challenge; so the verifier's output is outside \texttt{Seam.\allowbreak{}pub} on every outcome and \texttt{not\_\allowbreak{}perfectCompleteness} applies. \emph{Knowledge soundness is refuted too}: the stack \texttt{[1,\allowbreak{} 1,\allowbreak{} 0,\allowbreak{} 0]} violates the line and its claims at \texttt{(\allowbreak{}r,\allowbreak{} 1)} hold with value \texttt{0} at every challenge (\texttt{not\_\allowbreak{}rbr}). Both refutations need the evaluation of a concrete two-variable table at \texttt{(\allowbreak{}r,\allowbreak{} 1)} with \texttt{r} symbolic, which the toy's \texttt{read\_\allowbreak{}eval} (\texttt{Spine/\allowbreak{}Toy.\allowbreak{}lean:52-{}57}) shows \texttt{simp} can do on tables of this size. \textbf{Unverified} in Lean.

\subsubsection{Mutation 4b: the deployed verifier (the words' equation, the prover's values pooled). Sound at \texttt{1/\allowbreak{}|E|}, by a different argument; the key lemma run, the phase not written (C.13)}\label{sec:gen-pubinput-C13}

The deployed verifier is the mutation \enquote{4a plus 2}: the check \texttt{checkWeak} of C.5 (mutation 4a) and the pool \texttt{pooledFrom} of C.6. Its completeness holds by the argument of C.6 (the honest message makes \texttt{pooledFrom} equal to \texttt{pooled}, and it passes the words' equation, as 4a's \texttt{complete} shows). Its knowledge soundness at \texttt{1/\allowbreak{}|E|} is not a corollary of the phase's: the bad event is now \enquote{some message passes the words' equation and makes all pooled claims true of a stack outside the table seam}, and the message is free. The key lemma is proved (probe \texttt{ProbeWordsLemma.\allowbreak{}lean}, a plain file, \textbf{run before the environment change}, exit 0):

\begin{leancode}
/-- What an accepting run of the pinned verifiers says at the challenge `r`, when the three
pooled claims are true of a stack whose memory limbs have cells `a ℓ` at index 0 and `b ℓ` at
index 1: the scalars are the extensions of limbs 0 and 1, the extension of limb 2 is zero, and
the equation on the words holds. -/
structure Accepts (a b : Fin 3 → K) (w₀ w₁ r c₀ c₁ : E) : Prop where
  claim₀ : (1 - r) * ofK (a 0) + r * ofK (b 0) = c₀
  claim₁ : (1 - r) * ofK (a 1) + r * ofK (b 1) = c₁
  claim₂ : (1 - r) * ofK (a 2) + r * ofK (b 2) = 0
  words : c₀ + y * c₁ = w₀ + r * (w₀ + w₁)

/-- Accepting at two distinct challenges forces the cells: the public words are the two-limb
words of the cells, and the top cells are zero. -/
theorem accepts_two_challenges {a b : Fin 3 → K} {w₀ w₁ r₁ r₂ c₀ c₁ c₀' c₁' : E}
    (hr : r₁ ≠ r₂) (h₁ : Accepts a b w₀ w₁ r₁ c₀ c₁) (h₂ : Accepts a b w₀ w₁ r₂ c₀' c₁') :
    w₀ = E.ofLimbs (a 0) (a 1) 0 ∧ w₁ = E.ofLimbs (b 0) (b 1) 0 ∧ a 2 = 0 ∧ b 2 = 0

/-- For public words with a zero top limb, as leanISA's are by type: the six cells are the four
lanes and two zeros. -/
theorem cells_eq_lanes {a b : Fin 3 → K} {p : Fin 4 → K} {r₁ r₂ c₀ c₁ c₀' c₁' : E}
    (hr : r₁ ≠ r₂)
    (h₁ : Accepts a b (E.ofLimbs (p 0) (p 1) 0) (E.ofLimbs (p 2) (p 3) 0) r₁ c₀ c₁)
    (h₂ : Accepts a b (E.ofLimbs (p 0) (p 1) 0) (E.ofLimbs (p 2) (p 3) 0) r₂ c₀' c₁') :
    (a 0 = p 0 ∧ a 1 = p 1 ∧ a 2 = 0) ∧ (b 0 = p 2 ∧ b 1 = p 3 ∧ b 2 = 0)

/-- A public word with a nonzero top limb is accepted at one challenge at most, whatever the
stack: the binding itself excludes it, without the separate rejection of `read_public`. -/
theorem top_limb_zero_of_two_challenges {a b : Fin 3 → K} {w₀ w₁ r₁ r₂ c₀ c₁ c₀' c₁' : E}
    (hr : r₁ ≠ r₂) (h₁ : Accepts a b w₀ w₁ r₁ c₀ c₁) (h₂ : Accepts a b w₀ w₁ r₂ c₀' c₁') :
    w₀.limb 2 = 0 ∧ w₁.limb 2 = 0
\end{leancode}

\begin{srccode}
'WordsLemma.accepts_two_challenges' depends on axioms: [propext, Classical.choice, Quot.sound]
'WordsLemma.cells_eq_lanes' depends on axioms: [propext, Classical.choice, Quot.sound]
'WordsLemma.top_limb_zero_of_two_challenges' depends on axioms: [propext, Classical.choice, Quot.sound]
\end{srccode}

The proof is eleven lines: substituting the claims into the words' equation gives \texttt{(\allowbreak{}1\allowbreak{} -{}\allowbreak{} r\allowbreak{})\allowbreak{}·\allowbreak{}A\allowbreak{} +\allowbreak{} r\allowbreak{}·\allowbreak{}B\allowbreak{} =\allowbreak{} (\allowbreak{}1\allowbreak{} +\allowbreak{} r\allowbreak{})\allowbreak{}·\allowbreak{}w\allowbreak{}₀\allowbreak{} +\allowbreak{} r\allowbreak{}·\allowbreak{}w\allowbreak{}₁} with \texttt{A = a₀ + y·a₁}, \texttt{B = b₀ + y·b₁} in \texttt{E}; two roots force \texttt{A = w₀}, \texttt{B = w₁} (the phase's own degree-one argument, \texttt{line\_\allowbreak{}unique}, over \texttt{E} instead of \texttt{K}); \texttt{A = ofLimbs a₀ a₁ 0} by \texttt{ofLimbs\_\allowbreak{}eq}; the third claim is a line through zero. No independence of \texttt{1,\allowbreak{} y} over \texttt{K} is used until the limbs are read off (\texttt{cells\_\allowbreak{}eq\_\allowbreak{}lanes}).

\emph{The phase, if built.} Over an abstract instance the words' equation needs a shape: it combines the sent values with the powers \texttt{1,\allowbreak{} y,\allowbreak{} y²,\allowbreak{} …}, and beyond three sent lines those are dependent over \texttt{K}, so the argument fails. The check of C.5 (mutation 4a) handles this by falling back to the per-value check when the number of sent lines is not two; on that verifier, with \texttt{pooledFrom}, the knowledge state function is

\begin{itemize}
\item round 1: \texttt{∃\allowbreak{} c\allowbreak{}s\allowbreak{},\allowbreak{} c\allowbreak{}h\allowbreak{}e\allowbreak{}c\allowbreak{}k\allowbreak{}W\allowbreak{}e\allowbreak{}a\allowbreak{}k\allowbreak{} s\allowbreak{} r\allowbreak{} c\allowbreak{}s\allowbreak{} =\allowbreak{} t\allowbreak{}r\allowbreak{}u\allowbreak{}e\allowbreak{} ∧\allowbreak{} (\allowbreak{}(\allowbreak{}p\allowbreak{}o\allowbreak{}o\allowbreak{}l\allowbreak{}e\allowbreak{}d\allowbreak{}F\allowbreak{}r\allowbreak{}o\allowbreak{}m\allowbreak{} s\allowbreak{} r\allowbreak{} c\allowbreak{}s\allowbreak{},\allowbreak{} o\allowbreak{})\allowbreak{},\allowbreak{} (\allowbreak{})\allowbreak{})\allowbreak{} ∈\allowbreak{} S\allowbreak{}e\allowbreak{}a\allowbreak{}m\allowbreak{}.\allowbreak{}p\allowbreak{}u\allowbreak{}b\allowbreak{} I},
\item round 2: the same for the message sent,
\end{itemize}

and the bad-challenge uniqueness splits into: two sent lines \texttt{l₀,\allowbreak{} l₁} (then the pooled claims force \texttt{cs} to be the true extensions, and \texttt{accepts\_\allowbreak{}two\_\allowbreak{}challenges} with the third claim replaced by the unsent lines' claims), or not (then \texttt{cs = expectedValues}, \texttt{pooledFrom = pooled}, and the existing \texttt{bad\_\allowbreak{}challenge\_\allowbreak{}unique}). The \texttt{List} plumbing (that \texttt{claimsFrom} holding of the stack forces \texttt{cs} to be the true extensions of the sent lines and the unsent lines' claims) is an induction on the lines. Estimated at 120 lines; \textbf{not written}, since the environment was lost. Error: \texttt{1/\allowbreak{}|E|} on the one challenge, as for the specification's verifier.

\subsection{The public-input phase: the top limb, along the chain (D)}\label{sec:gen-probes-pubinput-toplimb}

The user's example: \enquote{If the verifier does not check, through some reduction or claim, that the top lane of both public inputs is 0 (the column \texttt{mem\_\allowbreak{}2} at cells 0 and 1), will the knowledge-soundness and completeness theorems catch this?}

\subsubsection{In the phase and the spine: no (D.1)}\label{sec:gen-pubinput-D1}

\textbf{The third claim exists because the instance lists a third line.} The phase pools one claim per element of \texttt{I.\allowbreak{}publicLines s.\allowbreak{}1} (\texttt{pooled}, C.0); \texttt{p\allowbreak{}u\allowbreak{}b\allowbreak{}l\allowbreak{}i\allowbreak{}c\allowbreak{}L\allowbreak{}i\allowbreak{}n\allowbreak{}e\allowbreak{}s\allowbreak{} :\allowbreak{} S\allowbreak{}t\allowbreak{}m\allowbreak{}t\allowbreak{} →\allowbreak{} L\allowbreak{}i\allowbreak{}s\allowbreak{}t\allowbreak{} (\allowbreak{}P\allowbreak{}u\allowbreak{}b\allowbreak{}l\allowbreak{}i\allowbreak{}c\allowbreak{}L\allowbreak{}i\allowbreak{}n\allowbreak{}e\allowbreak{} t\allowbreak{}o\allowbreak{}S\allowbreak{}h\allowbreak{}a\allowbreak{}p\allowbreak{}e\allowbreak{})} is a field of \texttt{M3Instance} (\texttt{Spine/\allowbreak{}Instance.\allowbreak{}lean:141-{}142} at \texttt{b435631}):

\begin{leancode}
  /-- The columns whose first two cells the statement fixes. -/
  publicLines : Stmt → List (PublicLine toShape)
\end{leancode}

and the relation's clause is over the same list (\texttt{Spine/\allowbreak{}Instance.\allowbreak{}lean:207-{}211}):

\begin{leancode}
/-- Cells 0 and 1 of every column the statement fixes hold their values (§8.2). -/
def PublicLinesHold (input : I.Stmt) (q : Column I.μ) : Prop :=
  ∀ l ∈ I.publicLines input,
    (I.column q l.col).values.get ⟨0, Nat.two_pow_pos _⟩ = l.cell0 ∧
      (I.column q l.col).values.get ⟨1, Nat.one_lt_two_pow l.pos.ne'⟩ = l.cell1
\end{leancode}

Every theorem of the phase and of the spine is stated for an arbitrary \texttt{I : M3Instance}: \texttt{p\allowbreak{}u\allowbreak{}b\allowbreak{}l\allowbreak{}i\allowbreak{}c\allowbreak{}I\allowbreak{}n\allowbreak{}p\allowbreak{}u\allowbreak{}t\allowbreak{}S\allowbreak{}e\allowbreak{}c\allowbreak{}u\allowbreak{}r\allowbreak{}i\allowbreak{}t\allowbreak{}y\allowbreak{} :\allowbreak{} (\allowbreak{}I\allowbreak{} :\allowbreak{} M\allowbreak{}3\allowbreak{}I\allowbreak{}n\allowbreak{}s\allowbreak{}t\allowbreak{}a\allowbreak{}n\allowbreak{}c\allowbreak{}e\allowbreak{})\allowbreak{} →\allowbreak{} P\allowbreak{}h\allowbreak{}a\allowbreak{}s\allowbreak{}e\allowbreak{}.\allowbreak{}S\allowbreak{}e\allowbreak{}c\allowbreak{}u\allowbreak{}r\allowbreak{}i\allowbreak{}t\allowbreak{}y\allowbreak{} I\allowbreak{} (\allowbreak{}p\allowbreak{}u\allowbreak{}b\allowbreak{}l\allowbreak{}i\allowbreak{}c\allowbreak{}I\allowbreak{}n\allowbreak{}p\allowbreak{}u\allowbreak{}t\allowbreak{}P\allowbreak{}h\allowbreak{}a\allowbreak{}s\allowbreak{}e\allowbreak{} I\allowbreak{})\allowbreak{} (\allowbreak{}S\allowbreak{}e\allowbreak{}a\allowbreak{}m\allowbreak{}.\allowbreak{}t\allowbreak{}a\allowbreak{}b\allowbreak{}l\allowbreak{}e\allowbreak{} I\allowbreak{})\allowbreak{} (\allowbreak{}S\allowbreak{}e\allowbreak{}a\allowbreak{}m\allowbreak{}.\allowbreak{}p\allowbreak{}u\allowbreak{}b\allowbreak{} I\allowbreak{})} (\texttt{PublicInput.\allowbreak{}lean:391-{}397}), \texttt{p\allowbreak{}i\allowbreak{}o\allowbreak{}p\allowbreak{}\_\allowbreak{}p\allowbreak{}e\allowbreak{}r\allowbreak{}f\allowbreak{}e\allowbreak{}c\allowbreak{}t\allowbreak{}C\allowbreak{}o\allowbreak{}m\allowbreak{}p\allowbreak{}l\allowbreak{}e\allowbreak{}t\allowbreak{}e\allowbreak{}n\allowbreak{}e\allowbreak{}s\allowbreak{}s\allowbreak{} (\allowbreak{}P\allowbreak{} :\allowbreak{} P\allowbreak{}h\allowbreak{}a\allowbreak{}s\allowbreak{}e\allowbreak{}s\allowbreak{} I\allowbreak{})\allowbreak{} (\allowbreak{}C\allowbreak{} :\allowbreak{} P\allowbreak{}.\allowbreak{}C\allowbreak{}o\allowbreak{}m\allowbreak{}p\allowbreak{}l\allowbreak{}e\allowbreak{}t\allowbreak{}e\allowbreak{})\allowbreak{} …}, \texttt{p\allowbreak{}i\allowbreak{}o\allowbreak{}p\allowbreak{}\_\allowbreak{}r\allowbreak{}b\allowbreak{}r\allowbreak{}K\allowbreak{}n\allowbreak{}o\allowbreak{}w\allowbreak{}l\allowbreak{}e\allowbreak{}d\allowbreak{}g\allowbreak{}e\allowbreak{}S\allowbreak{}o\allowbreak{}u\allowbreak{}n\allowbreak{}d\allowbreak{}n\allowbreak{}e\allowbreak{}s\allowbreak{}s\allowbreak{} (\allowbreak{}P\allowbreak{} :\allowbreak{} P\allowbreak{}h\allowbreak{}a\allowbreak{}s\allowbreak{}e\allowbreak{}s\allowbreak{} I\allowbreak{})\allowbreak{} (\allowbreak{}S\allowbreak{} :\allowbreak{} P\allowbreak{}.\allowbreak{}S\allowbreak{}e\allowbreak{}c\allowbreak{}u\allowbreak{}r\allowbreak{}i\allowbreak{}t\allowbreak{}y\allowbreak{})\allowbreak{} …} (\texttt{Spine/\allowbreak{}Compose.\allowbreak{}lean:162-{}177}). An instance whose \texttt{publicLines} lists \texttt{mem\_\allowbreak{}0} and \texttt{mem\_\allowbreak{}1} only is an \texttt{M3Instance} like any other; on it \texttt{M3Holds} says nothing of \texttt{mem\_\allowbreak{}2}, the phase pools two claims, and every theorem instantiates by \texttt{publicInputSecurity I\textquotesingle{}}, \texttt{p\allowbreak{}i\allowbreak{}o\allowbreak{}p\allowbreak{}\_\allowbreak{}r\allowbreak{}b\allowbreak{}r\allowbreak{}K\allowbreak{}n\allowbreak{}o\allowbreak{}w\allowbreak{}l\allowbreak{}e\allowbreak{}d\allowbreak{}g\allowbreak{}e\allowbreak{}S\allowbreak{}o\allowbreak{}u\allowbreak{}n\allowbreak{}d\allowbreak{}n\allowbreak{}e\allowbreak{}s\allowbreak{}s\allowbreak{} P\allowbreak{} S}. No probe is needed to see this (the quantification is in the types quoted); the probe that was planned (an instance \texttt{twoLines} with \texttt{e\allowbreak{}x\allowbreak{}a\allowbreak{}m\allowbreak{}p\allowbreak{}l\allowbreak{}e\allowbreak{} :\allowbreak{}=\allowbreak{} p\allowbreak{}u\allowbreak{}b\allowbreak{}l\allowbreak{}i\allowbreak{}c\allowbreak{}I\allowbreak{}n\allowbreak{}p\allowbreak{}u\allowbreak{}t\allowbreak{}S\allowbreak{}e\allowbreak{}c\allowbreak{}u\allowbreak{}r\allowbreak{}i\allowbreak{}t\allowbreak{}y\allowbreak{} t\allowbreak{}w\allowbreak{}o\allowbreak{}L\allowbreak{}i\allowbreak{}n\allowbreak{}e\allowbreak{}s} and a \texttt{\#guard} that a stack with a nonzero \texttt{mem\_\allowbreak{}2} cell 0 satisfies \texttt{M3Holds twoLines}) was not run (C.9) and would only have displayed it. Mutation 5 (C.8) is the same fact from the other side: given a listed line, dropping its claim is caught; whether it is listed is not the phase's business.

\textbf{So the oracle protocol's theorems do not catch a missing third line.} What does is the adaptor, not built (blueprint Layer 3, \texttt{p\allowbreak{}r\allowbreak{}o\allowbreak{}t\allowbreak{}o\allowbreak{}c\allowbreak{}o\allowbreak{}l\allowbreak{}-{}\allowbreak{}b\allowbreak{}l\allowbreak{}u\allowbreak{}e\allowbreak{}p\allowbreak{}r\allowbreak{}i\allowbreak{}n\allowbreak{}t\allowbreak{}.\allowbreak{}m\allowbreak{}d\allowbreak{}:\allowbreak{}8\allowbreak{}1\allowbreak{}3\allowbreak{}-{}\allowbreak{}8\allowbreak{}1\allowbreak{}6} at \texttt{b435631}):

\begin{leancode}
def witnessOf (prog) (s) (q : Column (leanIsaInstance prog s).μ) :
    EnsembleWitness (leanIsaEnsemble prog)                            -- total and computable
theorem satisfiedBy_witnessOf (hs : s.Admissible prog)                 -- the caps: a hypothesis
    (h : M3Holds (leanIsaInstance prog s) input q) : SatisfiedBy prog input (witnessOf prog s q)
\end{leancode}

and behind it leanISA's \texttt{constraintSoundness} (leanISA blueprint \texttt{:1163-{}1164}, not built), \texttt{S\allowbreak{}a\allowbreak{}t\allowbreak{}i\allowbreak{}s\allowbreak{}f\allowbreak{}i\allowbreak{}e\allowbreak{}d\allowbreak{}B\allowbreak{}y\allowbreak{} →\allowbreak{} ∃\allowbreak{} t\allowbreak{},\allowbreak{} A\allowbreak{}s\allowbreak{}s\allowbreak{}i\allowbreak{}g\allowbreak{}n\allowbreak{}m\allowbreak{}e\allowbreak{}n\allowbreak{}t\allowbreak{}R\allowbreak{}e\allowbreak{}p\allowbreak{}r\allowbreak{}e\allowbreak{}s\allowbreak{}e\allowbreak{}n\allowbreak{}t\allowbreak{}s\allowbreak{} w\allowbreak{} t\allowbreak{} ∧\allowbreak{} V\allowbreak{}a\allowbreak{}l\allowbreak{}i\allowbreak{}d\allowbreak{}E\allowbreak{}x\allowbreak{}e\allowbreak{}c\allowbreak{}u\allowbreak{}t\allowbreak{}i\allowbreak{}o\allowbreak{}n\allowbreak{} p\allowbreak{}r\allowbreak{}o\allowbreak{}g\allowbreak{} i\allowbreak{}n\allowbreak{}p\allowbreak{}u\allowbreak{}t\allowbreak{} t}, and T4's \texttt{b\allowbreak{}a\allowbreak{}s\allowbreak{}e\allowbreak{}V\allowbreak{}e\allowbreak{}r\allowbreak{}i\allowbreak{}f\allowbreak{}i\allowbreak{}e\allowbreak{}r\allowbreak{}\_\allowbreak{}e\allowbreak{}x\allowbreak{}t\allowbreak{}r\allowbreak{}a\allowbreak{}c\allowbreak{}t\allowbreak{}s\allowbreak{}E\allowbreak{}x\allowbreak{}e\allowbreak{}c\allowbreak{}u\allowbreak{}t\allowbreak{}i\allowbreak{}o\allowbreak{}n} (\texttt{p\allowbreak{}r\allowbreak{}o\allowbreak{}t\allowbreak{}o\allowbreak{}c\allowbreak{}o\allowbreak{}l\allowbreak{}-{}\allowbreak{}b\allowbreak{}l\allowbreak{}u\allowbreak{}e\allowbreak{}p\allowbreak{}r\allowbreak{}i\allowbreak{}n\allowbreak{}t\allowbreak{}.\allowbreak{}m\allowbreak{}d\allowbreak{}:\allowbreak{}1\allowbreak{}2\allowbreak{}4\allowbreak{}7\allowbreak{}-{}\allowbreak{}1\allowbreak{}2\allowbreak{}4\allowbreak{}8}). The links, read in the Lean sources at \texttt{b435631}:

\begin{enumerate}
\item \emph{The type of the public words.} \texttt{PublicInput} is four lanes of \texttt{K}, and its words are built with a zero top limb (\texttt{L\allowbreak{}e\allowbreak{}a\allowbreak{}n\allowbreak{}e\allowbreak{}r\allowbreak{}V\allowbreak{}M\allowbreak{}/\allowbreak{}S\allowbreak{}e\allowbreak{}m\allowbreak{}a\allowbreak{}n\allowbreak{}t\allowbreak{}i\allowbreak{}c\allowbreak{}s\allowbreak{}/\allowbreak{}M\allowbreak{}e\allowbreak{}m\allowbreak{}o\allowbreak{}r\allowbreak{}y\allowbreak{}.\allowbreak{}l\allowbreak{}e\allowbreak{}a\allowbreak{}n\allowbreak{}:\allowbreak{}1\allowbreak{}0\allowbreak{}0\allowbreak{}-{}\allowbreak{}1\allowbreak{}1\allowbreak{}0}):

\begin{leancode}
/-- The 256-bit public input as four `K` lanes (specification §2). -/
structure PublicInput where
  /-- `input₀, …, input₃`. -/
  lanes : Fin 4 → K
  deriving DecidableEq

/-- The first memory word, `input₀ + input₁·y`. -/
def PublicInput.word0 (p : PublicInput) : E := E.ofLimbs (p.lanes 0) (p.lanes 1) 0

/-- The second memory word, `input₂ + input₃·y`. -/
def PublicInput.word1 (p : PublicInput) : E := E.ofLimbs (p.lanes 2) (p.lanes 3) 0
\end{leancode}

with \texttt{E\allowbreak{}.\allowbreak{}o\allowbreak{}f\allowbreak{}L\allowbreak{}i\allowbreak{}m\allowbreak{}b\allowbreak{}s\allowbreak{} c\allowbreak{}0\allowbreak{} c\allowbreak{}1\allowbreak{} c\allowbreak{}2\allowbreak{} =\allowbreak{} c\allowbreak{}0\allowbreak{} +\allowbreak{} c\allowbreak{}1\allowbreak{}·\allowbreak{}y\allowbreak{} +\allowbreak{} c\allowbreak{}2\allowbreak{}·\allowbreak{}y\allowbreak{}²} (\texttt{Parameters/\allowbreak{}Field.\allowbreak{}lean:92}, \texttt{:145-{}146}). \textbf{A public word with a nonzero top limb is not representable} as a statement. This matches the Rust, whose \texttt{read\_\allowbreak{}public} rejects it before the protocol starts (\texttt{cpu/\allowbreak{}mod.\allowbreak{}rs:139-{}143}: \enquote{The transcript binds a public input as two 128-bit halves, so a third limb would be dropped and two statements would share a transcript. \texttt{i\allowbreak{}f\allowbreak{} p\allowbreak{}u\allowbreak{}b\allowbreak{}l\allowbreak{}i\allowbreak{}c\allowbreak{}\_\allowbreak{}i\allowbreak{}n\allowbreak{}p\allowbreak{}u\allowbreak{}t\allowbreak{}.\allowbreak{}i\allowbreak{}t\allowbreak{}e\allowbreak{}r\allowbreak{}(\allowbreak{})\allowbreak{}.\allowbreak{}a\allowbreak{}n\allowbreak{}y\allowbreak{}(\allowbreak{}|\allowbreak{}h\allowbreak{}a\allowbreak{}l\allowbreak{}f\allowbreak{}|\allowbreak{} h\allowbreak{}a\allowbreak{}l\allowbreak{}f\allowbreak{}.\allowbreak{}c\allowbreak{}2\allowbreak{} !\allowbreak{}=\allowbreak{} 0\allowbreak{})\allowbreak{} \{\allowbreak{} r\allowbreak{}e\allowbreak{}t\allowbreak{}u\allowbreak{}r\allowbreak{}n\allowbreak{} E\allowbreak{}r\allowbreak{}r\allowbreak{}(\allowbreak{}C\allowbreak{}p\allowbreak{}u\allowbreak{}E\allowbreak{}r\allowbreak{}r\allowbreak{}o\allowbreak{}r\allowbreak{}:\allowbreak{}:\allowbreak{}P\allowbreak{}u\allowbreak{}b\allowbreak{}l\allowbreak{}i\allowbreak{}c\allowbreak{}I\allowbreak{}n\allowbreak{}p\allowbreak{}u\allowbreak{}t\allowbreak{})\allowbreak{};\allowbreak{} \}}}), and the Python, whose \texttt{Digest.\allowbreak{}halves} builds \texttt{E(\allowbreak{}w0,\allowbreak{} w1),\allowbreak{} E(\allowbreak{}w2,\allowbreak{} w3)} with no third limb (\texttt{verifier.\allowbreak{}py:216-{}219}).
\item \emph{What \texttt{SatisfiedBy} says of cells 0 and 1.} Two conjuncts, on the image the prover data names (\texttt{L\allowbreak{}e\allowbreak{}a\allowbreak{}n\allowbreak{}e\allowbreak{}r\allowbreak{}V\allowbreak{}M\allowbreak{}/\allowbreak{}A\allowbreak{}r\allowbreak{}i\allowbreak{}t\allowbreak{}h\allowbreak{}m\allowbreak{}e\allowbreak{}t\allowbreak{}i\allowbreak{}z\allowbreak{}a\allowbreak{}t\allowbreak{}i\allowbreak{}o\allowbreak{}n\allowbreak{}/\allowbreak{}S\allowbreak{}t\allowbreak{}a\allowbreak{}t\allowbreak{}e\allowbreak{}m\allowbreak{}e\allowbreak{}n\allowbreak{}t\allowbreak{}.\allowbreak{}l\allowbreak{}e\allowbreak{}a\allowbreak{}n\allowbreak{}:\allowbreak{}3\allowbreak{}3\allowbreak{}7\allowbreak{}-{}\allowbreak{}3\allowbreak{}4\allowbreak{}0}):

\begin{leancode}
  /-- The first word of the image is `input₀ + input₁·y` (§8.2). -/
  word0_eq : (imageOf w.data).2.read (gpow 0) = some input.word0
  /-- The second word of the image is `input₂ + input₃·y` (§8.2). -/
  word1_eq : (imageOf w.data).2.read (gpow 1) = some input.word1
\end{leancode}

The image is \texttt{MemImage κ = Fin (\allowbreak{}2 \textasciicircum{} κ) → E} (\texttt{Memory.\allowbreak{}lean:77}): the word at index 0 is one element of \texttt{E}, \textbf{three limbs}, and it must equal \texttt{input.\allowbreak{}word0}, whose limb 2 is \texttt{0}. The memory block's rows are tied to that image limb by limb by \texttt{SeedRowsAreTheImage} (\texttt{Statement.\allowbreak{}lean:279-{}283}: the row of index \texttt{i} is \texttt{(\allowbreak{}i\allowbreak{}d\allowbreak{}x\allowbreak{} i\allowbreak{},\allowbreak{} c\allowbreak{}n\allowbreak{}t\allowbreak{}F\allowbreak{}i\allowbreak{}n\allowbreak{} i\allowbreak{},\allowbreak{} [\allowbreak{}l\allowbreak{}i\allowbreak{}m\allowbreak{}b\allowbreak{} 0\allowbreak{},\allowbreak{} l\allowbreak{}i\allowbreak{}m\allowbreak{}b\allowbreak{} 1\allowbreak{},\allowbreak{} l\allowbreak{}i\allowbreak{}m\allowbreak{}b\allowbreak{} 2\allowbreak{}]\allowbreak{})} of \texttt{(\allowbreak{}imageOf w.\allowbreak{}data).\allowbreak{}2 i}). So \texttt{SatisfiedBy} requires, of the committed columns \texttt{MEM\_\allowbreak{}LO,\allowbreak{} MEM\_\allowbreak{}HI,\allowbreak{} MEM\_\allowbreak{}TOP} at row 0, the values \texttt{(\allowbreak{}input₀,\allowbreak{} input₁,\allowbreak{} 0)}, and at row 1 \texttt{(\allowbreak{}input₂,\allowbreak{} input₃,\allowbreak{} 0)}. \textbf{The link holds}: the top cells are required to be zero, as three-limb values.
\item \emph{What \texttt{satisfiedBy\_\allowbreak{}witnessOf} must derive it from.} \texttt{witnessOf} \enquote{rebuilds … the image from the memory columns} (\texttt{p\allowbreak{}r\allowbreak{}o\allowbreak{}t\allowbreak{}o\allowbreak{}c\allowbreak{}o\allowbreak{}l\allowbreak{}-{}\allowbreak{}b\allowbreak{}l\allowbreak{}u\allowbreak{}e\allowbreak{}p\allowbreak{}r\allowbreak{}i\allowbreak{}n\allowbreak{}t\allowbreak{}.\allowbreak{}m\allowbreak{}d\allowbreak{}:\allowbreak{}8\allowbreak{}6\allowbreak{}3\allowbreak{}-{}\allowbreak{}8\allowbreak{}6\allowbreak{}4}), so \texttt{(\allowbreak{}i\allowbreak{}m\allowbreak{}a\allowbreak{}g\allowbreak{}e\allowbreak{}O\allowbreak{}f\allowbreak{} (\allowbreak{}w\allowbreak{}i\allowbreak{}t\allowbreak{}n\allowbreak{}e\allowbreak{}s\allowbreak{}s\allowbreak{}O\allowbreak{}f\allowbreak{} q\allowbreak{})\allowbreak{}.\allowbreak{}d\allowbreak{}a\allowbreak{}t\allowbreak{}a\allowbreak{})\allowbreak{}.\allowbreak{}2\allowbreak{} 0} is the word \texttt{o\allowbreak{}f\allowbreak{}L\allowbreak{}i\allowbreak{}m\allowbreak{}b\allowbreak{}s\allowbreak{} (\allowbreak{}m\allowbreak{}e\allowbreak{}m\allowbreak{}\_\allowbreak{}0\allowbreak{}[\allowbreak{}0\allowbreak{}]\allowbreak{})\allowbreak{} (\allowbreak{}m\allowbreak{}e\allowbreak{}m\allowbreak{}\_\allowbreak{}1\allowbreak{}[\allowbreak{}0\allowbreak{}]\allowbreak{})\allowbreak{} (\allowbreak{}m\allowbreak{}e\allowbreak{}m\allowbreak{}\_\allowbreak{}2\allowbreak{}[\allowbreak{}0\allowbreak{}]\allowbreak{})} read off \texttt{q}. \texttt{word0\_\allowbreak{}eq} then needs \texttt{mem\_\allowbreak{}0[0] = input₀}, \texttt{mem\_\allowbreak{}1[0] = input₁} \textbf{and \texttt{mem\_\allowbreak{}2[0] = 0}}. \texttt{M\allowbreak{}3\allowbreak{}H\allowbreak{}o\allowbreak{}l\allowbreak{}d\allowbreak{}s\allowbreak{} (\allowbreak{}l\allowbreak{}e\allowbreak{}a\allowbreak{}n\allowbreak{}I\allowbreak{}s\allowbreak{}a\allowbreak{}I\allowbreak{}n\allowbreak{}s\allowbreak{}t\allowbreak{}a\allowbreak{}n\allowbreak{}c\allowbreak{}e\allowbreak{} p\allowbreak{}r\allowbreak{}o\allowbreak{}g\allowbreak{} s\allowbreak{})\allowbreak{} i\allowbreak{}n\allowbreak{}p\allowbreak{}u\allowbreak{}t\allowbreak{} q} supplies exactly \texttt{PublicLinesHold}, i.e. the cells of the listed lines. With three lines the three equalities are there. With two lines the third is not, and \texttt{satisfiedBy\_\allowbreak{}witnessOf} is \textbf{false}: the stack that is honest except for \texttt{mem\_\allowbreak{}2[0] = 1} satisfies \texttt{M3Holds} of the two-line instance (nothing else reads that cell as a constraint; the bus carries it consistently from the seed row through every read to the finalize row), and its witness has \texttt{(\allowbreak{}i\allowbreak{}m\allowbreak{}a\allowbreak{}g\allowbreak{}e\allowbreak{}O\allowbreak{}f\allowbreak{} …\allowbreak{})\allowbreak{}.\allowbreak{}2\allowbreak{}.\allowbreak{}r\allowbreak{}e\allowbreak{}a\allowbreak{}d\allowbreak{} (\allowbreak{}g\allowbreak{}p\allowbreak{}o\allowbreak{}w\allowbreak{} 0\allowbreak{})\allowbreak{} =\allowbreak{} s\allowbreak{}o\allowbreak{}m\allowbreak{}e\allowbreak{} (\allowbreak{}o\allowbreak{}f\allowbreak{}L\allowbreak{}i\allowbreak{}m\allowbreak{}b\allowbreak{}s\allowbreak{} i\allowbreak{}n\allowbreak{}p\allowbreak{}u\allowbreak{}t\allowbreak{}₀\allowbreak{} i\allowbreak{}n\allowbreak{}p\allowbreak{}u\allowbreak{}t\allowbreak{}₁\allowbreak{} 1\allowbreak{})\allowbreak{} ≠\allowbreak{} s\allowbreak{}o\allowbreak{}m\allowbreak{}e\allowbreak{} i\allowbreak{}n\allowbreak{}p\allowbreak{}u\allowbreak{}t\allowbreak{}.\allowbreak{}w\allowbreak{}o\allowbreak{}r\allowbreak{}d\allowbreak{}0} (\texttt{limb\_\allowbreak{}ofLimbs}, \texttt{Field.\allowbreak{}lean:120}). A \texttt{witnessOf} that zeroed the top limb itself would break \texttt{mem\_\allowbreak{}balanced} or \texttt{SeedRowsAreTheImage} on any stack where a table reads cell 0, so no definition of \texttt{witnessOf} rescues the theorem for all stacks.
\item \emph{What the semantics says.} \texttt{ValidExecution prog input t} (\texttt{S\allowbreak{}e\allowbreak{}m\allowbreak{}a\allowbreak{}n\allowbreak{}t\allowbreak{}i\allowbreak{}c\allowbreak{}s\allowbreak{}/\allowbreak{}E\allowbreak{}x\allowbreak{}e\allowbreak{}c\allowbreak{}u\allowbreak{}t\allowbreak{}i\allowbreak{}o\allowbreak{}n\allowbreak{}.\allowbreak{}l\allowbreak{}e\allowbreak{}a\allowbreak{}n\allowbreak{}:\allowbreak{}1\allowbreak{}1\allowbreak{}0\allowbreak{}-{}\allowbreak{}1\allowbreak{}1\allowbreak{}3}) contains \texttt{HasPublicBoundary}: \texttt{t\allowbreak{}.\allowbreak{}i\allowbreak{}m\allowbreak{}a\allowbreak{}g\allowbreak{}e\allowbreak{}.\allowbreak{}r\allowbreak{}e\allowbreak{}a\allowbreak{}d\allowbreak{} (\allowbreak{}g\allowbreak{}p\allowbreak{}o\allowbreak{}w\allowbreak{} 0\allowbreak{})\allowbreak{} =\allowbreak{} s\allowbreak{}o\allowbreak{}m\allowbreak{}e\allowbreak{} i\allowbreak{}n\allowbreak{}p\allowbreak{}u\allowbreak{}t\allowbreak{}.\allowbreak{}w\allowbreak{}o\allowbreak{}r\allowbreak{}d\allowbreak{}0}, the same three-limb equality on the trace's image; \texttt{constraintSoundness} transports \texttt{word0\_\allowbreak{}eq} to it through \texttt{A\allowbreak{}s\allowbreak{}s\allowbreak{}i\allowbreak{}g\allowbreak{}n\allowbreak{}m\allowbreak{}e\allowbreak{}n\allowbreak{}t\allowbreak{}R\allowbreak{}e\allowbreak{}p\allowbreak{}r\allowbreak{}e\allowbreak{}s\allowbreak{}e\allowbreak{}n\allowbreak{}t\allowbreak{}s\allowbreak{}.\allowbreak{}i\allowbreak{}m\allowbreak{}a\allowbreak{}g\allowbreak{}e\allowbreak{}\_\allowbreak{}e\allowbreak{}q} (\texttt{Statement.\allowbreak{}lean:360}). So a stack with \texttt{mem\_\allowbreak{}2[0] = 1} accepted by a two-line protocol would be extracted to a witness with no valid execution on \texttt{input} (a program that reads cell 0 and branches on its top limb has an accepting stack and no valid execution on any \texttt{PublicInput}), and \texttt{b\allowbreak{}a\allowbreak{}s\allowbreak{}e\allowbreak{}V\allowbreak{}e\allowbreak{}r\allowbreak{}i\allowbreak{}f\allowbreak{}i\allowbreak{}e\allowbreak{}r\allowbreak{}\_\allowbreak{}e\allowbreak{}x\allowbreak{}t\allowbreak{}r\allowbreak{}a\allowbreak{}c\allowbreak{}t\allowbreak{}s\allowbreak{}E\allowbreak{}x\allowbreak{}e\allowbreak{}c\allowbreak{}u\allowbreak{}t\allowbreak{}i\allowbreak{}o\allowbreak{}n} would be false.
\end{enumerate}

\textbf{Anchor.} If the third line were dropped from \texttt{leanIsaInstance}, the theorem left without a proof is \texttt{satisfiedBy\_\allowbreak{}witnessOf} (its conjunct \texttt{word0\_\allowbreak{}eq}; and \texttt{word1\_\allowbreak{}eq} for cell 1), and through it T4's \texttt{b\allowbreak{}a\allowbreak{}s\allowbreak{}e\allowbreak{}V\allowbreak{}e\allowbreak{}r\allowbreak{}i\allowbreak{}f\allowbreak{}i\allowbreak{}e\allowbreak{}r\allowbreak{}\_\allowbreak{}e\allowbreak{}x\allowbreak{}t\allowbreak{}r\allowbreak{}a\allowbreak{}c\allowbreak{}t\allowbreak{}s\allowbreak{}E\allowbreak{}x\allowbreak{}e\allowbreak{}c\allowbreak{}u\allowbreak{}t\allowbreak{}i\allowbreak{}o\allowbreak{}n}. Nothing in \texttt{LeanerVM/\allowbreak{}Protocol/} would change: the phase, the spine and both master theorems would still hold, on the two-line instance, verbatim. The check on the top limb is therefore \textbf{not} a verifier check that the proof-system theorems can catch; it is a fact about which lines the leanISA instance lists, caught one layer down, at the wall between the proof system and the arithmetization, by a theorem that is not yet built. This is the right place for it (the instance is where leanISA enters), and it means the audit of \enquote{does the verifier check the top limb} is an audit of \texttt{leanIsaInstance.\allowbreak{}publicLines} against \texttt{SatisfiedBy.\allowbreak{}word0\_\allowbreak{}eq}, i.e. of a definition, plus the existence of \texttt{satisfiedBy\_\allowbreak{}witnessOf}.

\subsubsection{In the tests: acceptance test 10's witness exists today (D.2)}\label{sec:gen-pubinput-D2}

Acceptance test 10 (\texttt{p\allowbreak{}r\allowbreak{}o\allowbreak{}t\allowbreak{}o\allowbreak{}c\allowbreak{}o\allowbreak{}l\allowbreak{}-{}\allowbreak{}b\allowbreak{}l\allowbreak{}u\allowbreak{}e\allowbreak{}p\allowbreak{}r\allowbreak{}i\allowbreak{}n\allowbreak{}t\allowbreak{}.\allowbreak{}m\allowbreak{}d\allowbreak{}:\allowbreak{}1\allowbreak{}2\allowbreak{}9\allowbreak{}3\allowbreak{}-{}\allowbreak{}1\allowbreak{}2\allowbreak{}9\allowbreak{}5} at \texttt{b435631}): \enquote{The claim \texttt{mem\_\allowbreak{}2(\allowbreak{}r\_\allowbreak{}m,\allowbreak{} 0…) = 0} is pooled although no scalar is sent for it; omitting it lets a non-canonical public word pass the line. Witness: the Layer 8 mutation with a nonzero \texttt{mem\_\allowbreak{}2} at cell 0 and the third claim dropped.} The test file at \texttt{b435631} (\texttt{t\allowbreak{}e\allowbreak{}s\allowbreak{}t\allowbreak{}s\allowbreak{}/\allowbreak{}L\allowbreak{}e\allowbreak{}a\allowbreak{}n\allowbreak{}e\allowbreak{}r\allowbreak{}V\allowbreak{}M\allowbreak{}T\allowbreak{}e\allowbreak{}s\allowbreak{}t\allowbreak{}s\allowbreak{}/\allowbreak{}P\allowbreak{}r\allowbreak{}o\allowbreak{}t\allowbreak{}o\allowbreak{}c\allowbreak{}o\allowbreak{}l\allowbreak{}/\allowbreak{}P\allowbreak{}u\allowbreak{}b\allowbreak{}l\allowbreak{}i\allowbreak{}c\allowbreak{}I\allowbreak{}n\allowbreak{}p\allowbreak{}u\allowbreak{}t\allowbreak{}.\allowbreak{}l\allowbreak{}e\allowbreak{}a\allowbreak{}n\allowbreak{}:\allowbreak{}1\allowbreak{}5\allowbreak{}1\allowbreak{}-{}\allowbreak{}1\allowbreak{}8\allowbreak{}0}) has the literal three-line configuration (the earlier review's item \enquote{exercised nowhere} was met):

\begin{leancode}
/-- Three lines shaped like the memory limbs: two with their value sent, and a third with cells
`(0, 0)` and no value sent. -/
abbrev threeLimbs : M3Instance :=
  { toy with publicLines := fun v ↦
      [⟨⟨0, 0⟩, v, 1, true, by decide⟩, ⟨⟨0, 1⟩, 1, 1, true, by decide⟩,
        ⟨⟨0, 2⟩, 0, 0, false, by decide⟩] }
…
/-- The top column's cell 0 changed to `1`. -/
def badTopLimb : Column 3 := ⟨#v[1, 1, 1, 1, 1, 0, 0, 0]⟩
…
-- The wrong top cell passes the check, which sees the two sent values only, and fails the third
-- claim; the pool without the third claim accepts the stack.
#guard check threeLimbs stmtLimbs y (trueValues threeLimbs badTopLimb 1 y)
#guard ¬ ∀ c ∈ (pooled threeLimbs stmtLimbs y).2.columns, c.Holds badTopLimb
#guard ∀ c ∈ ((pooled threeLimbs stmtLimbs y).2.columns.take 2), c.Holds badTopLimb
\end{leancode}

It rejects what the acceptance test says: the check accepts the stack with the wrong top cell (no scalar sees it), the third claim fails of it, and the pool without the third claim (the \enquote{mutation with the third claim dropped}, written as \texttt{.\allowbreak{}take 2}) accepts it. Two limits:

\begin{itemize}
\item it is a \texttt{\#guard} on the \emph{pool}, not on the theorems: it shows the third claim is what rejects the stack, not that a phase without it has no \texttt{Phase.\allowbreak{}Security} (that is the refutation of C.8, on \texttt{noneSent});
\item it is on a toy-shaped instance built in the test, not on \texttt{leanIsaInstance} (not built). The acceptance test's real content, that leanISA's instance lists \texttt{mem\_\allowbreak{}2} with cells \texttt{(\allowbreak{}0,\allowbreak{} 0)}, is checkable only when Layer 3 lands; the blueprint's prose says it (\texttt{:834-{}836}).
\end{itemize}

\subsubsection{In the deployed verifiers: the top limb rides no scalar and is pooled at \texttt{0} (D.3)}\label{sec:gen-pubinput-D3}

Rust prover (\texttt{cpu/\allowbreak{}mod.\allowbreak{}rs:606-{}616}): \enquote{The top limb of both public words is zero, so its evaluation is zero at every \texttt{r\_\allowbreak{}pi} and rides no scalar}; \texttt{p\allowbreak{}i\allowbreak{}\_\allowbreak{}l\allowbreak{}i\allowbreak{}m\allowbreak{}b\allowbreak{}s\allowbreak{} =\allowbreak{} [\allowbreak{}i\allowbreak{}n\allowbreak{}t\allowbreak{}e\allowbreak{}r\allowbreak{}p\allowbreak{}\_\allowbreak{}k\allowbreak{}(\allowbreak{}…\allowbreak{}c\allowbreak{}0\allowbreak{}…\allowbreak{})\allowbreak{},\allowbreak{} i\allowbreak{}n\allowbreak{}t\allowbreak{}e\allowbreak{}r\allowbreak{}p\allowbreak{}\_\allowbreak{}k\allowbreak{}(\allowbreak{}…\allowbreak{}c\allowbreak{}1\allowbreak{}…\allowbreak{})\allowbreak{},\allowbreak{} F\allowbreak{}1\allowbreak{}9\allowbreak{}2\allowbreak{}:\allowbreak{}:\allowbreak{}Z\allowbreak{}E\allowbreak{}R\allowbreak{}O\allowbreak{}]}, only \texttt{pi\_\allowbreak{}limbs[.\allowbreak{}.\allowbreak{}2]} are added to the stream. Rust verifier (\texttt{cpu/\allowbreak{}mod.\allowbreak{}rs:745-{}756}): two scalars read into \texttt{pi\_\allowbreak{}limbs[.\allowbreak{}.\allowbreak{}2]}, \texttt{pi\_\allowbreak{}limbs[2]} stays \texttt{F192::ZERO}, the combined check, then \texttt{finish\_\allowbreak{}claims}, whose \texttt{bind\_\allowbreak{}pi\_\allowbreak{}claim} (\texttt{:674-{}682}) pools \texttt{[MEM\_\allowbreak{}LO,\allowbreak{} MEM\_\allowbreak{}HI,\allowbreak{} MEM\_\allowbreak{}TOP]} at \texttt{(\allowbreak{}r,\allowbreak{} 0,\allowbreak{} …)} with the three \texttt{pi\_\allowbreak{}limbs}. Python (\texttt{verifier.\allowbreak{}py:1398-{}1402}): \texttt{p\allowbreak{}u\allowbreak{}b\allowbreak{}l\allowbreak{}i\allowbreak{}c\allowbreak{}\_\allowbreak{}l\allowbreak{}i\allowbreak{}m\allowbreak{}b\allowbreak{}s\allowbreak{} =\allowbreak{} (\allowbreak{}*\allowbreak{}t\allowbreak{}r\allowbreak{}a\allowbreak{}n\allowbreak{}s\allowbreak{}c\allowbreak{}r\allowbreak{}i\allowbreak{}p\allowbreak{}t\allowbreak{}.\allowbreak{}n\allowbreak{}e\allowbreak{}x\allowbreak{}t\allowbreak{}\_\allowbreak{}s\allowbreak{}c\allowbreak{}a\allowbreak{}l\allowbreak{}a\allowbreak{}r\allowbreak{}s\allowbreak{}(\allowbreak{}2\allowbreak{})\allowbreak{},\allowbreak{} Z\allowbreak{}E\allowbreak{}R\allowbreak{}O\allowbreak{})}, the combined check \texttt{p\allowbreak{}o\allowbreak{}l\allowbreak{}y\allowbreak{}\_\allowbreak{}e\allowbreak{}v\allowbreak{}a\allowbreak{}l\allowbreak{}(\allowbreak{}p\allowbreak{}u\allowbreak{}b\allowbreak{}l\allowbreak{}i\allowbreak{}c\allowbreak{}\_\allowbreak{}l\allowbreak{}i\allowbreak{}m\allowbreak{}b\allowbreak{}s\allowbreak{},\allowbreak{} Y\allowbreak{})\allowbreak{} =\allowbreak{}=\allowbreak{} m\allowbreak{}u\allowbreak{}l\allowbreak{}t\allowbreak{}i\allowbreak{}l\allowbreak{}i\allowbreak{}n\allowbreak{}e\allowbreak{}a\allowbreak{}r\allowbreak{}\_\allowbreak{}e\allowbreak{}v\allowbreak{}a\allowbreak{}l\allowbreak{}(\allowbreak{}p\allowbreak{}u\allowbreak{}b\allowbreak{}l\allowbreak{}i\allowbreak{}c\allowbreak{}\_\allowbreak{}i\allowbreak{}n\allowbreak{}p\allowbreak{}u\allowbreak{}t\allowbreak{}.\allowbreak{}h\allowbreak{}a\allowbreak{}l\allowbreak{}v\allowbreak{}e\allowbreak{}s\allowbreak{}(\allowbreak{})\allowbreak{},\allowbreak{} [\allowbreak{}p\allowbreak{}u\allowbreak{}b\allowbreak{}l\allowbreak{}i\allowbreak{}c\allowbreak{}\_\allowbreak{}c\allowbreak{}h\allowbreak{}a\allowbreak{}l\allowbreak{}l\allowbreak{}e\allowbreak{}n\allowbreak{}g\allowbreak{}e\allowbreak{}]\allowbreak{})}, then three claims on \texttt{MEMORY\_\allowbreak{}0,\allowbreak{} MEMORY\_\allowbreak{}1,\allowbreak{} MEMORY\_\allowbreak{}2}. Recursion guest (\texttt{c\allowbreak{}r\allowbreak{}a\allowbreak{}t\allowbreak{}e\allowbreak{}s\allowbreak{}/\allowbreak{}r\allowbreak{}e\allowbreak{}c\allowbreak{}\_\allowbreak{}a\allowbreak{}g\allowbreak{}g\allowbreak{}r\allowbreak{}e\allowbreak{}g\allowbreak{}a\allowbreak{}t\allowbreak{}i\allowbreak{}o\allowbreak{}n\allowbreak{}/\allowbreak{}g\allowbreak{}u\allowbreak{}e\allowbreak{}s\allowbreak{}t\allowbreak{}s\allowbreak{}/\allowbreak{}a\allowbreak{}g\allowbreak{}g\allowbreak{}r\allowbreak{}e\allowbreak{}g\allowbreak{}a\allowbreak{}t\allowbreak{}e\allowbreak{}.\allowbreak{}p\allowbreak{}y\allowbreak{}:\allowbreak{}1\allowbreak{}6\allowbreak{}8\allowbreak{}0\allowbreak{}-{}\allowbreak{}1\allowbreak{}6\allowbreak{}8\allowbreak{}9}): \texttt{a\allowbreak{}s\allowbreak{}s\allowbreak{}e\allowbreak{}r\allowbreak{}t\allowbreak{} m\allowbreak{}e\allowbreak{}m\allowbreak{} =\allowbreak{}=\allowbreak{} m\allowbreak{}e\allowbreak{}m\allowbreak{}\_\allowbreak{}l\allowbreak{}o\allowbreak{} +\allowbreak{} m\allowbreak{}e\allowbreak{}m\allowbreak{}\_\allowbreak{}h\allowbreak{}i\allowbreak{} *\allowbreak{} Y\allowbreak{}\_\allowbreak{}T\allowbreak{}O\allowbreak{}W\allowbreak{}E\allowbreak{}R} and three pool entries, the third \texttt{0}. The statement itself cannot have a nonzero top limb (\texttt{read\_\allowbreak{}public}, D.1 item 1; the Python's \texttt{Digest} is 256 bits, \texttt{verifier.\allowbreak{}py:204-{}219}). The Lean phase pools the third claim at \texttt{l\allowbreak{}i\allowbreak{}n\allowbreak{}e\allowbreak{}V\allowbreak{}a\allowbreak{}l\allowbreak{}u\allowbreak{}e\allowbreak{} r\allowbreak{} l\allowbreak{}₂\allowbreak{} =\allowbreak{} (\allowbreak{}1\allowbreak{} +\allowbreak{} r\allowbreak{})\allowbreak{}·\allowbreak{}0\allowbreak{} +\allowbreak{} r\allowbreak{}·\allowbreak{}0\allowbreak{} =\allowbreak{} 0} (test \texttt{limbValues}), the same claim.

So the deployed verifiers do check the top limb, through the third pooled claim; the specification says so too (\texttt{0\allowbreak{}8\allowbreak{}-{}\allowbreak{}e\allowbreak{}n\allowbreak{}d\allowbreak{}-{}\allowbreak{}t\allowbreak{}o\allowbreak{}-{}\allowbreak{}e\allowbreak{}n\allowbreak{}d\allowbreak{}-{}\allowbreak{}p\allowbreak{}r\allowbreak{}o\allowbreak{}t\allowbreak{}o\allowbreak{}c\allowbreak{}o\allowbreak{}l\allowbreak{}.\allowbreak{}t\allowbreak{}e\allowbreak{}x\allowbreak{}:\allowbreak{}8\allowbreak{}3\allowbreak{}-{}\allowbreak{}8\allowbreak{}4}: \enquote{pools them, with \$0\$ for the top limb, as claims on \$\textbackslash{}mem\_0,\textbackslash{}mem\_1,\textbackslash{}mem\_2\$}). The Lean matches. What no document says in one place is that the check lives in the instance's list of lines and is enforced by the adaptor's theorem, not by the phase's; D.1 is that sentence.

\subsection{Layer 1's probes (I)}\label{sec:gen-probes-layer1}

\par\noindent{\small\textit{Cross-references.} A reference such as \enquote{(D.5)} or \enquote{section B} in this part is to a section of the dossier \file{dossiers/code-layer1.md}, whose codes the headings keep in parentheses: A (catalogue) is \cref{sec:gen-cat-layer1}; E (audit surface) is \cref{sec:gen-surface-layer1}; F (conformance) is \cref{sec:gen-conf-layer1}; G (findings) is \cref{sec:gen-findings-layer1}; I (probes) is \cref{sec:gen-probes-layer1}; B (faithfulness), C (statements), D (the wall and the To-folders) and H (the upgrade) is not reproduced in these fragments; they are in the dossier.}\par

All under \texttt{.\allowbreak{}c\allowbreak{}l\allowbreak{}a\allowbreak{}u\allowbreak{}d\allowbreak{}e\allowbreak{}/\allowbreak{}r\allowbreak{}e\allowbreak{}p\allowbreak{}o\allowbreak{}r\allowbreak{}t\allowbreak{}s\allowbreak{}/\allowbreak{}b\allowbreak{}l\allowbreak{}u\allowbreak{}e\allowbreak{}p\allowbreak{}r\allowbreak{}i\allowbreak{}n\allowbreak{}t\allowbreak{}-{}\allowbreak{}r\allowbreak{}e\allowbreak{}v\allowbreak{}i\allowbreak{}e\allowbreak{}w\allowbreak{}/\allowbreak{}p\allowbreak{}r\allowbreak{}o\allowbreak{}b\allowbreak{}e\allowbreak{}s\allowbreak{}/\allowbreak{}c\allowbreak{}o\allowbreak{}d\allowbreak{}e\allowbreak{}-{}\allowbreak{}l\allowbreak{}a\allowbreak{}y\allowbreak{}e\allowbreak{}r\allowbreak{}1\allowbreak{}/}, written by the interrupted first agent of this task and run once each by this one on 2026-09-30, each output saved next to its probe (\texttt{*.\allowbreak{}out}). \textbf{How they ran}, given that the owner moved the checkout at 09:18 (reflog: \texttt{reset} 09:18:19, \texttt{checkout main} 09:18:20, \texttt{pull} to \texttt{144c5aa} 09:18:26, back to the review branch 09:18:33, \texttt{merge main} 09:18:36; the dependency sources under \texttt{.\allowbreak{}lake/\allowbreak{}packages/} were switched to the new pins at 09:18:47-09:19:05, their \texttt{.\allowbreak{}olean} files were not rebuilt):

\begin{itemize}
\item \texttt{offsets.\allowbreak{}py}, \texttt{values.\allowbreak{}py}: \texttt{python3 -{}B <file>} at 09:16, exit 0. They import the pinned Python verifier read-only from \texttt{/\allowbreak{}h\allowbreak{}o\allowbreak{}m\allowbreak{}e\allowbreak{}/\allowbreak{}s\allowbreak{}c\allowbreak{}a\allowbreak{}r\allowbreak{}a\allowbreak{}v\allowbreak{}e\allowbreak{}n\allowbreak{}/\allowbreak{}D\allowbreak{}o\allowbreak{}c\allowbreak{}u\allowbreak{}m\allowbreak{}e\allowbreak{}n\allowbreak{}t\allowbreak{}s\allowbreak{}/\allowbreak{}l\allowbreak{}e\allowbreak{}a\allowbreak{}n\allowbreak{}E\allowbreak{}t\allowbreak{}h\allowbreak{}e\allowbreak{}r\allowbreak{}e\allowbreak{}u\allowbreak{}m\allowbreak{}/\allowbreak{}l\allowbreak{}e\allowbreak{}a\allowbreak{}n\allowbreak{}V\allowbreak{}M\allowbreak{}/\allowbreak{}p\allowbreak{}y\allowbreak{}t\allowbreak{}h\allowbreak{}o\allowbreak{}n\allowbreak{}-{}\allowbreak{}v\allowbreak{}e\allowbreak{}r\allowbreak{}i\allowbreak{}f\allowbreak{}i\allowbreak{}e\allowbreak{}r} (at \texttt{a386121f}).
\item \texttt{OffsetsProbe.\allowbreak{}lean} (finished 09:17:02) and \texttt{ValuesProbe.\allowbreak{}lean} (finished 09:18:17.99), with the brief's command \texttt{f\allowbreak{}l\allowbreak{}o\allowbreak{}c\allowbreak{}k\allowbreak{} .\allowbreak{}c\allowbreak{}l\allowbreak{}a\allowbreak{}u\allowbreak{}d\allowbreak{}e\allowbreak{}/\allowbreak{}r\allowbreak{}e\allowbreak{}p\allowbreak{}o\allowbreak{}r\allowbreak{}t\allowbreak{}s\allowbreak{}/\allowbreak{}b\allowbreak{}l\allowbreak{}u\allowbreak{}e\allowbreak{}p\allowbreak{}r\allowbreak{}i\allowbreak{}n\allowbreak{}t\allowbreak{}-{}\allowbreak{}r\allowbreak{}e\allowbreak{}v\allowbreak{}i\allowbreak{}e\allowbreak{}w\allowbreak{}/\allowbreak{}l\allowbreak{}o\allowbreak{}g\allowbreak{}s\allowbreak{}/\allowbreak{}l\allowbreak{}e\allowbreak{}a\allowbreak{}n\allowbreak{}.\allowbreak{}l\allowbreak{}o\allowbreak{}c\allowbreak{}k\allowbreak{} l\allowbreak{}a\allowbreak{}k\allowbreak{}e\allowbreak{} e\allowbreak{}n\allowbreak{}v\allowbreak{} l\allowbreak{}e\allowbreak{}a\allowbreak{}n\allowbreak{} <\allowbreak{}f\allowbreak{}i\allowbreak{}l\allowbreak{}e\allowbreak{}>}, both before the checkout moved: sources at \texttt{b435631}, libraries at the old pins, exit 0.
\item \texttt{DuplicatesProbe.\allowbreak{}lean}: the first attempt with \texttt{lake env lean} (09:18:22) found the file missing (the owner's checkout of \texttt{main} had removed the tracked probe files for thirteen seconds) and did nothing else. Both it and \texttt{StridedProbe.\allowbreak{}lean} were then run at 09:20 \textbf{not through \texttt{lake}} but with the old pin's Lean binary directly on the \texttt{b435631} build artefacts, which the move had not touched (no \texttt{.\allowbreak{}olean} under \texttt{.\allowbreak{}lake/} was newer than 2026-09-30 00:00, checked just before; \texttt{.\allowbreak{}lake/\allowbreak{}build} dates from the orchestrator's build of 2026-09-29 17:18, the packages' builds from 2026-09-28/29):

\begin{shellcode}
LP="$(pwd)/.lake/build/lib/lean"; for d in .lake/packages/*/.lake/build/lib/lean; do LP="$LP:$(pwd)/$d"; done
LEAN_PATH="$LP" timeout 590 flock .claude/reports/blueprint-review/logs/lean.lock \
  /home/scaraven/.elan/toolchains/leanprover--lean4---v4.33.1/bin/lean <probe>.lean
\end{shellcode}

Lean loads compiled \texttt{.\allowbreak{}olean} files only and never reads the (switched) library sources, so these runs elaborate against \texttt{b435631} and the old pins exactly as \texttt{lake env lean} did before the move; both exit 0. They were made before the coordinator's instruction to stop running Lean arrived; no Lean process was started after it. The orchestrator may prefer to treat these two as \enquote{written, not run through the brief's command}; the conclusions resting on them are marked in the text (\texttt{DuplicatesProbe}: D.2 and G.5; \texttt{StridedProbe}: B.8 and G.1) and each is also a short argument a reader can check by hand.
\end{itemize}

\subsubsection{\texttt{offsets.\allowbreak{}py} (I.1)}\label{sec:gen-layer1-I1}

\fhead{What it establishes}Runs the pinned Python verifier's own \texttt{stack\_\allowbreak{}offsets} and \texttt{build\_\allowbreak{}layout} (imported read-only from the leanVM checkout at \texttt{a386121f}) and a line-by-line transcription of the Rust \texttt{stack\_\allowbreak{}offsets} (\texttt{c\allowbreak{}r\allowbreak{}a\allowbreak{}t\allowbreak{}e\allowbreak{}s\allowbreak{}/\allowbreak{}l\allowbreak{}e\allowbreak{}a\allowbreak{}n\allowbreak{}\_\allowbreak{}v\allowbreak{}m\allowbreak{}/\allowbreak{}s\allowbreak{}r\allowbreak{}c\allowbreak{}/\allowbreak{}w\allowbreak{}i\allowbreak{}t\allowbreak{}n\allowbreak{}e\allowbreak{}s\allowbreak{}s\allowbreak{}.\allowbreak{}r\allowbreak{}s\allowbreak{}:\allowbreak{}6\allowbreak{}7\allowbreak{}-{}\allowbreak{}7\allowbreak{}9}), and prints the offsets and selectors of two configurations: the dossier's small configuration, and the pinned verifier's real layout for one admissible announcement (92 committed blocks, \texttt{stack\_\allowbreak{}log = 22}). Its numbers are the expected values of \texttt{OffsetsProbe.\allowbreak{}lean}.

\textbf{Command.} \texttt{p\allowbreak{}y\allowbreak{}t\allowbreak{}h\allowbreak{}o\allowbreak{}n\allowbreak{}3\allowbreak{} -{}\allowbreak{}B\allowbreak{} .\allowbreak{}c\allowbreak{}l\allowbreak{}a\allowbreak{}u\allowbreak{}d\allowbreak{}e\allowbreak{}/\allowbreak{}r\allowbreak{}e\allowbreak{}p\allowbreak{}o\allowbreak{}r\allowbreak{}t\allowbreak{}s\allowbreak{}/\allowbreak{}b\allowbreak{}l\allowbreak{}u\allowbreak{}e\allowbreak{}p\allowbreak{}r\allowbreak{}i\allowbreak{}n\allowbreak{}t\allowbreak{}-{}\allowbreak{}r\allowbreak{}e\allowbreak{}v\allowbreak{}i\allowbreak{}e\allowbreak{}w\allowbreak{}/\allowbreak{}p\allowbreak{}r\allowbreak{}o\allowbreak{}b\allowbreak{}e\allowbreak{}s\allowbreak{}/\allowbreak{}c\allowbreak{}o\allowbreak{}d\allowbreak{}e\allowbreak{}-{}\allowbreak{}l\allowbreak{}a\allowbreak{}y\allowbreak{}e\allowbreak{}r\allowbreak{}1\allowbreak{}/\allowbreak{}o\allowbreak{}f\allowbreak{}f\allowbreak{}s\allowbreak{}e\allowbreak{}t\allowbreak{}s\allowbreak{}.\allowbreak{}p\allowbreak{}y}

\begin{srccode}
"""Probe for task code-layer1, deliverable B.2.

Runs the pinned Python verifier's own `stack_offsets` and `build_layout` (imported read-only
from the leanVM checkout at a386121f; no bytecode cache is written) and a line-by-line
transcription of the Rust `stack_offsets` (crates/lean_vm/src/witness.rs:67-79), and prints
the offsets and selectors of two configurations.

Run:  python3 -B .claude/reports/blueprint-review/probes/code-layer1/offsets.py
"""

import sys

sys.dont_write_bytecode = True
sys.path.insert(0, "/home/scaraven/Documents/leanEthereum/leanVM/python-verifier")
import verifier as V  # noqa: E402


def rust_stack_offsets(kappas):
    """witness.rs:67-79, transcribed: `order.sort_by(|a,b| kappas[b].cmp(kappas[a]).then(a.cmp(b)))`."""
    n = len(kappas)
    order = [i for i in range(n) if kappas[i] is not None]
    order.sort(key=lambda i: (-kappas[i], i))
    offsets = [0] * n
    off = 0
    for i in order:
        offsets[i] = off
        off += 1 << kappas[i]
    return offsets, off, order


def report(name, kappas, names=None):
    print(f"== {name}")
    print("sizes in column-index order:", kappas)
    r_off, r_total, order = rust_stack_offsets(kappas)
    p_off, p_log = V.stack_offsets(kappas)
    print("Rust transcription offsets :", r_off, "total", r_total)
    print("Python verifier offsets    :", p_off, "log2_ceil(total)", p_log)
    assert r_off == p_off
    print("stacking order (column indices, first block first):", order)
    print("sorted sizes (the Blocks.size sequence)           :", [kappas[i] for i in order])
    print("offset per block, in stacking order               :", [r_off[i] for i in order])
    print("selector per block (offset >> size)               :", [r_off[i] >> kappas[i] for i in order])
    if names:
        for i in order:
            print(f"   block of column {i:3d} {names[i]:<14s} size {kappas[i]:2d} offset {r_off[i]:8d} selector {r_off[i] >> kappas[i]}")
    print()


# 1. The small configuration of the dossier: six shared columns and two table columns.
names = ["mem_0", "mem_1", "mem_2", "cntfin_mem", "cntfin_bc", "q_flock", "table col a", "table col b"]
report("small configuration", [4, 4, 4, 4, 2, 5, 4, 4], names)

# 2. The pinned verifier's real layout for one admissible announcement:
#    log_memory 16, table log-heights (xor, mul, set, deref, jump, blake2s) = (3, 16, 0, 5, 16, 3),
#    a bytecode of 2^4 instructions (the table has 2^(4+4) words; its content is irrelevant here).
bytecode = [V.K(0)] * (1 << 8)
layout = V.build_layout(bytecode, 16, (3, 16, 0, 5, 16, 3))
cols = []
for column, p in enumerate(layout.placements):
    cols.append((column, p.variables, p.index, p.low))
committed = [(c, k, off) for (c, k, off, low) in cols if low == 0]
virtual = [(c, k, off, low) for (c, k, off, low) in cols if low != 0]
print("== the pinned Python verifier's build_layout(log_memory=16, taus=(3,16,0,5,16,3), kbc=4)")
print("columns:", len(cols), "committed blocks:", len(committed), "strided (BLAKE2S limb) columns:", len(virtual))
print("stack_log:", layout.stack_log)
order = sorted(committed, key=lambda t: t[2])
print("stacking order, as (column index, size, offset):")
print(order)
print("sorted sizes:", [k for (_, k, _) in order])
print("offsets     :", [off for (_, _, off) in order])
print("column index of each block:", [c for (c, _, _) in order])
# the same sizes through the Rust transcription, virtual columns as None
kappas = [None if low != 0 else k for (_, k, _, low) in cols]
r_off, r_total, r_order = rust_stack_offsets(kappas)
assert [r_off[c] for (c, _, _) in committed] == [off for (_, _, off) in committed]
assert r_order == [c for (c, _, _) in order]
print("Rust transcription agrees on every committed column; total", r_total, "log2_ceil", V.log2_ceil(r_total))
print("global column bases of the six tables:", V.GLOBAL_COLUMN_BASES, "widths", V.TABLE_WIDTHS)
\end{srccode}

\textbf{Recorded output.}

\begin{srccode}
== small configuration
sizes in column-index order: [4, 4, 4, 4, 2, 5, 4, 4]
Rust transcription offsets : [32, 48, 64, 80, 128, 0, 96, 112] total 132
Python verifier offsets    : [32, 48, 64, 80, 128, 0, 96, 112] log2_ceil(total) 8
stacking order (column indices, first block first): [5, 0, 1, 2, 3, 6, 7, 4]
sorted sizes (the Blocks.size sequence)           : [5, 4, 4, 4, 4, 4, 4, 2]
offset per block, in stacking order               : [0, 32, 48, 64, 80, 96, 112, 128]
selector per block (offset >> size)               : [0, 2, 3, 4, 5, 6, 7, 32]
   block of column   5 q_flock        size  5 offset        0 selector 0
   block of column   0 mem_0          size  4 offset       32 selector 2
   block of column   1 mem_1          size  4 offset       48 selector 3
   block of column   2 mem_2          size  4 offset       64 selector 4
   block of column   3 cntfin_mem     size  4 offset       80 selector 5
   block of column   6 table col a    size  4 offset       96 selector 6
   block of column   7 table col b    size  4 offset      112 selector 7
   block of column   4 cntfin_bc      size  2 offset      128 selector 32

== the pinned Python verifier's build_layout(log_memory=16, taus=(3,16,0,5,16,3), kbc=4)
columns: 110 committed blocks: 92 strided (BLAKE2S limb) columns: 18
stack_log: 22
stacking order, as (column index, size, offset):
[(0, 16, 0), (1, 16, 65536), (2, 16, 131072), (3, 16, 196608), (21, 16, 262144), (22, 16, 327680), (23, 16, 393216), (24, 16, 458752), (25, 16, 524288), (26, 16, 589824), (27, 16, 655360), (28, 16, 720896), (29, 16, 786432), (30, 16, 851968), (31, 16, 917504), (32, 16, 983040), (33, 16, 1048576), (34, 16, 1114112), (35, 16, 1179648), (59, 16, 1245184), (60, 16, 1310720), (61, 16, 1376256), (62, 16, 1441792), (63, 16, 1507328), (64, 16, 1572864), (65, 16, 1638400), (66, 16, 1703936), (67, 16, 1769472), (68, 16, 1835008), (69, 16, 1900544), (70, 16, 1966080), (71, 16, 2031616), (72, 16, 2097152), (5, 11, 2162688), (44, 5, 2164736), (45, 5, 2164768), (46, 5, 2164800), (47, 5, 2164832), (48, 5, 2164864), (49, 5, 2164896), (50, 5, 2164928), (51, 5, 2164960), (52, 5, 2164992), (53, 5, 2165024), (54, 5, 2165056), (55, 5, 2165088), (56, 5, 2165120), (57, 5, 2165152), (58, 5, 2165184), (4, 4, 2165216), (6, 3, 2165232), (7, 3, 2165240), (8, 3, 2165248), (9, 3, 2165256), (10, 3, 2165264), (11, 3, 2165272), (12, 3, 2165280), (13, 3, 2165288), (14, 3, 2165296), (15, 3, 2165304), (16, 3, 2165312), (17, 3, 2165320), (18, 3, 2165328), (19, 3, 2165336), (20, 3, 2165344), (73, 3, 2165352), (74, 3, 2165360), (75, 3, 2165368), (76, 3, 2165376), (77, 3, 2165384), (78, 3, 2165392), (79, 3, 2165400), (80, 3, 2165408), (81, 3, 2165416), (100, 3, 2165424), (101, 3, 2165432), (102, 3, 2165440), (103, 3, 2165448), (104, 3, 2165456), (105, 3, 2165464), (106, 3, 2165472), (107, 3, 2165480), (108, 3, 2165488), (109, 3, 2165496), (36, 0, 2165504), (37, 0, 2165505), (38, 0, 2165506), (39, 0, 2165507), (40, 0, 2165508), (41, 0, 2165509), (42, 0, 2165510), (43, 0, 2165511)]
sorted sizes: [16, 16, 16, 16, 16, 16, 16, 16, 16, 16, 16, 16, 16, 16, 16, 16, 16, 16, 16, 16, 16, 16, 16, 16, 16, 16, 16, 16, 16, 16, 16, 16, 16, 11, 5, 5, 5, 5, 5, 5, 5, 5, 5, 5, 5, 5, 5, 5, 5, 4, 3, 3, 3, 3, 3, 3, 3, 3, 3, 3, 3, 3, 3, 3, 3, 3, 3, 3, 3, 3, 3, 3, 3, 3, 3, 3, 3, 3, 3, 3, 3, 3, 3, 3, 0, 0, 0, 0, 0, 0, 0, 0]
offsets     : [0, 65536, 131072, 196608, 262144, 327680, 393216, 458752, 524288, 589824, 655360, 720896, 786432, 851968, 917504, 983040, 1048576, 1114112, 1179648, 1245184, 1310720, 1376256, 1441792, 1507328, 1572864, 1638400, 1703936, 1769472, 1835008, 1900544, 1966080, 2031616, 2097152, 2162688, 2164736, 2164768, 2164800, 2164832, 2164864, 2164896, 2164928, 2164960, 2164992, 2165024, 2165056, 2165088, 2165120, 2165152, 2165184, 2165216, 2165232, 2165240, 2165248, 2165256, 2165264, 2165272, 2165280, 2165288, 2165296, 2165304, 2165312, 2165320, 2165328, 2165336, 2165344, 2165352, 2165360, 2165368, 2165376, 2165384, 2165392, 2165400, 2165408, 2165416, 2165424, 2165432, 2165440, 2165448, 2165456, 2165464, 2165472, 2165480, 2165488, 2165496, 2165504, 2165505, 2165506, 2165507, 2165508, 2165509, 2165510, 2165511]
column index of each block: [0, 1, 2, 3, 21, 22, 23, 24, 25, 26, 27, 28, 29, 30, 31, 32, 33, 34, 35, 59, 60, 61, 62, 63, 64, 65, 66, 67, 68, 69, 70, 71, 72, 5, 44, 45, 46, 47, 48, 49, 50, 51, 52, 53, 54, 55, 56, 57, 58, 4, 6, 7, 8, 9, 10, 11, 12, 13, 14, 15, 16, 17, 18, 19, 20, 73, 74, 75, 76, 77, 78, 79, 80, 81, 100, 101, 102, 103, 104, 105, 106, 107, 108, 109, 36, 37, 38, 39, 40, 41, 42, 43]
Rust transcription agrees on every committed column; total 2165512 log2_ceil 22
global column bases of the six tables: (6, 21, 36, 44, 59, 73) widths (15, 15, 8, 15, 14, 37)
\end{srccode}

\subsubsection{\texttt{OffsetsProbe.\allowbreak{}lean} (I.2)}\label{sec:gen-layer1-I2}

\fhead{What it establishes}\texttt{Blocks.\allowbreak{}offset} and \texttt{Blocks.\allowbreak{}selector} on sorted sizes reproduce leanVM's offsets once the columns are mapped by leanVM's order, and the tables-first order gives the same \texttt{Blocks} and other offsets (B.2).

\textbf{Command.} \texttt{f\allowbreak{}l\allowbreak{}o\allowbreak{}c\allowbreak{}k\allowbreak{} .\allowbreak{}c\allowbreak{}l\allowbreak{}a\allowbreak{}u\allowbreak{}d\allowbreak{}e\allowbreak{}/\allowbreak{}r\allowbreak{}e\allowbreak{}p\allowbreak{}o\allowbreak{}r\allowbreak{}t\allowbreak{}s\allowbreak{}/\allowbreak{}b\allowbreak{}l\allowbreak{}u\allowbreak{}e\allowbreak{}p\allowbreak{}r\allowbreak{}i\allowbreak{}n\allowbreak{}t\allowbreak{}-{}\allowbreak{}r\allowbreak{}e\allowbreak{}v\allowbreak{}i\allowbreak{}e\allowbreak{}w\allowbreak{}/\allowbreak{}l\allowbreak{}o\allowbreak{}g\allowbreak{}s\allowbreak{}/\allowbreak{}l\allowbreak{}e\allowbreak{}a\allowbreak{}n\allowbreak{}.\allowbreak{}l\allowbreak{}o\allowbreak{}c\allowbreak{}k\allowbreak{} l\allowbreak{}a\allowbreak{}k\allowbreak{}e\allowbreak{} e\allowbreak{}n\allowbreak{}v\allowbreak{} l\allowbreak{}e\allowbreak{}a\allowbreak{}n\allowbreak{} .\allowbreak{}c\allowbreak{}l\allowbreak{}a\allowbreak{}u\allowbreak{}d\allowbreak{}e\allowbreak{}/\allowbreak{}r\allowbreak{}e\allowbreak{}p\allowbreak{}o\allowbreak{}r\allowbreak{}t\allowbreak{}s\allowbreak{}/\allowbreak{}b\allowbreak{}l\allowbreak{}u\allowbreak{}e\allowbreak{}p\allowbreak{}r\allowbreak{}i\allowbreak{}n\allowbreak{}t\allowbreak{}-{}\allowbreak{}r\allowbreak{}e\allowbreak{}v\allowbreak{}i\allowbreak{}e\allowbreak{}w\allowbreak{}/\allowbreak{}p\allowbreak{}r\allowbreak{}o\allowbreak{}b\allowbreak{}e\allowbreak{}s\allowbreak{}/\allowbreak{}c\allowbreak{}o\allowbreak{}d\allowbreak{}e\allowbreak{}-{}\allowbreak{}l\allowbreak{}a\allowbreak{}y\allowbreak{}e\allowbreak{}r\allowbreak{}1\allowbreak{}/\allowbreak{}O\allowbreak{}f\allowbreak{}f\allowbreak{}s\allowbreak{}e\allowbreak{}t\allowbreak{}s\allowbreak{}P\allowbreak{}r\allowbreak{}o\allowbreak{}b\allowbreak{}e\allowbreak{}.\allowbreak{}l\allowbreak{}e\allowbreak{}a\allowbreak{}n}

\begin{leancode}
/-
  Probe for task code-layer1, deliverable B.2.

  leanVM's `stack_offsets` (`crates/lean_vm/src/witness.rs:67-79`, `python-verifier/verifier.py:
  305-311`) transcribed as a list function, and Layer 1's `Blocks.offset` / `Blocks.selector`
  fed with the order it produces. The expected numbers are the output of `offsets.py`, which
  runs the pinned Python verifier.

  Run from the repository root:
    flock .claude/reports/blueprint-review/logs/lean.lock \
      lake env lean .claude/reports/blueprint-review/probes/code-layer1/OffsetsProbe.lean
-/
import LeanerVM.Protocol.ToCompPoly.Stacking

open LeanerVM.Protocol

namespace Probe

/-! ## `stack_offsets`, transcribed -/

/-- The stacking order: column indices sorted by size, largest first, ties by index. -/
def stackOrder (kappas : List ℕ) : List ℕ :=
  (List.range kappas.length).mergeSort fun a b ↦
    decide (kappas[b]! < kappas[a]!) || (kappas[a]! == kappas[b]! && decide (a ≤ b))

/-- Prefix sums of the heights along a list of sizes. -/
def prefixOffsets : List ℕ → ℕ → List ℕ
  | [], _ => []
  | k :: ks, off => off :: prefixOffsets ks (off + 2 ^ k)

/-- The offset of every column, in column-index order (`offsets[i] = off` for `i` in order). -/
def stackOffsets (kappas : List ℕ) : List ℕ :=
  let order := stackOrder kappas
  let offs := prefixOffsets (order.map (kappas[·]!)) 0
  (List.range kappas.length).map fun i ↦ (offs[order.idxOf i]?).getD 0

/-- A `Blocks` from a list of sizes that is already largest first. -/
def blocksOfList (l : List ℕ) (h : l.Pairwise (· ≥ ·)) : Blocks where
  n := l.length
  size := fun i ↦ l[i]
  descending := by
    intro a b hab
    rcases eq_or_lt_of_le hab with rfl | hlt
    · exact le_rfl
    · exact (List.pairwise_iff_getElem.mp h) a.val b.val a.isLt b.isLt hlt

/-- Every offset of a `Blocks`, first block first. -/
def offsetsOf (B : Blocks) : List ℕ := (List.finRange B.n).map B.offset

/-- Every selector of a `Blocks`, first block first. -/
def selectorsOf (B : Blocks) {μ : ℕ} (hμ : B.total ≤ 2 ^ μ) : List ℕ :=
  (List.finRange B.n).map fun b ↦ (B.selector hμ b).val

/-! ## The small configuration

Column-index order: `mem_0, mem_1, mem_2, cntfin_mem, cntfin_bc, q_flock, table a, table b`. -/

def small : List ℕ := [4, 4, 4, 4, 2, 5, 4, 4]

-- The transcription gives what the pinned Python verifier and the Rust give (offsets.py).
#guard stackOrder small = [5, 0, 1, 2, 3, 6, 7, 4]
#guard stackOffsets small = [32, 48, 64, 80, 128, 0, 96, 112]

/-- The sorted sizes, as a `Blocks`. -/
def smallBlocks : Blocks := blocksOfList [5, 4, 4, 4, 4, 4, 4, 2] (by decide)

#guard (stackOrder small).map (small[·]!) = [5, 4, 4, 4, 4, 4, 4, 2]

theorem smallBlocks_total_le : smallBlocks.total ≤ 2 ^ 8 := by decide

example : smallBlocks.total = 132 := by decide

-- `Blocks.offset` and `Blocks.selector`, block by block, against the pinned verifiers' values.
#guard offsetsOf smallBlocks = [0, 32, 48, 64, 80, 96, 112, 128]
#guard selectorsOf smallBlocks smallBlocks_total_le = [0, 2, 3, 4, 5, 6, 7, 32]

-- Column `i` sits on block `(stackOrder small).idxOf i`; through that map the offsets are
-- leanVM's, column by column.
#guard (List.range 8).map (fun i ↦ (offsetsOf smallBlocks)[(stackOrder small).idxOf i]!) =
  stackOffsets small

/-- The order the blueprint had before finding F17: ties broken with the tables' columns
first. It sorts to the same sizes, so it is the same `Blocks`. -/
def tablesFirst : List ℕ := [5, 6, 7, 0, 1, 2, 3, 4]

#guard tablesFirst.map (small[·]!) = (stackOrder small).map (small[·]!)

-- Near miss: through that map `mem_0` sits at offset 64, not 32, and nothing about the
-- `Blocks` differs: the order of equal sizes lives in the map from columns to blocks only.
#guard (List.range 8).map (fun i ↦ (offsetsOf smallBlocks)[tablesFirst.idxOf i]!) =
  [64, 80, 96, 112, 128, 0, 32, 48]
#guard (List.range 8).map (fun i ↦ (offsetsOf smallBlocks)[tablesFirst.idxOf i]!) ≠
  stackOffsets small

-- The column-index order itself is not largest first, so it is no `Blocks`.
example : ¬ small.Pairwise (· ≥ ·) := by decide

/-! ## The pinned verifier's layout for one admissible announcement

`build_layout(log_memory = 16, taus = (3, 16, 0, 5, 16, 3), kbc = 4)`: 110 columns, 92 of
them committed. The three lists are copied from the output of `offsets.py`. -/

/-- The sizes of the 92 committed columns, in column-index order (the eighteen BLAKE2S limb
columns, global indices 82 to 99, left out). -/
def committedSizes : List ℕ :=
  [16, 16, 16, 16, 4, 11] ++ List.replicate 15 3 ++ List.replicate 15 16 ++
    List.replicate 8 0 ++ List.replicate 15 5 ++ List.replicate 14 16 ++
    List.replicate 9 3 ++ List.replicate 10 3

/-- The global column index of each committed column, in the same order. -/
def committedIndex : List ℕ := List.range 82 ++ (List.range 10).map (· + 100)

/-- Python: "column index of each block". -/
def pyOrder : List ℕ :=
  [0, 1, 2, 3, 21, 22, 23, 24, 25, 26, 27, 28, 29, 30, 31, 32, 33, 34, 35, 59, 60, 61, 62, 63,
   64, 65, 66, 67, 68, 69, 70, 71, 72, 5, 44, 45, 46, 47, 48, 49, 50, 51, 52, 53, 54, 55, 56,
   57, 58, 4, 6, 7, 8, 9, 10, 11, 12, 13, 14, 15, 16, 17, 18, 19, 20, 73, 74, 75, 76, 77, 78,
   79, 80, 81, 100, 101, 102, 103, 104, 105, 106, 107, 108, 109, 36, 37, 38, 39, 40, 41, 42, 43]

/-- Python: "offsets", first block first. -/
def pyOffsets : List ℕ :=
  [0, 65536, 131072, 196608, 262144, 327680, 393216, 458752, 524288, 589824, 655360, 720896,
   786432, 851968, 917504, 983040, 1048576, 1114112, 1179648, 1245184, 1310720, 1376256,
   1441792, 1507328, 1572864, 1638400, 1703936, 1769472, 1835008, 1900544, 1966080, 2031616,
   2097152, 2162688, 2164736, 2164768, 2164800, 2164832, 2164864, 2164896, 2164928, 2164960,
   2164992, 2165024, 2165056, 2165088, 2165120, 2165152, 2165184, 2165216, 2165232, 2165240,
   2165248, 2165256, 2165264, 2165272, 2165280, 2165288, 2165296, 2165304, 2165312, 2165320,
   2165328, 2165336, 2165344, 2165352, 2165360, 2165368, 2165376, 2165384, 2165392, 2165400,
   2165408, 2165416, 2165424, 2165432, 2165440, 2165448, 2165456, 2165464, 2165472, 2165480,
   2165488, 2165496, 2165504, 2165505, 2165506, 2165507, 2165508, 2165509, 2165510, 2165511]

/-- The sorted sizes. -/
def sortedSizes : List ℕ :=
  List.replicate 33 16 ++ [11] ++ List.replicate 15 5 ++ [4] ++ List.replicate 34 3 ++
    List.replicate 8 0

#guard committedSizes.length = 92
#guard (stackOrder committedSizes).map (committedIndex[·]!) = pyOrder
#guard (stackOrder committedSizes).map (committedSizes[·]!) = sortedSizes

/-- The sorted sizes as a `Blocks`. -/
def realBlocks : Blocks := blocksOfList sortedSizes (by decide +kernel)

#guard offsetsOf realBlocks = pyOffsets
#guard realBlocks.total = 2165512
example : realBlocks.total ≤ 2 ^ 22 := by decide +kernel
example : ¬ realBlocks.total ≤ 2 ^ 21 := by decide +kernel

end Probe
\end{leancode}

\textbf{Recorded output.}

\begin{srccode}
(no output)
exit=0
\end{srccode}

\subsubsection{\texttt{values.\allowbreak{}py} (I.3)}\label{sec:gen-layer1-I3}

\fhead{What it establishes}Reference numbers from the pinned leanVM, to be compared with Layer 1's Lean definitions in \texttt{ValuesProbe.\allowbreak{}lean}. Every evaluation is made by a function of the pinned Python verifier (\texttt{multilinear\_\allowbreak{}eval}, \texttt{index\_\allowbreak{}mle}, \texttt{eq\_\allowbreak{}eval}, \texttt{stack\_\allowbreak{}offsets}, \texttt{Placement}), imported read-only; the bytecode table is built by a line-by-line transcription of the Rust \texttt{bytecode\_\allowbreak{}columns} and \texttt{stacked\_\allowbreak{}bytecode\_\allowbreak{}table}, since the Python verifier takes that table as an input file and has no encoder of its own.

\textbf{Command.} \texttt{p\allowbreak{}y\allowbreak{}t\allowbreak{}h\allowbreak{}o\allowbreak{}n\allowbreak{}3\allowbreak{} -{}\allowbreak{}B\allowbreak{} .\allowbreak{}c\allowbreak{}l\allowbreak{}a\allowbreak{}u\allowbreak{}d\allowbreak{}e\allowbreak{}/\allowbreak{}r\allowbreak{}e\allowbreak{}p\allowbreak{}o\allowbreak{}r\allowbreak{}t\allowbreak{}s\allowbreak{}/\allowbreak{}b\allowbreak{}l\allowbreak{}u\allowbreak{}e\allowbreak{}p\allowbreak{}r\allowbreak{}i\allowbreak{}n\allowbreak{}t\allowbreak{}-{}\allowbreak{}r\allowbreak{}e\allowbreak{}v\allowbreak{}i\allowbreak{}e\allowbreak{}w\allowbreak{}/\allowbreak{}p\allowbreak{}r\allowbreak{}o\allowbreak{}b\allowbreak{}e\allowbreak{}s\allowbreak{}/\allowbreak{}c\allowbreak{}o\allowbreak{}d\allowbreak{}e\allowbreak{}-{}\allowbreak{}l\allowbreak{}a\allowbreak{}y\allowbreak{}e\allowbreak{}r\allowbreak{}1\allowbreak{}/\allowbreak{}v\allowbreak{}a\allowbreak{}l\allowbreak{}u\allowbreak{}e\allowbreak{}s\allowbreak{}.\allowbreak{}p\allowbreak{}y}

\begin{srccode}
"""Probe for task code-layer1, deliverables B.3, B.5, B.6, B.7 and C.

Reference numbers from the pinned leanVM (a386121f), to be compared with Layer 1's Lean
definitions in `ValuesProbe.lean`. Every evaluation below is made by a function of the pinned
Python verifier (`multilinear_eval`, `index_mle`, `eq_eval`, `stack_offsets`, `Placement`),
imported read-only. The bytecode table is built by a line-by-line transcription of the Rust
`bytecode_columns` (crates/lean_vm/src/cpu/layout.rs:229-309) and `stacked_bytecode_table`
(crates/lean_vm/src/leaf.rs:585-604): the Python verifier takes that table as an input file and
has no encoder of its own.

Run:  python3 -B .claude/reports/blueprint-review/probes/code-layer1/values.py
"""

import sys

sys.dont_write_bytecode = True
sys.path.insert(0, "/home/scaraven/Documents/leanEthereum/leanVM/python-verifier")
import verifier as V  # noqa: E402

K, E = V.K, V.E


def limbs(x):
    return (x.c0.value, x.c1.value, x.c2.value)


def gk(i):
    """g^i in K (primitives/field/mod.rs:82, g_pow)."""
    r = K(1)
    for _ in range(i):
        r = r * K(2)
    return r


# ---- the bytecode table ---------------------------------------------------------------
OP = {"xor": gk(0), "mul": gk(1), "set": gk(2), "deref": gk(3), "jump": gk(4), "blake2s": gk(5)}  # tables.rs:95-100


def bytecode_columns(prog):
    """layout.rs:229-309: (opcode, o1, o2, o3, fpc, ffp, extra0, extra1) per instruction."""
    cols = [[] for _ in range(8)]
    for op in prog:
        kind = op[0]
        z = K(0)
        if kind in ("xor", "mul"):
            _, a, b, c = op
            row = [OP[kind], gk(a), gk(b), gk(c), z, z, z, z]
        elif kind == "set":
            _, o, k = op  # k = (c0, c1, c2)
            row = [OP[kind], gk(o), K(k[0]), K(k[1]), K(k[2]), z, z, z]  # :257, :270
        elif kind == "deref":
            _, o1, o2, o3, mode = op
            f_pc = K(1) if mode == "pc" else z  # isa.rs:69-74
            f_fp = K(1) if mode == "fp" else z
            row = [OP[kind], gk(o1), gk(o2), gk(o3), f_pc, f_fp, z, z]
        elif kind == "jump":
            _, oc, od, of = op
            row = [OP[kind], gk(oc), gk(od), gk(of), z, z, z, z]
        elif kind == "blake2s":
            _, ins, cv, out, md = op
            row = [OP[kind], gk(ins[0]), gk(ins[1]), gk(ins[2]), gk(ins[3]), gk(cv), gk(out), gk(md)]
        for c in range(8):
            cols[c].append(row[c])
    return cols


def stacked_bytecode_table(cols, kbc):
    """leaf.rs:585-604: column i at slot 3 + i, `table[(slot << kbc)..((slot+1) << kbc)]`."""
    table = [K(0)] * (1 << (4 + kbc))
    for i, vals in enumerate(cols):
        slot = 3 + i
        assert len(vals) == 1 << kbc
        table[(slot << kbc):((slot + 1) << kbc)] = vals
    return table


PAD = ("set", 0, (0, 0, 0))  # lean_compiler/src/lib.rs:162
prog = [
    ("xor", 1, 2, 3),
    ("mul", 4, 5, 6),
    ("set", 7, (11, 12, 13)),
    ("deref", 8, 9, 10, "pc"),
    ("deref", 21, 22, 23, "fp"),
    ("deref", 24, 25, 26, "cell"),
    ("jump", 27, 28, 29),
    ("blake2s", (14, 15, 16, 17), 18, 19, 20),
] + [PAD] * 8
kbc = 4
table = stacked_bytecode_table(bytecode_columns(prog), kbc)
print("== bytecode table, 16 instructions, 256 cells (cell i + 16*s holds slot s of instruction i)")
print([w.value for w in table])

zeta = (E(3, 1, 0), E(5, 0, 7), E(2, 2, 2), E(0, 1, 0))
alpha = (E(9, 0, 1), E(0, 4, 0), E(6, 6, 0), E(1, 2, 3))
print("== bytecode multilinear at (zeta, alpha), as verifier.py:566 evaluates it")
print(limbs(V.multilinear_eval(table, (*zeta, *alpha))))
print("== with alpha reversed")
print(limbs(V.multilinear_eval(table, (*zeta, *reversed(alpha)))))

# ---- the index column ----------------------------------------------------------------
print("== index_mle at zeta (4 variables) and at its first 2 coordinates")
print(limbs(V.index_mle(zeta)), limbs(V.index_mle(zeta[:2])))
print("== the MLE of [g^0 .. g^15] at zeta, by multilinear_eval")
print(limbs(V.multilinear_eval([gk(i) for i in range(16)], zeta)))

# ---- stacking: selectors, the zero-padded and the one-padded stack ----------------------
sizes = [2, 1, 0]
tables = [[1, 2, 3, 4], [5, 6], [7]]
offsets, mu = V.stack_offsets(sizes)
print("== blocks of sizes 2, 1, 0: offsets", offsets, "mu", mu)
z3 = zeta[:3]
placements = [V.Placement(k, off) for k, off in zip(sizes, offsets)]
sel = [p.eq_above(z3) for p in placements]
print("eq(sel_b, zeta_hi) per block:", [limbs(s) for s in sel])
blocks_at = [V.multilinear_eval([K(v) for v in t], z3[:k]) for t, k in zip(tables, sizes)]
print("block b at zeta_lo:", [limbs(b) for b in blocks_at])
covered = E.sum(s * b for s, b in zip(sel, blocks_at))
ones_padding = E.sum(sel) + V.ONE  # verifier.py:589
print("covered part:", limbs(covered), " ones padding:", limbs(ones_padding))
stack0 = [K(v) for v in [1, 2, 3, 4, 5, 6, 7, 0]]
stack1 = [K(v) for v in [1, 2, 3, 4, 5, 6, 7, 1]]
print("zero-padded stack at zeta:", limbs(V.multilinear_eval(stack0, z3)))
print("one-padded stack at zeta :", limbs(V.multilinear_eval(stack1, z3)))
assert V.multilinear_eval(stack0, z3) == covered
assert V.multilinear_eval(stack1, z3) == covered + ones_padding

# ---- a column claim as a weight on the stack (verifier.py:520-522) -------------------------
point = (E(3, 1, 0),)  # a claim on block 1 (size 1)
lifted = placements[1].stack_point(point, mu)
print("== claim on block 1 at", [limbs(p) for p in point], "lifts to", [limbs(p) for p in lifted])
kernel = V.eq_kernel(lifted)
print("weight eq(lifted, .) on the cube paired with the arbitrary stack [9,8,7,6,5,4,3,2]:")
arbitrary = [K(v) for v in [9, 8, 7, 6, 5, 4, 3, 2]]
print(limbs(V.dot(arbitrary, kernel)), "  block 1 of it, [5,4], at the point:", limbs(V.multilinear_eval([K(5), K(4)], point)))
r = (E(1, 1, 0), E(0, 2, 0), E(7, 0, 0))
print("the weight's extension at r, eq_eval(lifted, r):", limbs(V.eq_eval(lifted, r)))

# ---- back-loaded padding (verifier.py:609-614) -------------------------------------------
# a table on 1 variable in a sumcheck on 3: the summand's weight at the final point carries
# `challenge` for each variable the table lacks.
print("== back-loaded padding: [3,5] lifted by two variables, at (7 | 11, 13) over K-embedded points")
pt = (E(7), E(11), E(13))
short = V.multilinear_eval([K(3), K(5)], pt[:1])
print("table at 7 times 11*13:", limbs(short * pt[1] * pt[2]))
print("MLE of [0,0,0,0,0,0,3,5] at the point:", limbs(V.multilinear_eval([K(v) for v in [0, 0, 0, 0, 0, 0, 3, 5]], pt)))
\end{srccode}

\textbf{Recorded output.}

\begin{srccode}
== bytecode table, 16 instructions, 256 cells (cell i + 16*s holds slot s of instruction i)
[0, 0, 0, 0, 0, 0, 0, 0, 0, 0, 0, 0, 0, 0, 0, 0, 0, 0, 0, 0, 0, 0, 0, 0, 0, 0, 0, 0, 0, 0, 0, 0, 0, 0, 0, 0, 0, 0, 0, 0, 0, 0, 0, 0, 0, 0, 0, 0, 1, 2, 4, 8, 8, 8, 16, 32, 4, 4, 4, 4, 4, 4, 4, 4, 2, 16, 128, 256, 2097152, 16777216, 134217728, 16384, 1, 1, 1, 1, 1, 1, 1, 1, 4, 32, 11, 512, 4194304, 33554432, 268435456, 32768, 0, 0, 0, 0, 0, 0, 0, 0, 8, 64, 12, 1024, 8388608, 67108864, 536870912, 65536, 0, 0, 0, 0, 0, 0, 0, 0, 0, 0, 13, 1, 0, 0, 0, 131072, 0, 0, 0, 0, 0, 0, 0, 0, 0, 0, 0, 0, 1, 0, 0, 262144, 0, 0, 0, 0, 0, 0, 0, 0, 0, 0, 0, 0, 0, 0, 0, 524288, 0, 0, 0, 0, 0, 0, 0, 0, 0, 0, 0, 0, 0, 0, 0, 1048576, 0, 0, 0, 0, 0, 0, 0, 0, 0, 0, 0, 0, 0, 0, 0, 0, 0, 0, 0, 0, 0, 0, 0, 0, 0, 0, 0, 0, 0, 0, 0, 0, 0, 0, 0, 0, 0, 0, 0, 0, 0, 0, 0, 0, 0, 0, 0, 0, 0, 0, 0, 0, 0, 0, 0, 0, 0, 0, 0, 0, 0, 0, 0, 0, 0, 0, 0, 0, 0, 0, 0, 0, 0, 0, 0, 0, 0, 0, 0, 0, 0, 0, 0, 0, 0, 0, 0, 0]
== bytecode multilinear at (zeta, alpha), as verifier.py:566 evaluates it
(1219889995492, 3571792677380, 1279835512890)
== with alpha reversed
(1332010270628, 4168742050444, 1982034182034)
== index_mle at zeta (4 variables) and at its first 2 coordinates
(950617, 874814, 877803) (109, 29, 108)
== the MLE of [g^0 .. g^15] at zeta, by multilinear_eval
(950617, 874814, 877803)
== blocks of sizes 2, 1, 0: offsets [0, 4, 6] mu 3
eq(sel_b, zeta_hi) per block: [(3, 2, 2), (6, 8, 8), (2, 26, 30)]
block b at zeta_lo: [(46, 11, 42), (0, 3, 0), (7, 0, 0)]
covered part: (38, 3, 34)  ones padding: (6, 16, 20)
zero-padded stack at zeta: (38, 3, 34)
one-padded stack at zeta : (32, 19, 54)
== claim on block 1 at [(3, 1, 0)] lifts to [(3, 1, 0), (0, 0, 0), (1, 0, 0)]
weight eq(lifted, .) on the cube paired with the arbitrary stack [9,8,7,6,5,4,3,2]:
(6, 1, 0)   block 1 of it, [5,4], at the point: (6, 1, 0)
the weight's extension at r, eq_eval(lifted, r): (9, 18, 0)
== back-loaded padding: [3,5] lifted by two variables, at (7 | 11, 13) over K-embedded points
table at 7 times 11*13: (1935, 0, 0)
MLE of [0,0,0,0,0,0,3,5] at the point: (1935, 0, 0)
\end{srccode}

\subsubsection{\texttt{ValuesProbe.\allowbreak{}lean} (I.4)}\label{sec:gen-layer1-I4}

\fhead{What it establishes}The index column, the bytecode column (256 cells and the \texttt{(\allowbreak{}ζ,\allowbreak{} α)} evaluation), the selector weights, both paddings, the column-claim weight and back-loaded padding agree with the pinned Python's numbers, with their near misses; \texttt{bytecodeColumn\_\allowbreak{}slot}'s proof also proves the opposite layout's analogue; \texttt{Layout.\allowbreak{}comap} aliases; and the axioms of seventeen Layer 1 declarations and of the probe's \texttt{bytecodeColumn\_\allowbreak{}answer\_\allowbreak{}slots} (B.3-B.7, C.3, C.4).

\textbf{Command.} \texttt{f\allowbreak{}l\allowbreak{}o\allowbreak{}c\allowbreak{}k\allowbreak{} .\allowbreak{}c\allowbreak{}l\allowbreak{}a\allowbreak{}u\allowbreak{}d\allowbreak{}e\allowbreak{}/\allowbreak{}r\allowbreak{}e\allowbreak{}p\allowbreak{}o\allowbreak{}r\allowbreak{}t\allowbreak{}s\allowbreak{}/\allowbreak{}b\allowbreak{}l\allowbreak{}u\allowbreak{}e\allowbreak{}p\allowbreak{}r\allowbreak{}i\allowbreak{}n\allowbreak{}t\allowbreak{}-{}\allowbreak{}r\allowbreak{}e\allowbreak{}v\allowbreak{}i\allowbreak{}e\allowbreak{}w\allowbreak{}/\allowbreak{}l\allowbreak{}o\allowbreak{}g\allowbreak{}s\allowbreak{}/\allowbreak{}l\allowbreak{}e\allowbreak{}a\allowbreak{}n\allowbreak{}.\allowbreak{}l\allowbreak{}o\allowbreak{}c\allowbreak{}k\allowbreak{} l\allowbreak{}a\allowbreak{}k\allowbreak{}e\allowbreak{} e\allowbreak{}n\allowbreak{}v\allowbreak{} l\allowbreak{}e\allowbreak{}a\allowbreak{}n\allowbreak{} .\allowbreak{}c\allowbreak{}l\allowbreak{}a\allowbreak{}u\allowbreak{}d\allowbreak{}e\allowbreak{}/\allowbreak{}r\allowbreak{}e\allowbreak{}p\allowbreak{}o\allowbreak{}r\allowbreak{}t\allowbreak{}s\allowbreak{}/\allowbreak{}b\allowbreak{}l\allowbreak{}u\allowbreak{}e\allowbreak{}p\allowbreak{}r\allowbreak{}i\allowbreak{}n\allowbreak{}t\allowbreak{}-{}\allowbreak{}r\allowbreak{}e\allowbreak{}v\allowbreak{}i\allowbreak{}e\allowbreak{}w\allowbreak{}/\allowbreak{}p\allowbreak{}r\allowbreak{}o\allowbreak{}b\allowbreak{}e\allowbreak{}s\allowbreak{}/\allowbreak{}c\allowbreak{}o\allowbreak{}d\allowbreak{}e\allowbreak{}-{}\allowbreak{}l\allowbreak{}a\allowbreak{}y\allowbreak{}e\allowbreak{}r\allowbreak{}1\allowbreak{}/\allowbreak{}V\allowbreak{}a\allowbreak{}l\allowbreak{}u\allowbreak{}e\allowbreak{}s\allowbreak{}P\allowbreak{}r\allowbreak{}o\allowbreak{}b\allowbreak{}e\allowbreak{}.\allowbreak{}l\allowbreak{}e\allowbreak{}a\allowbreak{}n}

\begin{leancode}
/-
  Probe for task code-layer1, deliverables B.3, B.5, B.6, B.7 and C.

  Layer 1's definitions evaluated against numbers produced by the pinned leanVM: every
  expected value below is copied from the output of `values.py`, which calls the pinned Python
  verifier's own functions (and, for the bytecode table, a transcription of the Rust encoder).

  Run from the repository root:
    flock .claude/reports/blueprint-review/logs/lean.lock \
      lake env lean .claude/reports/blueprint-review/probes/code-layer1/ValuesProbe.lean
-/
import LeanerVM.Protocol.FixedColumns
import LeanerVM.Protocol.Stack
import LeanerVM.Protocol.Padding
import LeanerVM.Protocol.ClaimWeights
import LeanerVM.Protocol.BlockClaims

open LeanerVM.Protocol LeanerVM.Parameters LeanerVM.Semantics LeanerVM.Arithmetization
open CompPoly CMlPolynomialEval

namespace Probe

/-- The oracle's answer, as a function `#guard` can evaluate. -/
def answer {n : ℕ} (q : Column n) (z : Vector E n) : E := OracleInterface.answer q z

/-- The point `ζ` of `values.py`. -/
def zeta : Vector E 4 := #v[E.ofLimbs 3 1 0, E.ofLimbs 5 0 7, E.ofLimbs 2 2 2, E.ofLimbs 0 1 0]

/-- The point `α` of `values.py`. -/
def alpha : Vector E 4 := #v[E.ofLimbs 9 0 1, E.ofLimbs 0 4 0, E.ofLimbs 6 6 0, E.ofLimbs 1 2 3]

/-- `α` with its coordinates reversed. -/
def alphaRev : Vector E 4 := #v[E.ofLimbs 1 2 3, E.ofLimbs 6 6 0, E.ofLimbs 0 4 0, E.ofLimbs 9 0 1]

/-! ## B.6 The bytecode column -/

/-- Sixteen instructions: the six opcodes, the three `DEREF` modes, distinct operands, and the
compiler's padding instruction `SET_CONSTANT [g^0] 0` in the last eight cells, the last index
(the halting address) included. Operand `a` of the Rust is the field element `g^a` here. -/
def prog16 : Program where
  logSize := 4
  logSize_le := by decide
  code i := match i.val with
    | 0 => .xor (gpow 1) (gpow 2) (gpow 3)
    | 1 => .mulNative (gpow 4) (gpow 5) (gpow 6)
    | 2 => .setConstant (gpow 7) (E.ofLimbs 11 12 13)
    | 3 => .deref (gpow 8) (gpow 9) (gpow 10) .pc
    | 4 => .deref (gpow 21) (gpow 22) (gpow 23) .fp
    | 5 => .deref (gpow 24) (gpow 25) (gpow 26) .cell
    | 6 => .jump (gpow 27) (gpow 28) (gpow 29)
    | 7 => .blake2s ![gpow 14, gpow 15, gpow 16, gpow 17] (gpow 18) (gpow 19) (gpow 20)
    | _ => .setConstant (gpow 0) 0

/-- The table the transcription of `bytecode_columns` and `stacked_bytecode_table` gives. -/
def pyTable : List ℕ :=
  [0, 0, 0, 0, 0, 0, 0, 0, 0, 0, 0, 0, 0, 0, 0, 0, 0, 0, 0, 0, 0, 0, 0, 0, 0, 0, 0, 0, 0, 0, 0, 0, 0, 0, 0, 0, 0, 0, 0, 0, 0, 0, 0, 0, 0, 0, 0, 0, 1, 2, 4, 8, 8, 8, 16, 32, 4, 4, 4, 4, 4, 4, 4, 4, 2, 16, 128, 256, 2097152, 16777216, 134217728, 16384, 1, 1, 1, 1, 1, 1, 1, 1, 4, 32, 11, 512, 4194304, 33554432, 268435456, 32768, 0, 0, 0, 0, 0, 0, 0, 0, 8, 64, 12, 1024, 8388608, 67108864, 536870912, 65536, 0, 0, 0, 0, 0, 0, 0, 0, 0, 0, 13, 1, 0, 0, 0, 131072, 0, 0, 0, 0, 0, 0, 0, 0, 0, 0, 0, 0, 1, 0, 0, 262144, 0, 0, 0, 0, 0, 0, 0, 0, 0, 0, 0, 0, 0, 0, 0, 524288, 0, 0, 0, 0, 0, 0, 0, 0, 0, 0, 0, 0, 0, 0, 0, 1048576, 0, 0, 0, 0, 0, 0, 0, 0, 0, 0, 0, 0, 0, 0, 0, 0, 0, 0, 0, 0, 0, 0, 0, 0, 0, 0, 0, 0, 0, 0, 0, 0, 0, 0, 0, 0, 0, 0, 0, 0, 0, 0, 0, 0, 0, 0, 0, 0, 0, 0, 0, 0, 0, 0, 0, 0, 0, 0, 0, 0, 0, 0, 0, 0, 0, 0, 0, 0, 0, 0, 0, 0, 0, 0, 0, 0, 0, 0, 0, 0, 0, 0, 0, 0, 0, 0, 0, 0]

#guard pyTable.length = 256
-- All 256 cells: the encoding, the opcode values, the slot order and the cell order.
#guard (bytecodeColumn prog16).values.toList.map (·.toNat) = pyTable

/-- The oracle's answer at `(ζ, α)`. -/
def bcAnswer : E := answer (bytecodeColumn prog16) (zeta ++ alpha)

-- The evaluation the pinned verifier makes (`verifier.py:566`), and the native evaluator.
#guard bcAnswer = E.ofLimbs 1219889995492 3571792677380 1279835512890
#guard bytecodeColumnEval prog16 zeta alpha = E.ofLimbs 1219889995492 3571792677380 1279835512890
-- Near miss: with `α` reversed the pinned verifier gets another value, and so does the oracle.
#guard answer (bytecodeColumn prog16) (zeta ++ alphaRev) =
  E.ofLimbs 1332010270628 4168742050444 1982034182034
#guard bcAnswer ≠ E.ofLimbs 1332010270628 4168742050444 1982034182034

/-- The opposite layout: slot bits low, instruction bits high. -/
def wrongColumn (prog : Program) : Column (4 + prog.logSize) :=
  ⟨Vector.ofFn fun i ↦
    let p := (cubeSplit 4 prog.logSize).symm i
    (encodeSlots (prog.code p.2))[p.1]⟩

/-- The analogue of `bytecodeColumn_slot` for the opposite layout, by the same proof script:
the *shape* of that theorem holds of either layout. -/
theorem wrongColumn_slot (prog : Program) (i : Fin (2 ^ prog.logSize)) (s : Fin 16) :
    (wrongColumn prog).values[cubeIndex (k := 4) s i] = (encodeSlots (prog.code i))[s] := by
  simp [wrongColumn, ← cubeSplit_apply]

-- The statement of `bytecodeColumn_slot` itself, word for word, is false of the opposite
-- layout: cell `0 + 16 * 3` is instruction 0's opcode there, and slot 0 of instruction 3 here.
#guard (bytecodeColumn prog16).values.toList[48]! = Opcode.xor.code
#guard (wrongColumn prog16).values.toList[48]! = 0
#guard (wrongColumn prog16).values.toList.map (·.toNat) ≠ pyTable

/-- What the bus phase needs and Layer 1 does not state: the program's share of a bytecode
block is one evaluation of the bytecode column, the slots weighted by `eq(w, ·)`. -/
theorem bytecodeColumn_answer_slots (prog : Program) (z : Vector E prog.logSize)
    (w : Vector E 4) :
    eval₂Mle (bytecodeColumn prog).values (algebraMap K E) (z ++ w) =
      ∑ s : Fin 16, (lagrangeBasis w)[s] *
        eval₂Mle (bytecodeSlotColumn prog s).values (algebraMap K E) z := by
  rw [eval₂Mle, evalMle_split]
  refine Finset.sum_congr rfl fun s _ ↦ ?_
  congr 1
  rw [eval₂Mle, ← slice_bytecodeColumn]
  congr 1
  apply Vector.ext
  intro i hi
  simp [slice, CMlPolynomialEval.map]

/-! ## B.5 The index column -/

#guard answer (idxColumn 4) zeta = E.ofLimbs 950617 874814 877803
#guard idxColumnEval zeta = E.ofLimbs 950617 874814 877803
#guard idxColumnEval (#v[E.ofLimbs 3 1 0, E.ofLimbs 5 0 7] : Vector E 2) = E.ofLimbs 109 29 108
#guard (idxColumn 4).values.toList.map (·.toNat) = (List.range 16).map (2 ^ ·)

/-! ## B.2, B.3 Selectors and the two paddings -/

/-- Blocks on 2, 1 and 0 variables. -/
def blocks : Blocks where
  n := 3
  size := ![2, 1, 0]
  descending := by
    show ∀ a b : Fin 3, a ≤ b → ![2, 1, 0] b ≤ ![2, 1, 0] a
    decide

theorem blocks_total_le : blocks.total ≤ 2 ^ 3 := by decide

/-- The tables `[1, 2, 3, 4]`, `[5, 6]`, `[7]`. -/
def tables : blocks.Tables K :=
  show (b : Fin 3) → CMlPolynomialEval K (![2, 1, 0] b) from fun b ↦ match b with
    | 0 => (#v[1, 2, 3, 4] : CMlPolynomialEval K 2)
    | 1 => (#v[5, 6] : CMlPolynomialEval K 1)
    | 2 => (#v[7] : CMlPolynomialEval K 0)

/-- The tables over `E`. -/
def tablesE : blocks.Tables E := fun b ↦ CMlPolynomialEval.map (algebraMap K E) (tables b)

/-- The first three coordinates of `ζ`. -/
def z3 : Vector E 3 := #v[E.ofLimbs 3 1 0, E.ofLimbs 5 0 7, E.ofLimbs 2 2 2]

/-- The selector weights, block by block. -/
def weights : List E := (List.finRange 3).map fun b ↦ blocks.selectorWeight blocks_total_le b z3

/-- The blocks at the low coordinates, block by block. -/
def blocksAt : List E := (List.finRange 3).map fun b ↦
  eval₂Mle (tables b) (algebraMap K E) (blocks.lowPoint blocks_total_le b z3)

-- `eq(sel_b, ζ_hi)` as `Placement.eq_above` computes it (`verifier.py:299-302`).
#guard weights = [E.ofLimbs 3 2 2, E.ofLimbs 6 8 8, E.ofLimbs 2 26 30]
#guard blocksAt = [E.ofLimbs 46 11 42, E.ofLimbs 0 3 0, E.ofLimbs 7 0 0]
-- The witness stack (pad 0) and a leaf stack (pad 1) at `ζ`.
#guard eval₂Mle (blocks.stackColumn tables 3).values (algebraMap K E) z3 = E.ofLimbs 38 3 34
#guard evalMle (blocks.stackAt tablesE 3 1) z3 = E.ofLimbs 32 19 54
-- The padding term of equation (2), `1 + Σ_b eq(sel_b, ζ_hi)` (`verifier.py:589`).
#guard 1 + weights.sum = E.ofLimbs 6 16 20
#guard (List.zipWith (· * ·) weights blocksAt).sum + (1 + weights.sum) = E.ofLimbs 32 19 54
-- Near misses: the wrong pad, and the padding term read as the constant 1.
#guard evalMle (blocks.stackAt tablesE 3 0) z3 ≠ E.ofLimbs 32 19 54
#guard (List.zipWith (· * ·) weights blocksAt).sum + 1 ≠ E.ofLimbs 32 19 54

/-! ## B.7 A column claim as a weight -/

/-- The claim's point on block 1, lifted. -/
def lifted : Vector E 3 := blocks.extendPoint blocks_total_le (1 : Fin 3) #v[E.ofLimbs 3 1 0]

/-- A stack no honest prover commits. -/
def arbitrary : Column 3 := ⟨#v[9, 8, 7, 6, 5, 4, 3, 2]⟩

#guard lifted = #v[E.ofLimbs 3 1 0, 0, 1]
#guard (eqWeight lifted).pair arbitrary = E.ofLimbs 6 1 0
#guard (eqWeight lifted).mle #v[E.ofLimbs 1 1 0, E.ofLimbs 0 2 0, E.ofLimbs 7 0 0] =
  E.ofLimbs 9 18 0
-- Near miss: the selector bits reversed, `(1, 0)`, weigh cells 2 and 3.
#guard (eqWeight (#v[E.ofLimbs 3 1 0, 1, 0] : Vector E 3)).pair arbitrary ≠ E.ofLimbs 6 1 0

/-! ## B.4 Back-loaded padding -/

/-- The table `[3, 5]` on one variable. -/
def short : CMlPolynomialEval K 1 := #v[3, 5]

#guard evalMle (padHigh short 2) ((#v[7] : Vector K 1) ++ (#v[11, 13] : Vector K 2)) = 1935
#guard (padHigh short 2).toList = [0, 0, 0, 0, 0, 0, 3, 5]
-- Near miss: padding on the low variables (front-loaded) is another table.
#guard (padHigh short 2).toList ≠ [0, 0, 0, 3, 0, 0, 0, 5]

/-! ## C A layout need not separate its columns -/

/-- Three columns all read off block 0: `Layout.comap` asks nothing of the renaming. -/
def aliased : Layout 3 (Fin 3) (fun _ ↦ 2) :=
  (blocks.layout blocks_total_le).comap (fun _ ↦ (0 : Fin 3)) (fun _ ↦ rfl)

#guard (List.finRange 3).all fun c ↦ (aliased.read arbitrary c).values.toList = [9, 8, 7, 6]

/-! ## Axioms -/

#print axioms Blocks.unstack_eval₂
#print axioms Blocks.unstack_getElem
#print axioms Blocks.readColumn_eval
#print axioms Blocks.layout
#print axioms Layout.comap
#print axioms Blocks.stack_eval_ambient
#print axioms Blocks.stack_eval_ambient_one
#print axioms Blocks.stackColumn_eval_ambient
#print axioms sumCube_padHigh
#print axioms evalMle_padHigh
#print axioms ColumnClaim.holds_iff_weighted
#print axioms eqWeight
#print axioms idxColumn_eval
#print axioms idxColumnEval_eq
#print axioms bytecodeColumn_answer_boolVec
#print axioms bytecodeColumn_eval
#print axioms BlockClaim.isValid_iff_pairing
#print axioms bytecodeColumn_answer_slots

end Probe
\end{leancode}

\textbf{Recorded output.}

\begin{srccode}
'LeanerVM.Protocol.Blocks.unstack_eval₂' depends on axioms: [propext, Classical.choice, Quot.sound]
'LeanerVM.Protocol.Blocks.unstack_getElem' depends on axioms: [propext, Classical.choice, Quot.sound]
'LeanerVM.Protocol.Blocks.readColumn_eval' depends on axioms: [propext, Classical.choice, Quot.sound]
'LeanerVM.Protocol.Blocks.layout' depends on axioms: [propext, Classical.choice, Quot.sound]
'LeanerVM.Protocol.Layout.comap' depends on axioms: [propext, Classical.choice, Quot.sound]
'LeanerVM.Protocol.Blocks.stack_eval_ambient' depends on axioms: [propext, Classical.choice, Quot.sound]
'LeanerVM.Protocol.Blocks.stack_eval_ambient_one' depends on axioms: [propext, Classical.choice, Quot.sound]
'LeanerVM.Protocol.Blocks.stackColumn_eval_ambient' depends on axioms: [propext, Classical.choice, Quot.sound]
'LeanerVM.Protocol.sumCube_padHigh' depends on axioms: [propext, Classical.choice, Quot.sound]
'LeanerVM.Protocol.evalMle_padHigh' depends on axioms: [propext, Classical.choice, Quot.sound]
'LeanerVM.Protocol.ColumnClaim.holds_iff_weighted' depends on axioms: [propext, Classical.choice, Quot.sound]
'LeanerVM.Protocol.eqWeight' depends on axioms: [propext, Classical.choice, Quot.sound]
'LeanerVM.Protocol.idxColumn_eval' depends on axioms: [propext, Classical.choice, Quot.sound]
'LeanerVM.Protocol.idxColumnEval_eq' depends on axioms: [propext, Classical.choice, Quot.sound]
'LeanerVM.Protocol.bytecodeColumn_answer_boolVec' depends on axioms: [propext, Classical.choice, Quot.sound]
'LeanerVM.Protocol.bytecodeColumn_eval' depends on axioms: [propext, Classical.choice, Quot.sound]
'LeanerVM.Protocol.BlockClaim.isValid_iff_pairing' depends on axioms: [propext, Classical.choice, Quot.sound]
'Probe.bytecodeColumn_answer_slots' depends on axioms: [propext, Classical.choice, Quot.sound]
exit=0
\end{srccode}

\subsubsection{\texttt{DuplicatesProbe.\allowbreak{}lean}, run with the old pin's Lean binary, not through Lake (I.5)}\label{sec:gen-layer1-I5}

\fhead{What it establishes}\texttt{evalMle\_\allowbreak{}lagrangeBasis} is CompPoly's \texttt{eqTilde\_\allowbreak{}eq\_\allowbreak{}prod}, three in-layer duplicates, and largest-first is not necessary for alignment (D.2, G.5). Run with the old pin's Lean binary on the untouched \texttt{b435631} artefacts, not with \texttt{lake env lean} (see the introduction of this part).

\textbf{Command.} \texttt{L\allowbreak{}E\allowbreak{}A\allowbreak{}N\allowbreak{}\_\allowbreak{}P\allowbreak{}A\allowbreak{}T\allowbreak{}H\allowbreak{}=\allowbreak{}…\allowbreak{} f\allowbreak{}l\allowbreak{}o\allowbreak{}c\allowbreak{}k\allowbreak{} …\allowbreak{} l\allowbreak{}e\allowbreak{}a\allowbreak{}n\allowbreak{}p\allowbreak{}r\allowbreak{}o\allowbreak{}v\allowbreak{}e\allowbreak{}r\allowbreak{}-{}\allowbreak{}-{}\allowbreak{}l\allowbreak{}e\allowbreak{}a\allowbreak{}n\allowbreak{}4\allowbreak{}-{}\allowbreak{}-{}\allowbreak{}-{}\allowbreak{}v\allowbreak{}4\allowbreak{}.\allowbreak{}3\allowbreak{}3\allowbreak{}.\allowbreak{}1\allowbreak{}/\allowbreak{}b\allowbreak{}i\allowbreak{}n\allowbreak{}/\allowbreak{}l\allowbreak{}e\allowbreak{}a\allowbreak{}n\allowbreak{} .\allowbreak{}c\allowbreak{}l\allowbreak{}a\allowbreak{}u\allowbreak{}d\allowbreak{}e\allowbreak{}/\allowbreak{}r\allowbreak{}e\allowbreak{}p\allowbreak{}o\allowbreak{}r\allowbreak{}t\allowbreak{}s\allowbreak{}/\allowbreak{}b\allowbreak{}l\allowbreak{}u\allowbreak{}e\allowbreak{}p\allowbreak{}r\allowbreak{}i\allowbreak{}n\allowbreak{}t\allowbreak{}-{}\allowbreak{}r\allowbreak{}e\allowbreak{}v\allowbreak{}i\allowbreak{}e\allowbreak{}w\allowbreak{}/\allowbreak{}p\allowbreak{}r\allowbreak{}o\allowbreak{}b\allowbreak{}e\allowbreak{}s\allowbreak{}/\allowbreak{}c\allowbreak{}o\allowbreak{}d\allowbreak{}e\allowbreak{}-{}\allowbreak{}l\allowbreak{}a\allowbreak{}y\allowbreak{}e\allowbreak{}r\allowbreak{}1\allowbreak{}/\allowbreak{}D\allowbreak{}u\allowbreak{}p\allowbreak{}l\allowbreak{}i\allowbreak{}c\allowbreak{}a\allowbreak{}t\allowbreak{}e\allowbreak{}s\allowbreak{}P\allowbreak{}r\allowbreak{}o\allowbreak{}b\allowbreak{}e\allowbreak{}.\allowbreak{}l\allowbreak{}e\allowbreak{}a\allowbreak{}n}

\begin{leancode}
/-
  Probe for task code-layer1, deliverables D and E: which Layer 1 statements are already in
  CompPoly at the pin (3468b38c), and which are instances of one another.

  Run from the repository root:
    flock .claude/reports/blueprint-review/logs/lean.lock \
      lake env lean .claude/reports/blueprint-review/probes/code-layer1/DuplicatesProbe.lean
-/
import LeanerVM.Protocol.ToCompPoly.BitProductTable
import LeanerVM.Protocol.ToCompPoly.AmbientStacking

open LeanerVM.Protocol CompPoly CMlPolynomialEval

namespace Probe

variable {R : Type*} [CommRing R]

/-- `evalMle_lagrangeBasis` is CompPoly's `eqTilde_eq_prod` (`Multilinear/Basic.lean:600`),
with the factors written in the other order. -/
example {n : ℕ} (w x : Vector R n) :
    evalMle (lagrangeBasis w) x = ∏ k : Fin n, ((1 - x[k]) * (1 - w[k]) + x[k] * w[k]) := by
  have h := eqTilde_eq_prod w x
  rw [eqTilde, ← eval_mle_eq_eval] at h
  rw [h]
  exact Finset.prod_congr rfl fun i _ ↦ by ring

/-- `unstack_eval` is `unstack_eval₂` at the identity map. -/
example (B : Blocks) {μ : ℕ} (hμ : B.total ≤ 2 ^ μ) (q : CMlPolynomialEval R μ) (b : Fin B.n)
    (z : Vector R (B.size b)) :
    evalMle q (B.extendPoint hμ b z) = evalMle (B.unstack hμ q b) z := by
  have h := B.unstack_eval₂ (RingHom.id R) hμ q b z
  simpa [eval₂Mle, CMlPolynomialEval.map] using h

/-- `stack_eval` is `unstack_eval` on a stack. -/
example (B : Blocks) (t : B.Tables R) {μ : ℕ} (hμ : B.total ≤ 2 ^ μ) (pad : R) (b : Fin B.n)
    (z : Vector R (B.size b)) :
    evalMle (B.stackAt t μ pad) (B.extendPoint hμ b z) = evalMle (t b) z := by
  rw [B.unstack_eval hμ, B.unstack_stackAt]

/-- `stack_eval_ambient_zero` is `stack_eval_ambient` at pad zero. -/
example (B : Blocks) (t : B.Tables R) {μ : ℕ} (hμ : B.total ≤ 2 ^ μ) (z : Vector R μ) :
    evalMle (B.stackAt t μ 0) z =
      ∑ b : Fin B.n, B.selectorWeight hμ b z * evalMle (t b) (B.lowPoint hμ b z) := by
  simpa using B.stack_eval_ambient t hμ 0 z

/-- The largest-first order is sufficient for alignment, not necessary: sizes 1, 1, 2 are
aligned at offsets 0, 2, 4 and are no `Blocks`. -/
example : ¬ Antitone (![1, 1, 2] : Fin 3 → ℕ) := fun h ↦
  absurd (h (show (0 : Fin 3) ≤ 2 by decide)) (by decide)

#guard (2 ^ 1 ∣ 0) ∧ (2 ^ 1 ∣ 2) ∧ (2 ^ 2 ∣ 4)

end Probe
\end{leancode}

\textbf{Recorded output.}

\begin{srccode}
(no output)
exit=0
\end{srccode}

\subsubsection{\texttt{StridedProbe.\allowbreak{}lean}, run with the old pin's Lean binary, not through Lake (I.6)}\label{sec:gen-layer1-I6}

\fhead{What it establishes}The low-index selection identity holds and is twenty lines from Layer 1's lemmas (B.8, G.1). Run with the old pin's Lean binary on the untouched \texttt{b435631} artefacts, not with \texttt{lake env lean} (see the introduction of this part).

\textbf{Command.} \texttt{L\allowbreak{}E\allowbreak{}A\allowbreak{}N\allowbreak{}\_\allowbreak{}P\allowbreak{}A\allowbreak{}T\allowbreak{}H\allowbreak{}=\allowbreak{}…\allowbreak{} f\allowbreak{}l\allowbreak{}o\allowbreak{}c\allowbreak{}k\allowbreak{} …\allowbreak{} l\allowbreak{}e\allowbreak{}a\allowbreak{}n\allowbreak{}p\allowbreak{}r\allowbreak{}o\allowbreak{}v\allowbreak{}e\allowbreak{}r\allowbreak{}-{}\allowbreak{}-{}\allowbreak{}l\allowbreak{}e\allowbreak{}a\allowbreak{}n\allowbreak{}4\allowbreak{}-{}\allowbreak{}-{}\allowbreak{}-{}\allowbreak{}v\allowbreak{}4\allowbreak{}.\allowbreak{}3\allowbreak{}3\allowbreak{}.\allowbreak{}1\allowbreak{}/\allowbreak{}b\allowbreak{}i\allowbreak{}n\allowbreak{}/\allowbreak{}l\allowbreak{}e\allowbreak{}a\allowbreak{}n\allowbreak{} .\allowbreak{}c\allowbreak{}l\allowbreak{}a\allowbreak{}u\allowbreak{}d\allowbreak{}e\allowbreak{}/\allowbreak{}r\allowbreak{}e\allowbreak{}p\allowbreak{}o\allowbreak{}r\allowbreak{}t\allowbreak{}s\allowbreak{}/\allowbreak{}b\allowbreak{}l\allowbreak{}u\allowbreak{}e\allowbreak{}p\allowbreak{}r\allowbreak{}i\allowbreak{}n\allowbreak{}t\allowbreak{}-{}\allowbreak{}r\allowbreak{}e\allowbreak{}v\allowbreak{}i\allowbreak{}e\allowbreak{}w\allowbreak{}/\allowbreak{}p\allowbreak{}r\allowbreak{}o\allowbreak{}b\allowbreak{}e\allowbreak{}s\allowbreak{}/\allowbreak{}c\allowbreak{}o\allowbreak{}d\allowbreak{}e\allowbreak{}-{}\allowbreak{}l\allowbreak{}a\allowbreak{}y\allowbreak{}e\allowbreak{}r\allowbreak{}1\allowbreak{}/\allowbreak{}S\allowbreak{}t\allowbreak{}r\allowbreak{}i\allowbreak{}d\allowbreak{}e\allowbreak{}d\allowbreak{}P\allowbreak{}r\allowbreak{}o\allowbreak{}b\allowbreak{}e\allowbreak{}.\allowbreak{}l\allowbreak{}e\allowbreak{}a\allowbreak{}n}

\begin{leancode}
/-
  Probe for task code-layer1, deliverable C and the findings: the selection identity Layer 1
  does not have. leanVM reads the eighteen BLAKE2S limb columns off `q_flock` by freezing the
  LOW coordinates to a slot's bits (`crates/pcs/src/stack_open.rs:84-97`, `StackClaim::Strided`:
  "Equivalent to a `Point` with `low_point = slot_bits ++ point`"). Layer 1 slices on the high
  index only. This file states and proves the mirrored identity, to show what is missing and
  that it is a generic fact of the same size as `evalMle_append_boolVec`.

  Run from the repository root:
    flock .claude/reports/blueprint-review/logs/lean.lock \
      lake env lean .claude/reports/blueprint-review/probes/code-layer1/StridedProbe.lean
-/
import LeanerVM.Protocol.ToCompPoly.Multilinear
import LeanerVM.Parameters.Field

open LeanerVM.Protocol LeanerVM.Parameters CompPoly CMlPolynomialEval

namespace Probe

variable {R : Type*} [CommRing R]

/-- The strided slice of a table at the low index `i`: the entries whose low `k` bits are `i`. -/
def sliceLow {k m : ℕ} (t : CMlPolynomialEval R (k + m)) (i : Fin (2 ^ k)) :
    CMlPolynomialEval R m :=
  Vector.ofFn fun j ↦ t[cubeIndex i j]

/-- A Boolean low coordinate selects the strided slice at that index. -/
theorem evalMle_boolVec_append {k m : ℕ} (t : CMlPolynomialEval R (k + m)) (i : Fin (2 ^ k))
    (s : Vector R m) :
    evalMle t ((boolVec i : Vector R k) ++ s) = evalMle (sliceLow t i) s := by
  rw [evalMle_eq_sum, sum_cube_split, evalMle_eq_sum]
  refine Finset.sum_congr rfl fun j _ ↦ ?_
  rw [Finset.sum_eq_single i]
  · rw [lagrangeBasis_cubeIndex, lowVec_append, highVec_append]
    have h := lagrangeBasis_boolVec (R := R) i i
    simp only [if_true] at h
    simp only [Fin.getElem_fin] at h ⊢
    rw [h, one_mul]
    simp [sliceLow]
  · intro i' _ hi
    rw [lagrangeBasis_cubeIndex, lowVec_append, highVec_append]
    have h := lagrangeBasis_boolVec (R := R) i i'
    simp only [if_neg hi] at h
    simp only [Fin.getElem_fin] at h ⊢
    rw [h, zero_mul, mul_zero]
  · intro h
    exact absurd (Finset.mem_univ _) h

/-- Slot 1 of four (two low bits) of an eight-cell table: cells 1 and 5. -/
def packed : CMlPolynomialEval K 3 := #v[10, 11, 12, 13, 20, 21, 22, 23]

#guard (sliceLow (k := 2) (m := 1) packed (1 : Fin 4)).toList = [11, 21]
-- The point `(slot bits 1, 0 | z)` reads the strided slice at `z`.
#guard evalMle packed (#v[1, 0, 7] : Vector K 3) =
  evalMle (#v[11, 21] : CMlPolynomialEval K 1) #v[7]
-- Near miss: the aligned slice at high index 1 is cells 2 and 3, another column.
#guard (slice (k := 1) (m := 2) packed (1 : Fin 4)).toList = [12, 13]

#print axioms evalMle_boolVec_append

end Probe
\end{leancode}

\textbf{Recorded output.}

\begin{srccode}
'Probe.evalMle_boolVec_append' depends on axioms: [propext, Classical.choice, Quot.sound]
exit=0
\end{srccode}

\subsubsection{The counting script of section E (\texttt{count2.\allowbreak{}py}, scratch, not a probe) (I.7)}\label{sec:gen-layer1-I7}

\fhead{What it establishes}The counting script of the audit surface of Layer 1 (section E); scratch, not a probe.

\textbf{Command.} \texttt{python3 -{}B count2.\allowbreak{}py}

\begin{srccode}
import re, subprocess, json
files = ["LeanerVM/Protocol/ToCompPoly/Multilinear.lean","LeanerVM/Protocol/ToCompPoly/BitProductTable.lean",
 "LeanerVM/Protocol/ToCompPoly/Stacking.lean","LeanerVM/Protocol/ToCompPoly/AmbientStacking.lean",
 "LeanerVM/Protocol/Stack.lean","LeanerVM/Protocol/Padding.lean","LeanerVM/Protocol/ClaimWeights.lean",
 "LeanerVM/Protocol/BlockClaims.lean","LeanerVM/Protocol/FixedColumns.lean"]
decl = re.compile(r'^(?:@\[[^\]]*\]\s*)?(private\s+)?(def|theorem|abbrev|structure|lemma)\s+(\S+)')
tot = {}
for f in files:
    src = subprocess.run(["git","show","b435631:"+f],capture_output=True,text=True,cwd="/home/scaraven/Documents/Verified-zkEVM/leanerVM").stdout.split("\n")
    rows=[]
    for i,l in enumerate(src):
        m = decl.match(l)
        if not m or m.group(1): continue
        kind=m.group(2); name=m.group(3)
        j=i; n=0
        if kind in ("def","abbrev","structure"):
            while j < len(src) and src[j].strip()!="":
                n+=1; j+=1
        else:
            while True:
                n+=1
                if ':=' in src[j]: break
                j+=1
        k=i-1; doc=0
        while k>=0 and src[k].startswith('omit'): k-=1
        if k>=0 and src[k].rstrip().endswith('-/'):
            while k>=0:
                doc+=1
                if src[k].lstrip().startswith('/--'): break
                k-=1
        rows.append((i+1,kind,name,n,doc))
    tot[f]=rows
    print(f, len(rows), sum(r[3] for r in rows), sum(r[3]+r[4] for r in rows))
json.dump(tot, open("decls2.json","w"))
\end{srccode}

\textbf{Recorded output.}

\begin{srccode}
LeanerVM/Protocol/ToCompPoly/Multilinear.lean 43 98 137
LeanerVM/Protocol/ToCompPoly/BitProductTable.lean 13 30 39
LeanerVM/Protocol/ToCompPoly/Stacking.lean 38 108 147
LeanerVM/Protocol/ToCompPoly/AmbientStacking.lean 10 35 48
LeanerVM/Protocol/Stack.lean 9 34 46
LeanerVM/Protocol/Padding.lean 6 12 20
LeanerVM/Protocol/ClaimWeights.lean 4 12 18
LeanerVM/Protocol/BlockClaims.lean 8 30 40
LeanerVM/Protocol/FixedColumns.lean 12 29 44
\end{srccode}

(columns: public declarations, statement lines, statement lines with docstrings; the split by class uses the lists of section A.)

\subsubsection{Not done, unverified, and contradictions with the brief (J)}\label{sec:gen-layer1-J}

\begin{itemize}
\item \textbf{Unverified: the Rust encoder was not executed.} The bytecode table \texttt{ValuesProbe} compares against is built by \texttt{values.\allowbreak{}py}'s transcription of \texttt{bytecode\_\allowbreak{}columns} and \texttt{stacked\_\allowbreak{}bytecode\_\allowbreak{}table}, checked by reading against \texttt{cpu/\allowbreak{}layout.\allowbreak{}rs:229-{}309} and \texttt{leaf.\allowbreak{}rs:585-{}604} (B.6). Running the pinned Rust \texttt{bytecode\_\allowbreak{}table} on the same program and comparing would verify it.
\item \textbf{Unverified: the proposed strided reader and layout combinator} (G.1) are sketched, not built; only the selection lemma is proved (\texttt{StridedProbe}).
\item \textbf{Not run at the new pin.} No probe was run against \texttt{144c5aa}; \texttt{ValuesProbe} and \texttt{StridedProbe} need their \texttt{K} numerals rewritten first (H). The conclusions are about \texttt{b435631} and the old pins, which is the review's object (brief §8).
\item \textbf{How two Lean probes ran} (I): \texttt{DuplicatesProbe} and \texttt{StridedProbe} were run with the old pin's Lean binary and \texttt{LEAN\_\allowbreak{}PATH} on the untouched \texttt{b435631} artefacts, not with \texttt{lake env lean}, because the checkout had moved; this was before the instruction to stop running Lean reached me. If the orchestrator discounts them, D.2's duplicate claims and B.8's lemma stand as short proofs on paper (a three-line rewrite with \texttt{eqTilde\_\allowbreak{}eq\_\allowbreak{}prod}; the mirror of \texttt{evalMle\_\allowbreak{}append\_\allowbreak{}boolVec}).
\item \textbf{Contradiction with the brief's §2} (not of my making): the leanerVM checkout was moved during this task (reflog in I), and \texttt{.\allowbreak{}lake/\allowbreak{}packages/} switched to the new pins; the coordinator's resume message (\enquote{nothing else changed}) predates it. I moved no checkout, ran no build, and wrote only under \texttt{.\allowbreak{}c\allowbreak{}l\allowbreak{}a\allowbreak{}u\allowbreak{}d\allowbreak{}e\allowbreak{}/\allowbreak{}r\allowbreak{}e\allowbreak{}p\allowbreak{}o\allowbreak{}r\allowbreak{}t\allowbreak{}s\allowbreak{}/\allowbreak{}b\allowbreak{}l\allowbreak{}u\allowbreak{}e\allowbreak{}p\allowbreak{}r\allowbreak{}i\allowbreak{}n\allowbreak{}t\allowbreak{}-{}\allowbreak{}r\allowbreak{}e\allowbreak{}v\allowbreak{}i\allowbreak{}e\allowbreak{}w\allowbreak{}/} (the dossier and the six \texttt{*.\allowbreak{}out} files).
\item \textbf{Facts of sibling dossiers used, and checked where I read the same lines}: the tie order and the six shared columns first (\texttt{gt-{}bus.\allowbreak{}md} A.5, G7; I read \texttt{witness.\allowbreak{}rs:63-{}102} and \texttt{cpu/\allowbreak{}layout.\allowbreak{}rs:9-{}49}); the leaf stacks' pad \texttt{1} and no floor (\texttt{gt-{}bus.\allowbreak{}md} A.5; not re-read); the limb columns as strided slots of \texttt{q\_\allowbreak{}flock} (\texttt{gt-{}flock-{}ring.\allowbreak{}md} 2.4, 5.4; I read \texttt{stack\_\allowbreak{}open.\allowbreak{}rs:73-{}98}, \texttt{hash\_\allowbreak{}flock.\allowbreak{}rs:82-{}115}, \texttt{py:878-{}895}); the table sumcheck binds the highest variable first and pads by the product of the high variables (\texttt{gt-{}table-{}pub.\allowbreak{}md}; I read \texttt{constraints.\allowbreak{}rs:250-{}290}, \texttt{py:598-{}626}). All agree with what I read. Nothing in the brief's §7 is contradicted by this task's sources.
\item \textbf{Not examined}: the rest of the blueprint's Layer 3 (beyond the lines that touch Layer 1), the leaf stacks' block order in code (not built), the pad cells of the committed stack past the last lane (\texttt{witness.\allowbreak{}rs:41-{}60}; no Layer 1 statement concerns them).
\end{itemize}
